% Mesures, incertitudes et qualités d'un instrument de mesure — Physique 1ère D — Cours signé OnBuch+
% Programme officiel MINESEC. Compiler avec : tectonic N01-mesures-incertitudes-qualites-instrument.tex
\newcommand{\DOCMATIERE}{Physique}
\newcommand{\DOCNIVEAU}{1ère D}
\newcommand{\DOCMODULE}{Module 1 : Mesures et incertitudes}
\newcommand{\DOCLECON}{N01}
\newcommand{\DOCTITRE}{Mesures, incertitudes et qualités d'un instrument de mesure}
% OnBuch+ — préambule « cours signés » ENRICHI (compilé avec tectonic / XeLaTeX)
% Système visuel « Pop » d'OnBuch+ : fond crème (jamais blanc), bordures encre
% épaisses, ombres « dures » décalées, titres Archivo Black en capitales.
% Version MATHÉMATIQUES : graphes (pgfplots/tikz), boîtes pédagogiques
% (définition, propriété, méthode, attention, le savais-tu, exercice résolu,
% exercice, TP, bilan), page de garde et signature OnBuch+ ancrée en bas.
%
% Macros attendues dans main.tex AVANT \input{preamble} :
%   \DOCMATIERE  \DOCNIVEAU  \DOCMODULE  \DOCLECON (numéro)  \DOCTITRE
\documentclass[11pt]{article}

\usepackage{fontspec}
\usepackage[french]{babel}
\usepackage[a4paper,margin=2cm,top=3.4cm,bottom=2.8cm,headheight=1.4cm,footskip=1.3cm]{geometry}
\usepackage{amsmath,amssymb,amsfonts}
\usepackage{mathtools} % \xmapsto, \coloneqq... souvent utilisés par les modèles
\usepackage{cancel} % simplifications barrées (bilans de forces)
% \abs{z} (module, valeur absolue) et \norm{u} (norme), utilisés par les modèles
\providecommand{\abs}{}\let\abs\relax\DeclarePairedDelimiter{\abs}{\lvert}{\rvert}
\providecommand{\norm}{}\let\norm\relax\DeclarePairedDelimiter{\norm}{\lVert}{\rVert}
\usepackage[table]{xcolor} % [table] : fournit aussi \rowcolors (alternance de lignes)
\usepackage{pagecolor}
\usepackage{tikz}
\usetikzlibrary{arrows.meta,positioning,calc,shapes.geometric,shapes.misc,shapes.arrows,decorations.pathmorphing,decorations.pathreplacing,decorations.markings,patterns,shadows,mindmap,trees,fit,backgrounds,matrix,intersections,angles,quotes}
\usepackage{pgfplots}
\usepackage[version=4]{mhchem} % équations nucléaires \ce{^{235}_{92}U}
\usepackage[european,siunitx]{circuitikz} % schémas électriques
\pgfplotsset{compat=1.18}
\usepgfplotslibrary{fillbetween,groupplots} % aires entre courbes (calcul intégral)
\usepackage{enumitem}
\usepackage{fancyhdr}
% [most] charge toutes les bibliothèques tcolorbox (listings, minted, raster,
% documentation...) et gonfle inutilement la mémoire TeX sur un document long
% (~300 boîtes) : on ne charge que ce qui est réellement utilisé.
\usepackage[skins,breakable]{tcolorbox}
\usepackage{graphicx}
\usepackage{colortbl}
\usepackage{booktabs}
\usepackage{tabularx}
\usepackage{diagbox} % case d'en-tête diagonale des tableaux à double entrée
% Colonnes extensibles centrée / gauche / droite (C, L, R), souvent utilisées par les modèles
\newcolumntype{C}{>{\centering\arraybackslash}X}
\newcolumntype{L}{>{\raggedright\arraybackslash}X}
\newcolumntype{R}{>{\raggedleft\arraybackslash}X}
\usepackage{array}
\usepackage{multicol}
\usepackage{siunitx}
% parse-numbers=false : accepte aussi \SI{2,5 \times 10^{-3}}{...}
\sisetup{locale=FR,output-decimal-marker={,},group-minimum-digits=4,parse-numbers=false}
\DeclareSIUnit\ohm{\ensuremath{\symup{\Omega}}} % U+2126 absent de la police
\DeclareSIUnit\division{div} % réglages d'oscilloscope (ms/div, V/div)
\DeclareSIUnit\tr{tr} % tours (vitesse de rotation, ex. \SI{1500}{\tr\per\minute})
\DeclareSIUnit\ch{ch} % cheval-vapeur (puissance moteur, unité usuelle hors SI)
\DeclareSIUnit\franccfa{FCFA} % franc CFA (coûts, contexte camerounais)
\DeclareSIUnit\dioptre{\ensuremath{\delta}} % vergence d'une lentille (dioptrie)
\usepackage{qrcode}
% \degree : commande fréquemment utilisée par les modèles pour un angle ou une
% température en dehors de \SI{}{} (non fournie par les paquets chargés).
\providecommand{\degree}{^{\circ}}
% \qed (fin de démonstration), utilisé par les modèles sans amsthm
\providecommand{\qed}{\hfill$\square$}
% \barre : double barre des valeurs interdites dans un tableau de variations (tkz-tab)
\providecommand{\barre}{\ensuremath{\|}}
% environnement proof (amsthm non chargé) utilisé par les modèles
\ifdefined\proof\else\newenvironment{proof}[1][Démonstration]{\par\noindent\textit{#1.}\ }{\qed\par}\fi
\usepackage{lastpage}
\usepackage{hyperref}
\hypersetup{hidelinks}

% ============================================================
% Palette « Pop » OnBuch+ — mêmes hex que la classe P (plus_ui.dart)
% ============================================================
\definecolor{poporange}{HTML}{F59321}
\definecolor{popdark}{HTML}{E07A0C}
\definecolor{popcream}{HTML}{FFF6E8}
\definecolor{popink}{HTML}{1C1714}
\definecolor{poppurple}{HTML}{7A5AE0}
\definecolor{popgreen}{HTML}{2E9E6A}
\definecolor{poppink}{HTML}{E0457B}
\definecolor{popblue}{HTML}{2F7DE1}
\definecolor{popgold}{HTML}{C98A12}
\definecolor{popmuted}{HTML}{8A7F76}
\definecolor{popline}{HTML}{EADFD2}
\definecolor{popgray}{HTML}{8A7F76}
\colorlet{popgrey}{popgray}
% Alias tolérants (le modèle nomme parfois les couleurs différemment)
\colorlet{popwhite}{white}
\colorlet{popblack}{popink}
\colorlet{popred}{poppink}
\colorlet{popyellow}{popgold}
% Teintes claires pour fonds de tableaux / zones
\colorlet{popblueL}{popblue!10!white}
\colorlet{poporangeL}{poporange!14!white}
\colorlet{popgreenL}{popgreen!12!white}
\colorlet{poppurpleL}{poppurple!12!white}
\colorlet{poppinkL}{poppink!12!white}
\colorlet{popgoldL}{popgold!14!white}

\pagecolor{popcream}

% ============================================================
% Typographie
% ============================================================
\defaultfontfeatures{Ligatures=TeX}
\setmainfont{PlusJakartaSans-Medium.ttf}[Path=../../../fonts/,BoldFont=PlusJakartaSans-Bold.ttf]
% Maths en sans-serif (Fira Math) avec chiffres et lettres droites de Plus
% Jakarta, pour que les formules se fondent dans le texte.
\usepackage[math-style=ISO,bold-style=ISO]{unicode-math}
\setmathfont{FiraMath-Regular.otf}[Path=../../../fonts/]
\setmathfont{PlusJakartaSans-Medium.ttf}[Path=../../../fonts/,range={up/num,up/{Latin,latin},\mathup}]
\usepackage{icomma} % virgule décimale sans espace en mode math
% Caractères mathématiques Unicode absents des polices de texte (Plus Jakarta,
% Archivo Black) : rendus par leur équivalent mathématique, en texte comme en
% formule — sinon ils s'affichent en carré vide (ex. « Arithmétique dans ℤ »).
\usepackage{newunicodechar}
\newunicodechar{ℝ}{\ensuremath{\mathbb{R}}}
\newunicodechar{ℤ}{\ensuremath{\mathbb{Z}}}
\newunicodechar{ℂ}{\ensuremath{\mathbb{C}}}
\newunicodechar{ℕ}{\ensuremath{\mathbb{N}}}
\newunicodechar{ℚ}{\ensuremath{\mathbb{Q}}}
\newunicodechar{𝔻}{\ensuremath{\mathbb{D}}}
\newunicodechar{α}{\ensuremath{\alpha}}
\newunicodechar{β}{\ensuremath{\beta}}
\newunicodechar{γ}{\ensuremath{\gamma}}
\newunicodechar{δ}{\ensuremath{\delta}}
\newunicodechar{ε}{\ensuremath{\varepsilon}}
\newunicodechar{θ}{\ensuremath{\theta}}
\newunicodechar{λ}{\ensuremath{\lambda}}
\newunicodechar{μ}{\ensuremath{\mu}}
\newunicodechar{σ}{\ensuremath{\sigma}}
\newunicodechar{τ}{\ensuremath{\tau}}
\newunicodechar{φ}{\ensuremath{\varphi}}
\newunicodechar{ψ}{\ensuremath{\psi}}
\newunicodechar{ω}{\ensuremath{\omega}}
\newunicodechar{Σ}{\ensuremath{\Sigma}}
% majuscules produites par \MakeUppercase dans les titres (α-aminés -> Α)
\newunicodechar{Α}{\ensuremath{\alpha}}
\newunicodechar{Β}{\ensuremath{\beta}}
\newunicodechar{Γ}{\ensuremath{\Gamma}}
\newunicodechar{Δ}{\ensuremath{\Delta}}
\newunicodechar{Ε}{\ensuremath{\varepsilon}}
\newunicodechar{Θ}{\ensuremath{\Theta}}
\newunicodechar{Λ}{\ensuremath{\Lambda}}
\newunicodechar{Μ}{\ensuremath{\mu}}
\newunicodechar{Τ}{\ensuremath{\tau}}
\newunicodechar{Φ}{\ensuremath{\Phi}}
\newunicodechar{Ψ}{\ensuremath{\Psi}}
\newunicodechar{Ω}{\ensuremath{\Omega}}
\newunicodechar{↦}{\ensuremath{\mapsto}}
\newunicodechar{∈}{\ensuremath{\in}}
\newunicodechar{∉}{\ensuremath{\notin}}
\newunicodechar{∧}{\ensuremath{\wedge}}
\newunicodechar{⇔}{\ensuremath{\Leftrightarrow}}
\newunicodechar{⟺}{\ensuremath{\Longleftrightarrow}}
\newunicodechar{⇒}{\ensuremath{\Rightarrow}}
\newunicodechar{⟹}{\ensuremath{\Longrightarrow}}
\newunicodechar{∘}{\ensuremath{\circ}}
\newunicodechar{∪}{\ensuremath{\cup}}
\newunicodechar{∩}{\ensuremath{\cap}}
\newunicodechar{≡}{\ensuremath{\equiv}}
\newunicodechar{⊂}{\ensuremath{\subset}}
\newunicodechar{⊥}{\ensuremath{\perp}}
\newunicodechar{∃}{\ensuremath{\exists}}
\newunicodechar{∀}{\ensuremath{\forall}}
\newunicodechar{∛}{\ensuremath{\sqrt[3]{\,}}}
\newunicodechar{∜}{\ensuremath{\sqrt[4]{\,}}}
\newunicodechar{⃗}{\ensuremath{{}^{\to}}}
\newunicodechar{ⁿ}{\ifmmode^{n}\else\textsuperscript{\ensuremath{n}}\fi}
\newunicodechar{ˣ}{\ifmmode^{x}\else\textsuperscript{\ensuremath{x}}\fi}
\newunicodechar{ᵇ}{\ifmmode^{b}\else\textsuperscript{\ensuremath{b}}\fi}
\newunicodechar{ʳ}{\ifmmode^{r}\else\textsuperscript{\ensuremath{r}}\fi}
\newunicodechar{ᵘ}{\ifmmode^{u}\else\textsuperscript{\ensuremath{u}}\fi}
\newunicodechar{ʸ}{\ifmmode^{y}\else\textsuperscript{\ensuremath{y}}\fi}
\newunicodechar{ᵃ}{\ifmmode^{a}\else\textsuperscript{\ensuremath{a}}\fi}
\newunicodechar{ᵉ}{\ifmmode^{e}\else\textsuperscript{\ensuremath{e}}\fi}
\newunicodechar{ᵏ}{\ifmmode^{k}\else\textsuperscript{\ensuremath{k}}\fi}
\newunicodechar{ᵖ}{\ifmmode^{p}\else\textsuperscript{\ensuremath{p}}\fi}
\newunicodechar{ᴺ}{\ifmmode^{N}\else\textsuperscript{\ensuremath{N}}\fi}
\newunicodechar{ᶜ}{\ifmmode^{c}\else\textsuperscript{\ensuremath{c}}\fi}
\newunicodechar{ʰ}{\ifmmode^{h}\else\textsuperscript{\ensuremath{h}}\fi}
\newunicodechar{ᵠ}{\ifmmode^{q}\else\textsuperscript{\ensuremath{q}}\fi}
\newunicodechar{ₙ}{\ifmmode_{n}\else\textsubscript{\ensuremath{n}}\fi}
\newunicodechar{ₐ}{\ifmmode_{a}\else\textsubscript{\ensuremath{a}}\fi}
\newunicodechar{ᵢ}{\ifmmode_{i}\else\textsubscript{\ensuremath{i}}\fi}
\newunicodechar{ᵣ}{\ifmmode_{r}\else\textsubscript{\ensuremath{r}}\fi}
\newunicodechar{ₖ}{\ifmmode_{k}\else\textsubscript{\ensuremath{k}}\fi}
\newunicodechar{ᵦ}{\ifmmode_{\beta}\else\textsubscript{\ensuremath{\beta}}\fi}
\newunicodechar{ₓ}{\ifmmode_{x}\else\textsubscript{\ensuremath{x}}\fi}
\newfontfamily\popdisplay{ArchivoBlack-Regular.ttf}[Path=../../../fonts/]
\newfontfamily\poplogo{SpaceGrotesk-Bold.ttf}[Path=../../../fonts/]
\newfontfamily\popsemi{PlusJakartaSans-SemiBold.ttf}[Path=../../../fonts/]
\newfontfamily\popxbold{PlusJakartaSans-ExtraBold.ttf}[Path=../../../fonts/]
\newfontfamily\popxboldls{PlusJakartaSans-ExtraBold.ttf}[Path=../../../fonts/,LetterSpace=3.0]

\setlist[enumerate]{leftmargin=1.7em,itemsep=5pt,topsep=4pt}
\setlist[itemize]{leftmargin=1.7em,itemsep=4pt,topsep=4pt,label={\color{poporange}\textbullet}}
\setlength{\parindent}{0pt}
\setlength{\parskip}{5pt}
\renewcommand{\arraystretch}{1.25}

% pgfplots : style Pop par défaut
\pgfplotsset{
  every axis/.append style={axis line style={popink,line width=1pt},tick style={popink},
    grid style={popline,line width=0.5pt},label style={font=\small},tick label style={font=\footnotesize}},
  popcurve/.style={poporange,line width=1.8pt,no markers,smooth},
}
% Même style disponible dans un \draw TikZ simple (hors pgfplots)
\tikzset{popcurve/.style={poporange,line width=1.8pt}}
\tikzset{popgrid/.style={popline,very thin}}
\colorlet{popcurve}{poporange} % le nom du style est parfois utilisé comme couleur

% ============================================================
% Pill / squiggle
% ============================================================
\newcommand{\poppill}[3]{%
  \tikz[baseline=-0.55ex]\node[draw=popink,line width=1.2pt,rounded corners=2.6pt,%
    fill=#2,inner xsep=6.5pt,inner ysep=3pt]{%
    \popxboldls\fontsize{7}{7}\selectfont\color{#3}\MakeUppercase{#1}};%
}
\newcommand{\popsquiggle}[1]{%
  \tikz[baseline=-0.4ex]{
    \draw[#1,line width=2.6pt,line cap=round]
      plot [smooth] coordinates {(0,0.09) (0.34,0.01) (0.68,0.21) (1.02,0.01) (1.36,0.10)};
  }%
}
% Mot-clé surligné (marqueur orange sous le texte)
\newcommand{\cle}[1]{{\popxbold\color{popdark}#1}}

% ============================================================
% En-tête / pied
% ============================================================
\pagestyle{fancy}
\fancyhf{}
\renewcommand{\headrulewidth}{0pt}
\renewcommand{\footrulewidth}{0pt}
\lhead{\raisebox{-0.15ex}{\poplogo\fontsize{13}{13}\selectfont\color{popink}OnBuch\color{poporange}+}}
\rhead{\poppill{\DOCMATIERE\ \textbullet\ \DOCNIVEAU\ \textbullet\ Leçon \DOCLECON}{white}{popink}}
\lfoot{\popxbold\fontsize{7.6}{7.6}\selectfont\color{popmuted}\MakeUppercase{Cours signé OnBuch+}\ \textbullet\ {\popsemi\DOCTITRE}}
\rfoot{\tikz[baseline=-0.6ex]\node[rounded corners=8.5pt,fill=popink,minimum height=17pt,minimum width=17pt,inner xsep=5pt,inner sep=0]{\popxbold\fontsize{8}{8}\selectfont\color{popcream}\thepage\,/\,\pageref*{LastPage}};}

% ============================================================
% Titres
% ============================================================
\newcommand{\chaptitre}[2]{%
  \begin{tcolorbox}[enhanced,colback=poporange,colframe=popink,boxrule=2.6pt,arc=11pt,
      drop shadow=popink,left=15pt,right=15pt,top=12pt,bottom=11pt,before skip=3pt,after skip=18pt]
    \poppill{#2}{popink}{popcream}\par\vspace{9pt}
    {\raggedright\hyphenpenalty=10000\exhyphenpenalty=10000\popdisplay\fontsize{22}{24}\selectfont\color{popcream}\MakeUppercase{#1}\par}
  \end{tcolorbox}
}
% Section numérotée : pastille orange avec le numéro + titre Archivo
\newcounter{popsec}
\newcommand{\coursec}[1]{%
  \stepcounter{popsec}\setcounter{popsub}{0}%
  \par\vspace{16pt}\noindent
  \begin{minipage}{\linewidth}
  \tikz[baseline=-0.7ex]\node[draw=popink,line width=1.6pt,fill=poporange,rounded corners=4pt,minimum size=22pt,inner sep=0]{\popdisplay\fontsize{12}{12}\selectfont\color{popcream}\thepopsec};%
  \hspace{8pt}{\popdisplay\fontsize{15}{17}\selectfont\color{popink}\MakeUppercase{#1}}\par
  \vspace{4pt}\noindent{\color{poporange}\rule{36pt}{3.6pt}}
  \end{minipage}\par\nopagebreak\vspace{8pt}
}
\newcounter{popsub}
\newcommand{\courssub}[1]{%
  \stepcounter{popsub}%
  \par\vspace{9pt}\noindent{\popxbold\fontsize{11.5}{13}\selectfont\color{popink}{\color{poporange}\thepopsec.\thepopsub}\hspace{6pt}#1}\par\nopagebreak\vspace{4pt}
}

% ============================================================
% Boîtes pédagogiques — même recette : fond blanc, bordure épaisse colorée,
% ombre dure encre, badge plein. Titre optionnel [#1].
% ============================================================
\tcbset{popbase/.style={breakable,enhanced,colback=white,boxrule=2.3pt,arc=9pt,
  shadow={3pt}{-3pt}{0pt}{popink},left=14pt,right=14pt,top=11pt,bottom=11pt,before skip=11pt,after skip=15pt,
  fontupper=\popsemi\fontsize{10.6}{14.5}\selectfont\color{popink},
  fontlower=\popsemi\fontsize{10.6}{14.5}\selectfont\color{popink}}}
% Coche dessinée (le glyphe ✓ n'existe pas dans les polices du document)
\newcommand{\popcheck}[1]{\tikz[baseline=-0.5ex]\draw[#1,line width=1.6pt,line cap=round,line join=round](-0.13,0.0)--(-0.04,-0.09)--(0.13,0.1);}
\providecommand{\coche}{\popcheck{popgreen}}
\providecommand{\resultat}[1]{\par\smallskip\noindent\fcolorbox{popgreen}{popgreenL}{\parbox{\dimexpr\linewidth-2\fboxsep-2\fboxrule\relax}{\textbf{Résultat.} #1}}\par\smallskip}
\newcommand{\popboxhead}[3]{\poppill{#1}{#2}{white}\ifx\relax#3\relax\else\hspace{6pt}{\popxbold\fontsize{10.6}{12}\selectfont\color{popink}#3}\fi\par\vspace{7pt}}

\newtcolorbox{aretenir}[1][]{popbase,colframe=popgreen,before upper={\popboxhead{À retenir}{popgreen}{#1}}}
\newtcolorbox{exemplebox}[1][]{popbase,colframe=poporange,before upper={\popboxhead{Exemple}{poporange}{#1}}}
\newtcolorbox{definition}[1][]{popbase,colframe=popblue,before upper={\popboxhead{Définition}{popblue}{#1}}}
\newtcolorbox{propriete}[1][]{popbase,colframe=poppurple,before upper={\popboxhead{Propriété}{poppurple}{#1}}}
\newtcolorbox{methode}[1][]{popbase,colframe=popgold,before upper={\popboxhead{Méthode}{popgold}{#1}}}
\newtcolorbox{attention}[1][]{popbase,colframe=poppink,before upper={\popboxhead{Attention}{poppink}{#1}}}
\newtcolorbox{experience}[1][]{popbase,colframe=popblue,colback=popblueL,before upper={\popboxhead{Expérience}{popblue}{#1}}}
% « Le savais-tu ? » : carte violette pleine (contexte, culture, Cameroun)
\newtcolorbox{savaistu}[1][]{popbase,colback=poppurple,colframe=popink,
  fontupper=\popsemi\fontsize{10.4}{14.2}\selectfont\color{white},
  before upper={\poppill{Le savais-tu\,?}{white}{poppurple}\ifx\relax#1\relax\else\hspace{6pt}{\popxbold\color{white}#1}\fi\par\vspace{7pt}}}
% Exercice résolu : énoncé en haut, \tcblower puis la solution sur fond crème
\newcounter{popexo}\newcounter{popexr}
\newtcolorbox{exoresolu}[1][]{popbase,colframe=poporange,colbacklower=popcream,
  segmentation style={popink,line width=1.2pt,dashed},
  before upper={\stepcounter{popexr}\popboxhead{Exercice résolu \thepopexr}{poporange}{#1}},
  before lower={\poppill{Solution}{popink}{popcream}\par\vspace{6pt}}}
% Exercice (non résolu en place) — la correction est dans la partie Corrigés
\newtcolorbox{exercice}[1][]{popbase,colframe=popink,
  before upper={\stepcounter{popexo}\popboxhead{Exercice \thepopexo}{popink}{#1}}}
% Corrigé
\newtcolorbox{corrige}[1][]{popbase,colframe=popgreen,colback=popgreenL,
  before upper={\popboxhead{Corrigé}{popgreen}{#1}}}
% Objectifs / prérequis / bilan
\newtcolorbox{objectifs}{popbase,colframe=popink,colback=poporangeL,
  before upper={\popboxhead{Objectifs de la leçon}{popink}{}}}
\newtcolorbox{prerequis}{popbase,colframe=popink,colback=popblueL,
  before upper={\popboxhead{Prérequis}{popblue}{}}}
\newtcolorbox{bilan}{popbase,colframe=popink,colback=white,boxrule=2.8pt,
  before upper={\popboxhead{Fiche bilan}{poporange}{L'essentiel à connaître pour l'examen}}}
% Figure encadrée avec légende
\newtcolorbox{popfigure}[1][]{enhanced,breakable=false,colback=white,colframe=popline,boxrule=1.4pt,arc=8pt,
  left=10pt,right=10pt,top=10pt,bottom=8pt,before skip=10pt,after skip=12pt,halign=center,
  lower separated=false,halign lower=center,
  fontlower=\popsemi\fontsize{9}{11.5}\selectfont\color{popmuted},
  #1}
\newcommand{\legende}[1]{\par\vspace{6pt}{\popsemi\fontsize{9}{11.5}\selectfont\color{popmuted}#1}}

% ============================================================
% Signature OnBuch+
% ============================================================
\newcommand{\poplogoinline}[1]{{\poplogo\fontsize{#1}{#1}\selectfont\color{popink}OnBuch\color{poporange}+}}
\newcommand{\signatureonbuch}{%
  \begin{tcolorbox}[enhanced,colback=popink,colframe=popink,boxrule=2.2pt,arc=10pt,
      drop shadow=poporange,left=14pt,right=14pt,top=10pt,bottom=10pt,before skip=0pt,after skip=0pt]
    \tikz[baseline=-0.7ex]\node[circle,fill=popgreen,draw=popcream,line width=1.3pt,minimum size=22pt,inner sep=0]{\popcheck{white}};%
    \hspace{9pt}\begin{minipage}[c]{0.62\linewidth}
      {\popxbold\fontsize{12}{13}\selectfont\color{popcream}Signé\ \textbullet\ {\poplogo OnBuch\color{poporange}+}}\par\vspace{2pt}
      {\raggedright\popsemi\fontsize{8}{10.5}\selectfont\color{popline}\MakeUppercase{Équipe pédagogique OnBuch+}\\\MakeUppercase{Conforme au programme officiel MINESEC}\par}
    \end{minipage}\hfill
    {\poplogo\fontsize{22}{22}\selectfont\color{popcream}OnBuch\color{poporange}+}
  \end{tcolorbox}%
}

% ============================================================
% Page de garde
% #1 titre · #2 matière · #3 niveau · #4 module · #5 n° leçon · #6 durée ·
% #7 description / objectifs (texte ou itemize)
% ============================================================
\newcommand{\pagegarde}[7]{%
  \thispagestyle{empty}
  \newgeometry{a4paper,margin=1.7cm,top=1.6cm,bottom=1.4cm}%
  \begin{tikzpicture}[remember picture,overlay]
    \fill[poporange!18] ([xshift=-3.2cm,yshift=-3.6cm]current page.north east) circle (4.6cm);
    \fill[poppurple!14] ([xshift=2.4cm,yshift=7.6cm]current page.south west) circle (3.8cm);
  \end{tikzpicture}%
  {\poplogo\fontsize{30}{30}\selectfont\color{popink}OnBuch\color{poporange}+}\hfill
  \raisebox{0.6ex}{\poppill{Cours signé \textbullet\ Premium}{popink}{popcream}}\par
  \vspace{1pt}\popsquiggle{poppurple}\par
  \vspace{8pt}
  \begin{tcolorbox}[enhanced,colback=poporange,colframe=popink,boxrule=2.8pt,arc=13pt,
      drop shadow=popink,left=18pt,right=18pt,top=12pt,bottom=13pt,after skip=10pt]
    \poppill{#2 \textbullet\ #3}{popink}{popcream}\hspace{5pt}\poppill{#4}{white}{popink}\par\vspace{8pt}
    {\popdisplay\fontsize{14}{14}\selectfont\color{popink}LEÇON #5}\par\vspace{4pt}
    {\raggedright\hyphenpenalty=10000\exhyphenpenalty=10000\popdisplay\fontsize{25}{27}\selectfont\color{popcream}\MakeUppercase{#1}\par}
    \vspace{8pt}
    {\popxbold\fontsize{9.5}{9.5}\selectfont\color{popink}\MakeUppercase{Durée conseillée : #6}}
  \end{tcolorbox}
  \begin{tcolorbox}[enhanced,colback=white,colframe=popink,boxrule=2.2pt,arc=10pt,
      drop shadow=popink,left=16pt,right=16pt,top=10pt,bottom=10pt,after skip=8pt]
    \poppill{Dans cette leçon}{poppurple}{white}\par\vspace{5pt}
    \popsemi\fontsize{9.3}{12.6}\selectfont\color{popink}#7
  \end{tcolorbox}
  \noindent\begin{minipage}[t]{0.487\linewidth}
    \begin{tcolorbox}[enhanced,colback=popgreen,colframe=popink,boxrule=2.2pt,arc=10pt,
        drop shadow=popink,left=11pt,right=11pt,top=10pt,bottom=10pt,height=3.1cm,valign=center]
      \tcbox[colback=white,colframe=popink,boxrule=1.3pt,arc=5pt,boxsep=0pt,nobeforeafter,
        left=3pt,right=3pt,top=3pt,bottom=3pt,valign=center]{\qrcode[height=1.9cm]{https://play.google.com/store/apps/details?id=com.luvvix.onbuch&hl=fr}}%
      \hspace{8pt}\begin{minipage}[c]{0.5\linewidth}
        {\popdisplay\fontsize{11}{12.5}\selectfont\color{white}\MakeUppercase{Télécharge OnBuch}}\par\vspace{4pt}
        {\raggedright\popsemi\fontsize{8.4}{11}\selectfont\color{white}Gratuit \textbullet\ Google Play\\Tuteur IA, annales, concours, résultats\par}
      \end{minipage}
    \end{tcolorbox}
  \end{minipage}\hfill
  \begin{minipage}[t]{0.487\linewidth}
    \begin{tcolorbox}[enhanced,colback=poppurple,colframe=popink,boxrule=2.2pt,arc=10pt,
        drop shadow=popink,left=11pt,right=11pt,top=10pt,bottom=10pt,height=3.1cm,valign=center]
      \tikz[baseline=-0.7ex]\node[circle,fill=white,draw=popink,line width=1.3pt,minimum size=1.6cm,inner sep=0]{\popdisplay\fontsize{24}{24}\selectfont\color{poppurple}+};%
      \hspace{8pt}\begin{minipage}[c]{0.6\linewidth}
        {\popdisplay\fontsize{11}{12.5}\selectfont\color{white}\MakeUppercase{Passe à OnBuch+}}\par\vspace{4pt}
        {\raggedright\popsemi\fontsize{8.4}{11}\selectfont\color{white}Tous les cours signés, Tuteur IA illimité, fiches PDF — onglet OnBuch+ de l'app\par}
      \end{minipage}
    \end{tcolorbox}
  \end{minipage}
  \vspace{8pt}
  \popcontenu
  \vfill
  \signatureonbuch
  \restoregeometry
  \newpage
}

% Rangée « Ce cours contient » (page de garde)
\newcommand{\poptile}[3]{% #1 couleur #2 gros texte #3 légende
  \begin{tcolorbox}[enhanced,colback=#1,colframe=popink,boxrule=2pt,arc=9pt,shadow={2.5pt}{-2.5pt}{0pt}{popink},
      left=6pt,right=6pt,top=9pt,bottom=9pt,halign=center,valign=center,height=1.75cm,width=\linewidth,nobeforeafter]
    {\popdisplay\fontsize{12.5}{13}\selectfont\color{white}\MakeUppercase{#2}}\par\vspace{3pt}
    {\centering\popsemi\fontsize{7.8}{9.5}\selectfont\color{white}#3\par}
  \end{tcolorbox}}
\newcommand{\popcontenu}{%
  \par\noindent{\popxboldls\fontsize{7.5}{7.5}\selectfont\color{popmuted}CE COURS SIGNÉ CONTIENT}\par\vspace{7pt}
  \noindent\begin{minipage}[t]{0.235\linewidth}\poptile{popblue}{Cours}{complet, illustré, pas à pas}\end{minipage}\hfill
  \begin{minipage}[t]{0.235\linewidth}\poptile{poporange}{Méthodes}{et exercices résolus}\end{minipage}\hfill
  \begin{minipage}[t]{0.235\linewidth}\poptile{poppink}{Exercices}{type examen + corrigés}\end{minipage}\hfill
  \begin{minipage}[t]{0.235\linewidth}\poptile{popgreen}{Fiche bilan}{l'essentiel à retenir}\end{minipage}\par}

% Sommaire
\newtcolorbox{sommaire}{enhanced,colback=white,colframe=popink,boxrule=2.2pt,arc=10pt,
  drop shadow=popink,left=18pt,right=18pt,top=13pt,bottom=13pt,before skip=6pt,after skip=20pt,
  before upper={\poppill{Sommaire}{poppurple}{white}\par\vspace{9pt}\popsemi\fontsize{10.6}{14.5}\selectfont\color{popink}}}

% Signature de fin de cours — poussée tout en bas de la dernière page
\newcommand{\findecours}{%
  \par\vfill\nopagebreak
  \begin{center}\popsquiggle{poporange}\end{center}\vspace{4pt}
  \signatureonbuch
}
\AtBeginDocument{\def\llbracket{\mathopen{[\mkern-3.5mu[}}\def\rrbracket{\mathclose{]\mkern-3.5mu]}}} % intervalles d'entiers (glyphe ⟦ absent de la police)
% Glyphes absents de Fira Math : redéfinis à partir de glyphes disponibles
\AtBeginDocument{%
  \def\setminus{\mathbin{\backslash}}\let\smallsetminus\setminus
  \def\vdots{\mathord{\vbox{\baselineskip3.2pt\lineskiplimit0pt\kern4pt\hbox{.}\hbox{.}\hbox{.}}}}%
  \def\ddots{\mathinner{\mkern1mu\raise6.5pt\hbox{.}\mkern2mu\raise3.6pt\hbox{.}\mkern2mu\raise0.7pt\hbox{.}\mkern1mu}}%
  \def\triangle{\mathord{\scalebox{0.9}{$\symup{\Delta}$}}}%
  \def\nabla{\mathord{\raisebox{\depth}{\scalebox{1}[-1]{$\symup{\Delta}$}}}}%
  \def\checkmark{\popcheck{popdark}}%
  \def\star{\mathbin{\ast}}\def\bigstar{\mathord{\ast}}%
  \def\diamond{\mathbin{\scalebox{0.6}{$\Diamond$}}}%
  \def\bigcup{\mathop{\mathchoice{\raisebox{-0.3em}{\scalebox{1.7}{$\cup$}}}{\scalebox{1.25}{$\cup$}}{\cup}{\cup}}\displaylimits}%
  \def\bigcap{\mathop{\mathchoice{\raisebox{-0.3em}{\scalebox{1.7}{$\cap$}}}{\scalebox{1.25}{$\cap$}}{\cap}{\cap}}\displaylimits}%
  \def\lesseqqgtr{\mathrel{\lessgtr}}\def\bigoplus{\mathop{\oplus}\displaylimits}%
}
\providecommand{\ssi}{si et seulement si} % abréviation fréquente
\AtBeginDocument{\providecommand{\wideparen}{\overparen}} % arc de cercle

%% FIGFIX-BEGIN v4 — correctifs globaux des figures (docs/onbuch-generation/tools/figfix)
\usepackage{adjustbox}\usepackage{etoolbox}
% toute figure TikZ/pgfplots plus large que la ligne est réduite à la largeur disponible
\BeforeBeginEnvironment{tikzpicture}{\begin{adjustbox}{max width=\linewidth}}
\AfterEndEnvironment{tikzpicture}{\end{adjustbox}}
% légendes pgfplots lisibles par-dessus les courbes
\pgfplotsset{every axis/.append style={legend style={fill=white,fill opacity=0.92,text opacity=1,draw=black!15,inner sep=3pt}}}
% étiquettes d'arêtes (syntaxe edge["..."]) avec fond blanc
\tikzset{every edge quotes/.append style={fill=white,inner sep=1.2pt}}
%% FIGFIX-END
\begin{document}

\pagegarde{Mesures, incertitudes et qualités d'un instrument de mesure}{Physique}{1ère D}{Module 1 : Mesures et incertitudes}{N01}{4 h}{Cette leçon initie l'élève de Première C à la démarche métrologique du physicien : caractériser la qualité d'un instrument de mesure (sensibilité, fidélité, justesse, précision), quantifier l'incertitude sur une mesure unique ou sur une série de mesures, et exprimer correctement un résultat sous la forme valeur ± incertitude avec un niveau de confiance.
\vspace{4pt}
\begin{multicols}{2}\small
\begin{itemize}
\item À la fin de cette leçon, je sais distinguer erreur et incertitude sur une mesure
\item À la fin de cette leçon, je sais définir et illustrer la sensibilité, la fidélité, la justesse et la précision d'un instrument de mesure
\item À la fin de cette leçon, je sais estimer l'incertitude sur une mesure unique à partir des caractéristiques de l'instrument
\item À la fin de cette leçon, je sais calculer la moyenne et l'écart-type expérimental d'une série de mesures
\item À la fin de cette leçon, je sais calculer l'incertitude-type d'une série de mesures répétées
\item À la fin de cette leçon, je sais déterminer l'incertitude élargie connaissant l'incertitude-type et le niveau de confiance
\end{itemize}\end{multicols}}

\begin{sommaire}
\begin{multicols}{2}
\begin{enumerate}
\item Mesure, erreur et incertitude
\item Qualités d'un instrument de mesure
\item Incertitude sur une mesure unique
\item Série de mesures : moyenne et écart-type
\item Incertitude-type, incertitude élargie et intervalle de confiance
\item Expression finale d'un résultat de mesure
\item Activité d'intégration
\item Méthodes et exercices résolus
\item Exercices
\item Corrigés des exercices
\item Fiche bilan
\end{enumerate}
\end{multicols}
\end{sommaire}

\begin{objectifs}
\begin{itemize}
\item À la fin de cette leçon, je sais distinguer erreur et incertitude sur une mesure
\item À la fin de cette leçon, je sais définir et illustrer la sensibilité, la fidélité, la justesse et la précision d'un instrument de mesure
\item À la fin de cette leçon, je sais estimer l'incertitude sur une mesure unique à partir des caractéristiques de l'instrument
\item À la fin de cette leçon, je sais calculer la moyenne et l'écart-type expérimental d'une série de mesures
\item À la fin de cette leçon, je sais calculer l'incertitude-type d'une série de mesures répétées
\item À la fin de cette leçon, je sais déterminer l'incertitude élargie connaissant l'incertitude-type et le niveau de confiance
\item À la fin de cette leçon, je sais construire un intervalle de confiance et l'interpréter
\item À la fin de cette leçon, je sais organiser des données expérimentales en tableau et en représentation graphique
\item À la fin de cette leçon, je sais exprimer un résultat sous la forme valeur ± incertitude avec un nombre raisonnable de chiffres significatifs
\end{itemize}
\end{objectifs}

\begin{prerequis}
\begin{itemize}
\item Notion de mesure d'une grandeur physique et unités du Système International (Seconde)
\item Chiffres significatifs et écriture scientifique d'un nombre
\item Calcul d'une moyenne arithmétique (mathématiques, collège)
\item Lecture de graduations sur une échelle (règle, thermomètre, cadran)
\item Notions élémentaires de représentation graphique (repère, nuage de points)
\end{itemize}
\end{prerequis}

\newpage

\coursec{Mesures, incertitudes et qualités d'un instrument de mesure}

Au marché Mokolo de Yaoundé, une maman achète 5\,kg de riz : le commerçant utilise une balance à ressort, et la cliente exige de revérifier sur une autre balance. Les deux affichent 4,95\,kg et 5,02\,kg ! Qui a raison ? De même, lors des contrôles de température à l'hôpital de district de Bafoussam, deux thermomètres donnent 38,4\,°C et 38,9\,°C pour le même patient. Comment un scientifique décide-t-il de la valeur à retenir et de la confiance à accorder à une mesure ? C'est tout l'objet de cette leçon.

\courssub{Grandeur physique, mesurage et mesurande}

\begin{definition}[Vocabulaire de base de la métrologie]
Une \cle{grandeur physique} est une propriété d'un phénomène, d'un corps ou d'une substance qui peut être qualifiée et quantifiée par un nombre et une unité (ex. : masse, température, longueur, tension électrique).
Le \cle{mesurage} est l'ensemble des opérations ayant pour but de déterminer une valeur de grandeur.
Le \cle{mesurande} est la grandeur particulière sujette au mesurage. Par exemple : « la masse de ce sac de ciment CIMENCAM » ou « la température axillaire de ce patient à cet instant ».
\end{definition}

Toute grandeur physique se mesure par comparaison à une référence : l'\cle{étalon}. Au Cameroun, les étalons nationaux sont conservés au Laboratoire National de Métrologie (LNM) rattaché à l'ANOR (Agence des Normes et de la Qualité), qui assure la traçabilité vers les étalons internationaux du BIPM (Bureau International des Poids et Mesures) à Sèvres.

\courssub{La valeur vraie : un idéal inaccessible}

\begin{definition}[Valeur vraie]
La \cle{valeur vraie} d'une grandeur est la valeur qui caractériserait parfaitement le mesurande dans les conditions de sa définition. C'est une notion idéale, inaccessible en pratique : on ne peut jamais la connaître avec certitude.
\end{definition}

\begin{propriete}[Inaccessibilité de la valeur vraie]
Toute mesure est entachée d'imperfections. Quelle que soit la qualité de l'instrument, du protocole ou de l'opérateur, il existe toujours un écart entre la valeur mesurée et la valeur vraie. Ce n'est pas une « faute » : c'est une propriété fondamentale de la mesure physique.
\end{propriete}

\begin{savaistu}[L'histoire du kilogramme]
Au Bureau International des Poids et Mesures (BIPM) à Sèvres, le kilogramme a été défini de 1889 à 2019 par l'\cle{étalon international du kilogramme} (le « Grand K »), un cylindre en platine-iridium conservé sous triple cloche. Or, des comparaisons périodiques ont montré que sa masse dérivait de quelques \emph{microgrammes} par rapport aux témoins nationaux ! Depuis 2019, le kilogramme est défini par la constante de Planck $h = 6,626\,070\,15 \times 10^{-34}\ \mathrm{J\cdot s}$, fixant la valeur vraie par une constante fondamentale de la nature, immuable et universelle.
\end{savaistu}

\courssub{Erreur de mesure : aléatoire vs systématique}

\begin{definition}[Erreur de mesure]
L'\cle{erreur de mesure} $E$ est la différence entre la valeur mesurée $m$ et la valeur vraie $V_{\text{vraie}}$ :
\[
E = m - V_{\text{vraie}}
\]
Comme $V_{\text{vraie}}$ est inconnue, l'erreur \textbf{ne peut jamais être connue exactement}.
\end{definition}

On distingue deux composantes de l'erreur :

\begin{itemize}
\item \textbf{Erreur systématique} : composante d'erreur qui, lors de répétitions du mesurage dans les mêmes conditions, reste constante ou varie de manière prévisible. Elle décale \emph{toutes} les mesures dans le même sens (ex. : balance mal tarée affichant +20\,g, thermomètre non étalonné lisant 0,5\,°C de trop).
\item \textbf{Erreur aléatoire} : composante d'erreur qui varie de manière imprévisible d'une mesure à l'autre. Elle provoque une \emph{dispersion} des résultats autour d'une valeur centrale (ex. : fluctuations de lecture, vibrations, bruit électronique).
\end{itemize}

\begin{popfigure}
\begin{tikzpicture}[scale=0.9]
% Cible 1 : erreurs aléatoires seulement (centré, dispersé)
\begin{scope}[xshift=-6cm]
\draw[popblue, thick] (0,0) circle (2.5);
\draw[popblue] (0,0) circle (1.8);
\draw[popblue] (0,0) circle (1.1);
\draw[popblue] (0,0) circle (0.4);
\fill[popblue] (0,0) circle (0.1);
\node[popblue] at (0,-3) {\textbf{Erreurs aléatoires seulement}};
\node[popblue] at (0,-3.5) {Dispersion autour du centre};
% Points dispersés autour du centre
\foreach \x/\y in {0.3/0.2, -0.4/0.5, 0.6/-0.3, -0.7/-0.6, 0.2/0.8, -0.5/-0.8, 0.8/0.1, -0.2/-0.4, 0.5/0.6, -0.6/0.3} {
    \fill[popblue] (\x,\y) circle (0.08);
}
\draw[->, popblue, thick] (0,0) -- (0,1.5) node[midway, right] {$E_{\text{aléat.}}$};
\end{scope}

% Cible 2 : erreur systématique seulement (groupé, décalé)
\begin{scope}[xshift=0cm]
\draw[poporange, thick] (0,0) circle (2.5);
\draw[poporange] (0,0) circle (1.8);
\draw[poporange] (0,0) circle (1.1);
\draw[poporange] (0,0) circle (0.4);
\fill[poporange] (0,0) circle (0.1);
\node[poporange] at (0,-3) {\textbf{Erreur systématique seulement}};
\node[poporange] at (0,-3.5) {Groupe décalé du centre};
% Points groupés mais décalés
\foreach \x/\y in {1.2/0.3, 1.3/0.4, 1.1/0.2, 1.4/0.5, 1.2/0.1, 1.3/0.3, 1.1/0.4, 1.2/0.2} {
    \fill[poporange] (\x,\y) circle (0.08);
}
\draw[->, poporange, thick] (0,0) -- (1.2,0.3) node[midway, above] {$E_{\text{syst.}}$};
\draw[dashed, popmuted] (0,0) -- (1.2,0.3);
\end{scope}

% Cible 3 : les deux (décalé ET dispersé)
\begin{scope}[xshift=6cm]
\draw[poppurple, thick] (0,0) circle (2.5);
\draw[poppurple] (0,0) circle (1.8);
\draw[poppurple] (0,0) circle (1.1);
\draw[poppurple] (0,0) circle (0.4);
\fill[poppurple] (0,0) circle (0.1);
\node[poppurple] at (0,-3) {\textbf{Erreurs systématiques + aléatoires}};
\node[poppurple] at (0,-3.5) {Dispersé ET décalé};
% Points dispersés et décalés
\foreach \x/\y in {1.5/0.6, 1.2/0.8, 1.8/0.2, 1.0/0.4, 1.6/0.9, 1.3/0.1, 1.7/0.5, 1.1/0.7, 1.4/0.3, 1.6/0.6} {
    \fill[poppurple] (\x,\y) circle (0.08);
}
\draw[->, poppurple, thick] (0,0) -- (1.4,0.5) node[midway, above] {$E_{\text{totale}}$};
\draw[dashed, popmuted] (0,0) -- (1.4,0.5);
\end{scope}
\end{tikzpicture}
\legende{Analogie de la cible de tir : le centre représente la valeur vraie. (Gauche) Erreurs aléatoires : dispersion symétrique. (Milieu) Erreur systématique : groupement décalé. (Droite) Cas réel : les deux effets se combinent.}
\end{popfigure}

\begin{attention}[Erreur vs Incertitude — Piège classique]
Ne confondez jamais :
\begin{itemize}
\item \textbf{Erreur} : écart (inconnu) entre valeur mesurée et valeur vraie. C'est une \emph{grandeur physique} qui existe mais qu'on ne peut pas connaître.
\item \textbf{Incertitude} : estimation numérique (connue) de la dispersion des valeurs qu'on peut raisonnablement attribuer au mesurande. C'est un \emph{paramètre d'évaluation} qu'on calcule ou qu'on estime.
\end{itemize}
Dire « l'erreur est de 0,2\,kg » est \textbf{faux} (on ne la connaît pas). Dire « l'incertitude est de 0,2\,kg » est \textbf{correct} (c'est notre estimation du doute).
\end{attention}

\courssub{Incertitude de mesure : quantifier le doute}

\begin{definition}[Incertitude de mesure]
L'\cle{incertitude de mesure} est un paramètre, associé au résultat d'un mesurage, qui caractérise la dispersion des valeurs qu'on peut raisonnablement attribuer au mesurande. Elle s'exprime sous la forme d'un intervalle centré sur la valeur mesurée (ou la moyenne d'une série).
\end{definition}

\begin{propriete}[Distinction fondamentale]
\begin{itemize}
\item L'erreur est \textbf{inconnue} (valeur vraie inaccessible).
\item L'incertitude s'\textbf{estime} et s'\textbf{exprime} (c'est notre connaissance de la qualité de la mesure).
\end{itemize}
Une mesure sans incertitude n'a pas de sens scientifique. Une « mesure parfaite » (incertitude nulle) n'existe pas.
\end{propriete}

\courssub{Sources d'incertitude}

Toute incertitude provient de quatre grandes familles de causes (norme NF ISO/CEI Guide 98-3) :

\begin{center}
\begin{tabularx}{\linewidth}{|l|X|}\hline
\textbf{Source} & \textbf{Exemples concrets (contexte camerounais)} \\ \hline
\rowcolor{popblueL}
Instrument & Résolution limitée (balance à 10\,g près), dérive d'étalonnage (thermomètre non vérifié depuis 2 ans), non-linéarité, effet de charge (voltmètre qui perturbe le circuit) \\ \hline
Opérateur & Erreur de lecture (parallaxe sur aiguille), temps de réaction (chronomètre), fatigue, biais de confirmation \\ \hline
Méthode & Modèle imparfait (négliger la résistance interne d'une pile), approximation de formule, hypothèses non vérifiées \\ \hline
Environnement & Température (dilatation de la règle, dérive du capteur), humidité, vibrations (pont Wouri), champs électromagnétiques (lignes ENEO haute tension), pression atmosphérique \\ \hline
\end{tabularx}
\end{center}

\begin{experience}[Identifier les sources d'incertitude]
\textbf{Matériel} : Balance électronique de cuisine (précision 1\,g), règle graduée (mm), chronomètre smartphone, thermomètre à mercure.
\textbf{Protocole} : Mesurer la masse d'un œuf, la longueur d'un crayon, la durée de chute d'une balle de ping-pong sur 2\,m, la température d'un verre d'eau.
\textbf{Observations} : Noter pour chaque mesure : résolution de l'instrument, fluctuations observées, conditions ambiantes (T, humidité), difficultés de lecture.
\textbf{Interprétation} : Classer chaque source identifiée dans le tableau ci-dessus. Discuter : laquelle domine ? Peut-on la réduire ?
\end{experience}

\courssub{Écriture d'un résultat de mesure}

\begin{definition}[Écriture standardisée]
Le résultat d'une mesure s'exprime sous la forme :
\[
\boxed{m = m_0 \pm \Delta m\ \text{(unité)}}
\]
où :
\begin{itemize}
\item $m_0$ est la \cle{valeur mesurée} (ou la moyenne d'une série) — meilleure estimation de la valeur vraie.
\item $\Delta m$ est l'\cle{incertitude élargie} (souvent notée $U$) — demi-largeur de l'intervalle de confiance.
\item L'unité est obligatoire et doit être cohérente (SI de préférence).
\end{itemize}
\end{definition}

\begin{exemplebox}[Masse d'un sac de ciment CIMENCAM]
Un sac de ciment est pesé sur une balance industrielle étalonnée. L'affichage donne $m_0 = 50,0\ \mathrm{kg}$. Le certificat d'étalonnage indique une incertitude élargie $U = 0,2\ \mathrm{kg}$ pour un niveau de confiance de 95\,\%.
On écrit :
\[
m = 50,0 \pm 0,2\ \mathrm{kg}
\]
Ce qui signifie : \textbf{la valeur vraie de la masse du sac se situe avec une probabilité de 95\,\% dans l'intervalle}
\[
49,8\ \mathrm{kg} \le m \le 50,2\ \mathrm{kg}
\]
Cet intervalle s'appelle \cle{intervalle de confiance} (niveau 95\,\%).
\end{exemplebox}

\begin{popfigure}
\begin{tikzpicture}[node distance=1.2cm, >=Stealth]
% Nœuds
\node[draw, rounded corners, fill=popblueL, thick, minimum width=3cm, minimum height=1.2cm, align=center] (vraie){\textbf{Valeur vraie}\\(inconnue, idéale)};
\node[draw, rounded corners, fill=poporangeL, thick, minimum width=3.5cm, minimum height=1.2cm, right=of vraie, align=center] (mesurage){\textbf{Mesurage}\\(instrument + protocole + opérateur + environnement)};
\node[draw, rounded corners, fill=popgreenL, thick, minimum width=4cm, minimum height=1.5cm, right=of mesurage, align=center] (resultat){\textbf{Résultat exprimé}\\($m_0 \pm \Delta m$ avec unité)\\\textit{Intervalle de confiance}};

% Flèches
\draw[->, thick, popdark] (vraie.east) -- node[above] {\small{inaccessible}} (mesurage.west);
\draw[->, thick, popdark] (mesurage.east) -- node[above] {\small{produit}} (resultat.west);

% Annotations sous les nœuds
\node[below=0.3cm of vraie, text width=3.5cm, align=center, font=\small] {Cible idéale\\« Grand K », constante $h$, définition SI};
\node[below=0.3cm of mesurage, text width=4cm, align=center, font=\small] {Sources d'erreur :\newline systématiques + aléatoires};
\node[below=0.3cm of resultat, text width=4.5cm, align=center, font=\small] {Encadrement quantifié :\newline $m_0 - \Delta m \le V_{\text{vraie}} \le m_0 + \Delta m$};
\end{tikzpicture}
\legende{Chaîne métrologique : de la valeur vraie (inaccessible) au résultat exprimé avec son incertitude. Le mesurage introduit des erreurs ; l'incertitude quantifie notre ignorance résiduelle.}
\end{popfigure}

\courssub{Règles de présentation (chiffres significatifs)}

\begin{propriete}[Règles d'écriture $m_0 \pm \Delta m$]
\begin{enumerate}
\item L'incertitude $\Delta m$ s'exprime avec \textbf{1 ou 2 chiffres significatifs} (généralement 1, sauf si le premier chiffre est 1 ou 2, où on en garde 2).
\item La valeur mesurée $m_0$ s'arrondit à la \textbf{même décimale} que l'incertitude.
\item L'unité s'écrit une seule fois, après l'incertitude.
\end{enumerate}
\end{propriete}

\begin{exoresolu}[Écriture correcte d'un résultat]
\textbf{Énoncé.} Un élève mesure la longueur d'un fil de cuivre avec une règle au millimètre. Il obtient $L_0 = 1,234\ \mathrm{m}$ et estime son incertitude à $\Delta L = 0,003\ \mathrm{m}$.
\textbf{Solution.}
\begin{itemize}
\item $\Delta L = 0,003\ \mathrm{m}$ a 1 chiffre significatif (le 3).
\item $L_0$ doit s'arrondir au millimètre (même décimale) : $L_0 = 1,234\ \mathrm{m} \to 1,234\ \mathrm{m}$ (déjà au mm).
\item Résultat : $\boxed{L = 1,234 \pm 0,003\ \mathrm{m}}$.
\end{itemize}
\textbf{Variante.} Si $\Delta L = 0,0015\ \mathrm{m}$ (2 chiffres significatifs car commence par 1), on arrondit $\Delta L$ à $0,0015\ \mathrm{m}$ et $L_0$ au dixième de mm : $L = 1,2340 \pm 0,0015\ \mathrm{m}$.
\end{exoresolu}

\begin{attention}[Pièges de rédaction au Probatoire]
\begin{itemize}
\item \textbf{Oublier l'unité} : $m = 50,0 \pm 0,2$ \textcolor{red}{$\times$} $\to$ $m = 50,0 \pm 0,2\ \mathrm{kg}$ \textcolor{popgreen}{$\checkmark$}
\item \textbf{Décimales incohérentes} : $L = 1,23 \pm 0,005\ \mathrm{m}$ \textcolor{red}{$\times$} (valeur au cm, incertitude au mm) $\to$ $L = 1,230 \pm 0,005\ \mathrm{m}$ \textcolor{popgreen}{$\checkmark$}
\item \textbf{Trop de chiffres sur l'incertitude} : $\Delta m = 0,1234\ \mathrm{kg}$ \textcolor{red}{$\times$} $\to$ $\Delta m = 0,12\ \mathrm{kg}$ \textcolor{popgreen}{$\checkmark$}
\item \textbf{Confondre $\pm$ et $\mp$} : le symbole $\pm$ signifie « plus ou moins » (intervalle symétrique).
\end{itemize}
\end{attention}

\begin{methode}[Exprimer un résultat de mesure]
\begin{enumerate}
\item Identifier le mesurande et les conditions de mesure.
\item Choisir l'instrument adapté (résolution, domaine).
\item Effectuer le mesurage (une fois ou en série).
\item Estimer l'incertitude (voir sections 3 à 5).
\item Arrondir l'incertitude à 1 ou 2 chiffres significatifs.
\item Arrondir la valeur mesurée à la même décimale.
\item Écrire : $m = m_0 \pm \Delta m\ \text{(unité SI)}$.
\item Interpréter : donner l'intervalle de confiance et son niveau (généralement 95\,\%).
\end{enumerate}
\end{methode}

\begin{exemplebox}[Température à l'hôpital de Bafoussam]
Deux thermomètres numériques (résolution 0,1\,°C) donnent $T_1 = 38,4\ \mathrm{°C}$ et $T_2 = 38,9\ \mathrm{°C}$ pour le même patient.
\begin{itemize}
\item L'écart $|T_1 - T_2| = 0,5\ \mathrm{°C}$ révèle une différence \textbf{bien supérieure} à la résolution (0,1\,°C).
\item Cela signale une \textbf{erreur systématique} probable sur l'un des deux (dérive d'étalonnage, mauvais contact, pile faible).
\item Le protocole médical exige : \textbf{ne pas décider sur une seule mesure}, utiliser un thermomètre de référence étalonné, ou prendre une série de mesures pour estimer l'incertitude.
\item Résultat correct : $T = 38,7 \pm 0,3\ \mathrm{°C}$ (moyenne $\pm$ demi-écart) \emph{si} on n'a pas d'étalon, mais avec la mention « incertitude dominée par erreur systématique non corrigée ».
\end{itemize}
\end{exemplebox}

\begin{savaistu}[Métrologie légale au Cameroun]
L'ANOR (Agence des Normes et de la Qualité) impose la \textbf{vérification périodique} des instruments de mesure utilisés dans les transactions commerciales (balances de marché, compteurs d'eau SONATREL/ENEO, pompes à essence, taximètres). Un instrument non vérifié ou non conforme est saisi. C'est l'application concrète de la métrologie : garantir l'équité commerciale et la sécurité sanitaire par la \textbf{traçabilité} des mesures vers les étalons nationaux et internationaux.
\end{savaistu}

\coursec{Qualités d'un instrument de mesure}

Au-delà de l'incertitude, un instrument se juge à quatre qualités métrologiques fondamentales. Un bon instrument pour le laboratoire de physique du lycée de Ngaoundéré ou pour le contrôle qualité à la SONARA doit les satisfaire.

\courssub{Sensibilité}

\begin{definition}[Sensibilité]
La \cle{sensibilité} d'un instrument est le rapport de la variation de l'indication (sortie) à la variation correspondante de la grandeur mesurée (entrée). Pour un instrument à réponse linéaire : $S = \frac{\Delta I}{\Delta M}$.
\end{definition}

\begin{propriete}[Interprétation]
\begin{itemize}
\item Une \textbf{grande sensibilité} signifie qu'une petite variation du mesurande produit une grande variation de l'indication (ex. : galvanomètre à miroir, balance analytique). C'est souhaitable pour détecter de faibles signaux.
\item Une \textbf{faible sensibilité} signifie qu'il faut une grande variation du mesurande pour voir bouger l'aiguille (ex. : balance de marché pour sacs de 50\,kg).
\item La sensibilité ne garantit pas la justesse : un instrument très sensible peut être mal réglé (erreur systématique).
\end{itemize}
\end{propriete}

\begin{exemplebox}[Sensibilité d'un voltmètre analogique]
Un voltmètre à aiguille possède une échelle de 100 divisions pour une tension maximale de 10\,V. Son indication $I$ (en divisions) vaut $I = 10 \times U$ (U en V). Sa sensibilité est $S = 10\ \mathrm{div/V}$.
Si on change l'échelle pour 100\,V max (100 divisions), $S = 1\ \mathrm{div/V}$ : l'instrument est 10 fois moins sensible.
\end{exemplebox}

\courssub{Fidélité (ou répétabilité)}

\begin{definition}[Fidélité]
La \cle{fidélité} (ou \cle{répétabilité}) est l'aptitude d'un instrument à donner des indications voisines pour des répliques successives du mesurage, effectuées dans les \textbf{mêmes conditions} (même opérateur, même instrument, même lieu, court intervalle de temps).
\end{definition}

\begin{propriete}[Lien avec l'erreur aléatoire]
La fidélité caractérise la dispersion due aux \textbf{erreurs aléatoires} seulement. Elle s'évalue par l'écart-type expérimental $s$ (ou l'incertitude-type $u_A$) sur une série de mesures dans des conditions de répétabilité.
\begin{itemize}
\item Instrument \textbf{fidèle} : points groupés (faible dispersion).
\item Instrument \textbf{peu fidèle} : points dispersés (bruit, instabilité, frottements).
\end{itemize}
\end{propriete}

\courssub{Justesse (ou exactitude)}

\begin{definition}[Justesse]
La \cle{justesse} est la proximité entre la moyenne d'un grand nombre de résultats de mesure et la \textbf{valeur vraie} (ou une valeur de référence acceptée comme vraie). Elle caractérise l'absence d'\textbf{erreur systématique}.
\end{definition}

\begin{propriete}[Lien avec l'erreur systématique]
\begin{itemize}
\item Instrument \textbf{juste} : le centre du groupement coïncide avec la valeur vraie (erreur systématique nulle ou corrigée).
\item Instrument \textbf{peu juste} : le groupement est décalé (biais, dérive d'étalonnage, zéro défectueux).
\end{itemize}
\end{propriete}

\begin{attention}[Justesse vs Fidélité — Le piège de la cible]
\begin{itemize}
\item Un instrument peut être \textbf{fidèle mais pas juste} : groupement serré mais décalé (erreur systématique non corrigée). C'est le cas le plus fréquent en TP : on corrige la justesse par un étalonnage.
\item Un instrument peut être \textbf{juste mais pas fidèle} : moyenne correcte mais grande dispersion (bruit important). La moyenne de nombreuses mesures redonne la valeur vraie.
\item L'idéal : \textbf{juste ET fidèle}.
\end{itemize}
\end{attention}

\courssub{Précision}

\begin{definition}[Précision]
La \cle{précision} est l'aptitude à donner des résultats proches les uns des autres \textbf{ET} proches de la valeur vraie. Elle combine \textbf{fidélité} et \textbf{justesse}.
\end{definition}

\begin{propriete}[Incertitude et précision]
L'incertitude de mesure est l'indicateur quantitatif de la \textbf{précision}. Une incertitude faible signifie une haute précision.
\[
\text{Précision} \approx \frac{1}{\text{Incertitude}}
\]
Attention : dans le langage courant, « précis » veut souvent dire « sensible » (ex. : « une montre précise au 1/100e de seconde »). En métrologie rigoureuse, une montre qui avance de 5 min par jour est \emph{sensible} au 1/100e s mais \emph{peu juste} et donc \emph{peu précise}.
\end{propriete}

\begin{popfigure}
\begin{tikzpicture}[scale=0.85]
% Quatre cibles
\def\cible#1#2#3#4{
\begin{scope}[xshift=#1, yshift=#2]
\draw[popline, thick] (0,0) circle (2);
\draw[popline] (0,0) circle (1.4);
\draw[popline] (0,0) circle (0.8);
\draw[popline] (0,0) circle (0.2);
\fill[popdark] (0,0) circle (0.08) node[below=0.3cm] {Vraie};
\foreach \ptx/\pty in {#3} { \fill[popblue] (\ptx,\pty) circle (0.07); }
\node[popdark, font=\bfseries\small, align=center] at (0,-2.5) {#4};
\end{scope}
}
% Juste & Fidèle
\cible{-7.5}{0}{0.1/0.1, -0.1/-0.1, 0.05/0.15, -0.05/-0.15, 0.15/0.05, -0.15/-0.05}{Juste \& Fidèle\\ (Précis)}
% Juste & Pas Fidèle
\cible{-1.5}{0}{0.8/0.5, -0.7/-0.6, 1.0/-0.4, -0.9/0.7, 0.6/0.9, -0.5/-0.8}{Juste \& Pas Fidèle\\ (Moyenne vraie, dispersé)}
% Pas Juste & Fidèle
\cible{4.5}{0}{1.5/0.3, 1.6/0.4, 1.4/0.2, 1.7/0.5, 1.5/0.1, 1.6/0.3}{Pas Juste \& Fidèle\\ (Biais systématique)}
% Pas Juste & Pas Fidèle
\cible{10.5}{0}{1.2/0.8, 1.8/-0.5, 0.9/0.6, 1.5/1.0, 1.0/-0.7, 1.7/0.2}{Pas Juste \& Pas Fidèle\\ (Imprécis)}
\end{tikzpicture}
\legende{Les quatre combinaisons Justesse / Fidélité. Seule la première (en haut à gauche) correspond à un instrument \textbf{précis} (faible incertitude).}
\end{popfigure}

\begin{experience}[Tester les qualités d'un instrument]
\textbf{Matériel} : 3 balances électroniques (marques différentes, même portée 200\,g, résolution 0,01\,g), 1 masse étalon de 100\,g (classe F1 ou F2), 10 pièces de 100\,F CFA (masses voisines).
\textbf{Protocole} :
\begin{enumerate}
\item Peser 10 fois la masse étalon sur chaque balance $\to$ évaluer la \textbf{fidélité} (écart-type $s$).
\item Comparer la moyenne $\bar{m}$ à 100,00\,g $\to$ évaluer la \textbf{justesse} (biais $= \bar{m} - 100,00$).
\item Peser les 10 pièces sur la meilleure balance $\to$ vérifier la \textbf{sensibilité} (les masses diffèrent-elles de plus que la résolution ?).
\end{enumerate}
\textbf{Interprétation} : Classer les balances. Laquelle choisir pour de la chimie analytique ? Pour de la vente au détail ?
\end{experience}

\coursec{Incertitude sur une mesure unique (Évaluation de type B)}

Lorsqu'on ne peut pas répéter la mesure (coût, temps, destructif), on estime l'incertitude par une \textbf{évaluation de type B} : analyse des spécifications constructeur, certificat d'étalonnage, expérience antérieure, résolution.

\courssub{Incertitude due à la résolution (instrument numérique)}

\begin{propriete}[Règle de la résolution numérique]
Pour un afficheur numérique de résolution (pas) $d$ (ex. : 0,1\,g, 1\,mm, 0,01\,V), l'erreur d'arrondi suit une distribution \textbf{rectangulaire} (uniforme) de demi-largeur $a = d/2$.
L'incertitude-type associée est :
\[
\boxed{u_{\text{résol}} = \frac{d}{2\sqrt{3}} = \frac{d}{\sqrt{12}}}
\]
\textbf{Démonstration (exigible au Probatoire) :} Pour une loi uniforme sur $[-a, a]$, la variance est $\sigma^2 = \frac{a^2}{3}$. Ici $a = d/2$, donc $u = \frac{d/2}{\sqrt{3}} = \frac{d}{\sqrt{12}}$.
\end{propriete}

\begin{exemplebox}[Multimètre numérique en TP]
Un multimètre affiche $U = 4,57\ \mathrm{V}$ sur le calibre 20\,V (résolution $d = 0,01\ \mathrm{V}$).
\[
u_{\text{résol}} = \frac{0,01}{\sqrt{12}} \approx 0,0029\ \mathrm{V}
\]
On arrondit l'incertitude à 1 chiffre significatif : $u_{\text{résol}} \approx 0,003\ \mathrm{V}$.
La tension s'écrit $U = 4,570 \pm 0,003\ \mathrm{V}$ (justesse non vérifiée ici).
\end{exemplebox}

\courssub{Incertitude due à la lecture (instrument analogique)}

\begin{propriete}[Règle de la lecture analogique]
Pour un instrument à aiguille (règle, ampèremètre analogique), l'erreur de lecture (parallaxe, interpolation entre graduations) est souvent modélisée par une distribution rectangulaire de demi-largeur $a = \frac{\text{pas de graduation}}{2}$ (estimation optimiste) ou $a = \text{pas de graduation}$ (estimation prudente).
\[
u_{\text{lect}} = \frac{a}{\sqrt{3}}
\]
Souvent, par convention simplifiée au lycée : $u_{\text{lect}} \approx \frac{\text{pas}}{2}$ (demi-pas) ou $\frac{\text{pas}}{5}$ (interpolation visuelle). \textbf{Précisez toujours votre hypothèse.}
\end{propriete}

\begin{exemplebox}[Règle graduée au millimètre]
On mesure une longueur $L = 12,3\ \mathrm{cm}$. Le pas est 1\,mm = 0,1\,cm.
Estimation prudente : $a = 0,1\ \mathrm{cm}$ (on ne lit pas mieux que le trait).
$u_{\text{lect}} = \frac{0,1}{\sqrt{3}} \approx 0,058\ \mathrm{cm} \to 0,06\ \mathrm{cm}$.
Résultat : $L = 12,3 \pm 0,06\ \mathrm{cm}$.
\end{exemplebox}

\courssub{Incertitude due à l'étalonnage (certificat)}

\begin{propriete}[Incertitude d'étalonnage]
Le certificat d'étalonnage donne une \textbf{incertitude élargie} $U_{\text{étal}}$ pour un niveau de confiance (souvent 95\,\%, $k=2$).
L'incertitude-type correspondante est :
\[
u_{\text{étal}} = \frac{U_{\text{étal}}}{k}
\]
Si $k$ n'est pas indiqué, on suppose $k=2$ (loi normale) ou $k=\sqrt{3}$ (loi rectangulaire) selon le contexte.
\end{propriete}

\courssub{Autres sources (Type B) : température, alimentation, etc.}

Chaque influence (température, tension d'alimentation, vieillissement) est modélisée par une distribution de probabilité (souvent rectangulaire par défaut si on ne connaît que les bornes $\pm a$) et convertie en incertitude-type $u_i = a/\sqrt{3}$.

\begin{methode}[Bilan d'incertitude Type B (mesure unique)]
\begin{enumerate}
\item Lister toutes les sources d'incertitude identifiées (résolution, lecture, étalonnage, température, etc.).
\item Pour chaque source, estimer une demi-largeur $a_i$ (borne d'erreur maximale plausible) et choisir une distribution (rectangulaire par défaut $\to$ diviseur $\sqrt{3}$ ; normale si certificat $\to$ diviseur $k$ ; triangulaire $\to$ diviseur $\sqrt{6}$).
\item Calculer l'incertitude-type $u_i = a_i / \text{diviseur}$ pour chaque source.
\item Combiner les incertitudes-types (somme quadratique si sources indépendantes) :
\[
u_c = \sqrt{\sum u_i^2}
\]
\item Calculer l'incertitude élargie $U = k \cdot u_c$ (généralement $k=2$ pour 95\,\%).
\item Exprimer le résultat final.
\end{enumerate}
\end{methode}

\begin{exoresolu}[Bilan Type B : Mesure d'une résistance au multimètre]
\textbf{Énoncé.} On mesure $R = 100,5\ \Omega$ sur un multimètre (calibre 200\,$\Omega$, résolution 0,1\,$\Omega$). Certificat : erreur max $\pm 0,2\ \Omega$ ($k=2$). Température labo : $23 \pm 3\ \mathrm{°C}$. Coefficient de température : $100\ \mathrm{ppm/°C}$.
\textbf{Solution.}
\begin{itemize}
\item Résolution : $d=0,1\ \Omega \to u_1 = 0,1/\sqrt{12} = 0,029\ \Omega$.
\item Étalonnage : $U=0,2\ \Omega, k=2 \to u_2 = 0,2/2 = 0,10\ \Omega$.
\item Température : $\Delta T = \pm 3\ \mathrm{°C} \to \Delta R = 100,5 \times 100\times10^{-6} \times 3 = 0,030\ \Omega$. Loi rectangulaire $a=0,030 \to u_3 = 0,030/\sqrt{3} = 0,017\ \Omega$.
\item Combinée : $u_c = \sqrt{0,029^2 + 0,10^2 + 0,017^2} = \sqrt{0,00084 + 0,0100 + 0,00029} = \sqrt{0,0111} \approx 0,105\ \Omega$.
\item Élargie ($k=2$) : $U = 2 \times 0,105 = 0,21\ \Omega \to 0,2\ \Omega$ (1 ch. sig.).
\item Résultat : $\boxed{R = 100,5 \pm 0,2\ \Omega}$.
\end{itemize}
\textbf{Note :} L'étalonnage domine ($u_2=0,10$). Améliorer la résolution ne changerait presque rien.
\end{exoresolu}

\coursec{Série de mesures : moyenne et écart-type (Évaluation de type A)}

Lorsque le mesurage est répétable, on effectue une \textbf{série de $n$ mesures} ($n \ge 5$, idéalement $10 \dots 20$) pour évaluer la composante aléatoire par des méthodes statistiques (Type A).

\courssub{Moyenne arithmétique : meilleure estimation}

\begin{definition}[Moyenne d'une série]
Pour $n$ mesures $x_1, x_2, \dots, x_n$ du même mesurande dans les mêmes conditions :
\[
\boxed{\bar{x} = \frac{1}{n} \sum_{i=1}^n x_i}
\]
La moyenne $\bar{x}$ est la meilleure estimation de la valeur vraie (si les erreurs systématiques sont négligeables ou corrigées).
\end{definition}

\courssub{Écart-type expérimental : dispersion de la série}

\begin{definition}[Écart-type expérimental (échantillon)]
L'\cle{écart-type expérimental} $s$ (ou $\sigma_{n-1}$) quantifie la dispersion des valeurs individuelles autour de la moyenne :
\[
\boxed{s = \sqrt{\frac{1}{n-1} \sum_{i=1}^n (x_i - \bar{x})^2}}
\]
Le dénominateur $n-1$ (degrés de liberté) donne un estimateur non biaisé de l'écart-type de la population.
\end{definition}

\begin{propriete}[Écart-type de la moyenne (Incertitude-type de type A)]
La dispersion de la \textbf{moyenne} est plus faible que celle des valeurs individuelles. L'incertitude-type due aux effets aléatoires (Type A) est :
\[
\boxed{u_A(\bar{x}) = \frac{s}{\sqrt{n}}}
\]
C'est l'incertitude-type sur la meilleure estimation $\bar{x}$.
\end{propriete}

\begin{exemplebox}[Série de mesures : période d'un pendule]
Un élève mesure 10 périodes d'un pendule simple (longueur $\approx 1\ \mathrm{m}$) au chronomètre (résolution 0,01\,s).
Valeurs (en s) : 2,01 ; 2,03 ; 2,00 ; 2,02 ; 2,04 ; 1,99 ; 2,02 ; 2,01 ; 2,03 ; 2,00.
\begin{itemize}
\item $n=10$.
\item $\bar{T} = \frac{20,15}{10} = 2,015\ \mathrm{s}$.
\item Calcul de $s$ :
$\sum (T_i - \bar{T})^2 = (-0,005)^2 + (0,015)^2 + \dots = 0,00225\ \mathrm{s}^2$.
$s = \sqrt{\frac{0,00225}{9}} = \sqrt{0,00025} = 0,0158\ \mathrm{s}$.
\item Incertitude-type Type A : $u_A = \frac{0,0158}{\sqrt{10}} = 0,0050\ \mathrm{s}$.
\item Incertitude résolution (Type B) : $d=0,01\ \mathrm{s} \to u_{\text{résol}} = 0,01/\sqrt{12} = 0,0029\ \mathrm{s}$.
\item Incertitude temps de réaction (Type B, estimation prudente $\pm 0,1\ \mathrm{s}$ rectangulaire) : $u_{\text{réac}} = 0,1/\sqrt{3} = 0,058\ \mathrm{s}$.
\end{itemize}
\textbf{Analyse :} $u_{\text{réac}}$ domine largement. La série n'améliore pas l'incertitude totale si l'erreur systématique (réaction) n'est pas réduite (ex. : capteur photogate).
\end{exemplebox}

\begin{attention}[Nombre de mesures $n$]
\begin{itemize}
\item $n < 5$ : estimation de $s$ très peu fiable (incertitude sur $s$ elle-même > 20\,\%).
\item $n = 10$ : bon compromis temps/précision (incertitude sur $s \approx 15\%$).
\item $n \ge 30$ : loi normale bien approchée, $k=2$ exact pour 95\,\%. Au lycée, $n=10$ est la norme.
\end{itemize}
\end{attention}

\coursec{Incertitude-type, incertitude élargie et intervalle de confiance}

\courssub{Incertitude-type composée $u_c$}

Lorsque plusieurs sources (Type A et Type B) contribuent, on combine leurs incertitudes-types $u_i$ (supposées indépendantes) par la \textbf{somme quadratique} (propagation des variances) :

\begin{propriete}[Incertitude-type composée]
\[
\boxed{u_c = \sqrt{\sum_{i=1}^N u_i^2} = \sqrt{u_A^2 + u_{B1}^2 + u_{B2}^2 + \dots}}
\]
\textbf{Condition d'application} : les sources d'incertitude doivent être \textbf{indépendantes} (non corrélées). Si corrélées, il faut ajouter des termes de covariance (hors programme 1ère C).
\end{propriete}

\courssub{Incertitude élargie $U$ et facteur de couverture $k$}

\begin{definition}[Incertitude élargie]
L'\cle{incertitude élargie} $U$ est obtenue en multipliant l'incertitude-type composée $u_c$ par un \cle{facteur de couverture} $k$ :
\[
\boxed{U = k \cdot u_c}
\]
Le facteur $k$ est choisi pour obtenir le \textbf{niveau de confiance} souhaité (probabilité que la valeur vraie soit dans l'intervalle).
\end{definition}

\begin{propriete}[Facteur de couverture $k$ usuel]
\begin{itemize}
\item \textbf{Niveau de confiance 95\,\%} (standard industriel et scientifique) :
\begin{itemize}
\item Si $n$ grand ($>30$) ou incertitudes Type B dominantes (loi normale) : $k = 2$ (exact : 1,96).
\item Si $n$ petit (ex. $n=10$) et Type A dominant (loi de Student) : $k = t_{0,975}(n-1)$. Pour $n=10$, $k \approx 2,26$.
\end{itemize}
\item \textbf{Niveau de confiance 68\,\%} (1 sigma) : $k = 1$.
\item \textbf{Au Probatoire} : sauf précision contraire, on utilise \textbf{$k=2$} pour un niveau de confiance de 95\,\%.
\end{itemize}
\end{propriete}

\courssub{Intervalle de confiance}

\begin{definition}[Intervalle de confiance]
L'\cle{intervalle de confiance} au niveau $P$ (ex. 95\,\%) est l'intervalle :
\[
\boxed{[\bar{x} - U\ ;\ \bar{x} + U] \quad \text{ou} \quad \bar{x} \pm U}
\]
\textbf{Interprétation correcte} : « La valeur vraie du mesurande a une probabilité de 95\,\% de se situer dans cet intervalle » (ou : « Si on répétait l'expérience un grand nombre de fois, 95\,\% des intervalles calculés ainsi contiendraient la valeur vraie »).
\end{definition}

\begin{attention}[Interdits de formulation]
\begin{itemize}
\item \textcolor{red}{$\times$} « L'erreur est de $\pm U$ ».
\item \textcolor{red}{$\times$} « La valeur vraie EST dans l'intervalle ».
\item \textcolor{popgreen}{$\checkmark$} « Le résultat est $\bar{x} \pm U$ avec un niveau de confiance de 95\,\% ».
\item \textcolor{popgreen}{$\checkmark$} « L'intervalle de confiance à 95\,\% est $[\bar{x}-U, \bar{x}+U]$ ».
\end{itemize}
\end{attention}

\coursec{Expression finale d'un résultat de mesure}

\courssub{Synthèse : de la mesure brute au résultat final}

\begin{methode}[Procédure complète d'expression d'un résultat]
\begin{enumerate}
\item \textbf{Modéliser} : identifier le mesurande, l'équation de mesure (ex. $g = 4\pi^2 L/T^2$), les grandeurs d'entrée.
\item \textbf{Mesurer} : effectuer une série ($n\ge 10$) ou une mesure unique selon le contexte. Noter les conditions (T, P, opérateur).
\item \textbf{Analyser les sources} : lister Type A (dispersion série) et Type B (résolution, étalonnage, influences).
\item \textbf{Quantifier $u_i$} : convertir chaque source en incertitude-type (diviseurs $\sqrt{3}, \sqrt{6}, 2, \sqrt{n}$).
\item \textbf{Combiner} : $u_c = \sqrt{\sum u_i^2}$.
\item \textbf{Élargir} : $U = k \cdot u_c$ ($k=2$ pour 95\,\%).
\item \textbf{Arrondir} : $U$ à 1 ou 2 ch. sig. ; $\bar{x}$ à la même décimale.
\item \textbf{Rédiger} : $\boxed{\text{Mesurande} = \bar{x} \pm U\ \text{(unité SI)},\quad \text{niveau de confiance } 95\%}$.
\item \textbf{Conclure} : comparer à une valeur de référence, vérifier compatibilité, commenter la source dominante.
\end{enumerate}
\end{methode}

\begin{exoresolu}[Complet : Détermination de $g$ par pendule simple]
\textbf{Énoncé.} On mesure $L = 0,995\ \mathrm{m}$ (règle mm, $u_L = 0,001\ \mathrm{m}$) et 10 périodes $T_i$ donnant $\bar{T} = 2,005\ \mathrm{s}$, $s_T = 0,012\ \mathrm{s}$. Chronomètre résolution 0,01\,s. Calculer $g$ et son incertitude.
\textbf{Solution.}
\begin{enumerate}
\item Modèle : $g = \dfrac{4\pi^2 L}{\bar{T}^2}$.
\item Valeur nominale : $g_0 = \dfrac{4\pi^2 \times 0,995}{2,005^2} = 9,78\ \mathrm{m/s^2}$.
\item Incertitudes-types sur les entrées :
$u(L) = 0,001\ \mathrm{m}$ (Type B, lecture).
$u_A(\bar{T}) = s_T/\sqrt{10} = 0,012/\sqrt{10} = 0,0038\ \mathrm{s}$.
$u_{\text{résol}}(T) = 0,01/\sqrt{12} = 0,0029\ \mathrm{s}$.
$u(\bar{T}) = \sqrt{0,0038^2 + 0,0029^2} = 0,0048\ \mathrm{s}$.
\item Propagation (formule simplifiée puissances) : $\frac{u_c(g)}{g} = \sqrt{ \left(\frac{u(L)}{L}\right)^2 + \left(2\frac{u(\bar{T})}{\bar{T}}\right)^2 }$.
$\frac{u(L)}{L} = 0,001/0,995 = 0,0010$.
$2\frac{u(\bar{T})}{\bar{T}} = 2 \times 0,0048 / 2,005 = 0,0048$.
$u_c(g)/g = \sqrt{0,0010^2 + 0,0048^2} = 0,0049$.
$u_c(g) = 9,78 \times 0,0049 = 0,048\ \mathrm{m/s^2}$.
\item Élargie ($k=2$) : $U = 2 \times 0,048 = 0,096 \to 0,10\ \mathrm{m/s^2}$ (2 ch. sig. car commence par 1).
\item Arrondi $g_0$ : même décimale (centièmes) $\to 9,78\ \mathrm{m/s^2}$.
\item \textbf{Résultat final :} $\boxed{g = 9,78 \pm 0,10\ \mathrm{m/s^2}\ (95\%)}$.
\item Interprétation : Intervalle $[9,68 ; 9,88]$. La valeur de référence $9,81$ est dans l'intervalle $\to$ \textbf{compatibilité expérimentale validée}.
\end{enumerate}
\end{exoresolu}

\begin{exercice}[Exercices d'entraînement — Probatoire]
\begin{enumerate}
\item \textbf{Qualités métrologiques.} Un thermomètre à mercure (échelle 0,1\,°C) est utilisé pour mesurer la température de fusion de la glace (0,00\,°C). On fait 5 mesures : 0,2 ; 0,1 ; 0,3 ; 0,2 ; 0,1 °C.
\begin{itemize}
\item[a)] Calculer la moyenne et l'écart-type expérimental.
\item[b)] Ce thermomètre est-il fidèle ? Justifiez.
\item[c)] Est-il juste ? Justifiez.
\item[d)] Quelle est son incertitude-type de type A sur la moyenne ?
\end{itemize}
\item \textbf{Mesure unique (Type B).} On mesure une tension $U = 12,34\ \mathrm{V}$ sur un multimètre numérique (calibre 20\,V, résolution 0,01\,V). Le certificat d'étalonnage donne $U_{\text{étal}} = 0,05\ \mathrm{V}$ ($k=2$).
\begin{itemize}
\item[a)] Calculer $u_{\text{résol}}$ et $u_{\text{étal}}$.
\item[b)] Calculer $u_c$ et $U$ (95\,\%).
\item[c)] Exprimer le résultat final.
\end{itemize}
\item \textbf{Série de mesures (Type A + B).} Pour étalonner une balance, on pèse 10 fois une masse étalon de 200,00\,g. Résultats (en g) : 200,01 ; 199,99 ; 200,02 ; 200,00 ; 199,98 ; 200,01 ; 200,00 ; 199,99 ; 200,02 ; 200,00. Résolution balance : 0,01\,g.
\begin{itemize}
\item[a)] Calculer $\bar{m}$, $s$, $u_A$.
\item[b)] Estimer $u_{\text{résol}}$.
\item[c)] Calculer $u_c$, $U$ ($k=2$) et donner le résultat final.
\item[d)] La balance est-elle juste au sens métrologique ? (Seuil d'acceptation : $\pm 0,02\ \mathrm{g}$).
\end{itemize}
\end{enumerate}
\end{exercice}

\begin{corrige}[Exercices d'entraînement]
\textbf{Exercice 1.}
\begin{itemize}
\item[a)] $\bar{T} = 0,18\ \mathrm{°C}$ ; $s = \sqrt{\frac{(0,02)^2+(-0,08)^2+(0,12)^2+(0,02)^2+(-0,08)^2}{4}} = \sqrt{\frac{0,028}{4}} = 0,084\ \mathrm{°C}$.
\item[b)] Fidélité : $s = 0,08\ \mathrm{°C}$ (dispersion faible, de l'ordre de la résolution). \textbf{Oui, assez fidèle}.
\item[c)] Justesse : Biais $= \bar{T} - 0,00 = +0,18\ \mathrm{°C}$. Erreur systématique nette. \textbf{Non, pas juste}.
\item[d)] $u_A = s/\sqrt{5} = 0,084/\sqrt{5} = 0,037\ \mathrm{°C}$.
\end{itemize}
\textbf{Exercice 2.}
\begin{itemize}
\item[a)] $u_{\text{résol}} = 0,01/\sqrt{12} = 0,0029\ \mathrm{V}$. $u_{\text{étal}} = 0,05/2 = 0,025\ \mathrm{V}$.
\item[b)] $u_c = \sqrt{0,0029^2 + 0,025^2} = 0,025\ \mathrm{V}$ (étalonnage domine). $U = 2 \times 0,025 = 0,05\ \mathrm{V}$.
\item[c)] $U=0,05\ \mathrm{V}$ (1 ch. sig. $\to$ 0,05). $\bar{U}$ au centième : $12,34\ \mathrm{V}$. $\boxed{U = 12,34 \pm 0,05\ \mathrm{V}\ (95\%)}$.
\end{itemize}
\textbf{Exercice 3.}
\begin{itemize}
\item[a)] $\bar{m} = 200,002\ \mathrm{g} \approx 200,00\ \mathrm{g}$. $s = 0,014\ \mathrm{g}$. $u_A = 0,014/\sqrt{10} = 0,0044\ \mathrm{g}$.
\item[b)] $u_{\text{résol}} = 0,01/\sqrt{12} = 0,0029\ \mathrm{g}$.
\item[c)] $u_c = \sqrt{0,0044^2 + 0,0029^2} = 0,0053\ \mathrm{g}$. $U = 2 \times 0,0053 = 0,0106 \to 0,011\ \mathrm{g}$ (2 ch. sig. car 1). $\bar{m}$ au centième : $200,00\ \mathrm{g}$. $\boxed{m = 200,00 \pm 0,011\ \mathrm{g}\ (95\%)}$.
\item[d)] Biais $= 200,00 - 200,00 = 0,00\ \mathrm{g} < 0,02\ \mathrm{g}$. \textbf{Oui, la balance est juste} (dans les specs).
\end{itemize}
\end{corrige}

\begin{savaistu}[Le Guide pour l'expression de l'incertitude de mesure (GUM)]
Publié en 1993 par le BIPM (avec ISO, CEI, OIML, etc.), le GUM (NF ISO/CEI Guide 98-3) est la « bible » mondiale de l'évaluation des incertitudes. Il unifie les pratiques : fini les « erreurs maximales » arbitraires, place aux \textbf{incertitudes-types} combinées par la loi de propagation des variances. Au Cameroun, l'ANOR impose sa méthode pour l'accréditation des laboratoires (norme ISO/CEI 17025). Maîtriser ce vocabulaire (Type A, Type B, $u_c$, $U$, $k$) est indispensable pour tout ingénieur ou technicien supérieur.
\end{savaistu}

\coursec{Qualités d'un instrument de mesure}

Dans la section précédente, nous avons vu que toute mesure est entachée d'une erreur et que le résultat d'une mesure n'a de sens que s'il est accompagné d'une incertitude. Mais d'où viennent ces erreurs ? Une grande partie d'entre elles dépend de l'\cle{instrument de mesure} lui-même. Deux thermomètres posés côte à côte dans la même salle de classe à Yaoundé peuvent afficher \SI{25,4}{\celsius} et \SI{26,1}{\celsius} : lequel croire ? Pour répondre à cette question, les scientifiques ont défini des \cle{qualités métrologiques} qui permettent de caractériser et de comparer les instruments. C'est l'objet de cette section.

\courssub{Intuition : la métaphore de la cible de fléchettes}

Imagine un joueur de fléchettes qui lance plusieurs fléchettes sur une cible. Le centre de la cible représente la \cle{valeur vraie} de la grandeur à mesurer, et chaque fléchette plantée représente une mesure individuelle. Deux questions naturelles se posent :

\begin{itemize}
\item Les fléchettes sont-elles \textbf{groupées} les unes près des autres ? C'est la question de la \cle{fidélité}.
\item Les fléchettes sont-elles proches du \textbf{centre} de la cible ? C'est la question de la \cle{justesse}.
\end{itemize}

Ces deux qualités sont \textbf{indépendantes} : on peut être groupé mais loin du centre, ou dispersé autour du centre. Garde cette image en tête pendant toute la section : elle sera notre fil conducteur.

\courssub{La sensibilité}

\begin{definition}[Sensibilité]
La \cle{sensibilité} d'un instrument de mesure est son aptitude à détecter de \textbf{petites variations} de la grandeur mesurée. Plus un instrument est sensible, plus il réagit à de faibles changements de la grandeur.
\end{definition}

Concrètement, la sensibilité est liée à la \textbf{plus petite variation perceptible} sur l'instrument :

\begin{itemize}
\item pour un instrument \textbf{analogique} (à aiguille ou à graduation), elle est liée à la valeur de la \cle{plus petite graduation} : une règle graduée au millimètre est plus sensible qu'une règle graduée au centimètre ;
\item pour un instrument \textbf{numérique} (à affichage digital), elle est liée à la \cle{résolution}, c'est-à-dire la plus petite variation que l'afficheur peut montrer : un chronomètre affichant le centième de seconde est plus sensible qu'un chronomètre affichant la seconde.
\end{itemize}

\begin{exemplebox}[Comparaison de sensibilités]
Une balance de ménage affiche les masses au gramme près : sa résolution est \SI{1}{g}. La balance électronique du laboratoire de physique affiche le centigramme : sa résolution est \SI{0,01}{g}. Pour détecter la perte de masse d'une feuille qui sèche, seule la balance de laboratoire est assez \textbf{sensible}.
\end{exemplebox}

\begin{attention}[Sensibilité et qualité]
Un instrument très sensible n'est pas forcément de bonne qualité ! Un voltmètre qui affiche le millivolt mais dont l'aiguille tremble sans cesse, ou qui est mal étalonné, donnera des mesures fausses avec beaucoup de chiffres. La sensibilité ne dit rien sur la fidélité ni sur la justesse.
\end{attention}

\courssub{La fidélité}

\begin{definition}[Fidélité]
La \cle{fidélité} d'un instrument de mesure est son aptitude à donner des indications \textbf{très proches les unes des autres} lorsqu'on répète plusieurs fois la mesure d'une même grandeur, dans les \textbf{mêmes conditions}. Un instrument fidèle donne des mesures peu dispersées : l'\cle{erreur aléatoire} est faible.
\end{definition}

Sur la cible de fléchettes, un instrument fidèle correspond à des fléchettes \textbf{groupées}, que ce soit au centre ou non. La fidélité se quantifie par la \textbf{dispersion} des mesures répétées : plus l'écart-type de la série de mesures est petit (nous le calculerons dans la section 4), plus l'instrument est fidèle.

\begin{exemplebox}[Test de fidélité d'un chronomètre]
Un élève mesure dix fois de suite la durée de dix oscillations d'un pendule avec le même chronomètre. Il obtient (en secondes) : 20,12 ; 20,15 ; 20,11 ; 20,14 ; 20,13 ; 20,12 ; 20,15 ; 20,11 ; 20,14 ; 20,13. Les valeurs sont très groupées (elles s'écartent de moins de \SI{0,05}{s}) : le chronomètre est \textbf{fidèle}. Mais cela ne prouve pas qu'il est juste : s'il retarde systématiquement de \SI{0,3}{s}, toutes les mesures seront groupées... et fausses !
\end{exemplebox}

\courssub{La justesse}

\begin{definition}[Justesse]
La \cle{justesse} d'un instrument de mesure est son aptitude à donner des indications \textbf{proches de la valeur vraie} (ou d'une valeur de référence acceptée) de la grandeur mesurée. Un instrument juste est peu affecté par les \cle{erreurs systématiques} (défaut d'étalonnage, zéro décalé, parallaxe systématique, etc.).
\end{definition}

Sur la cible, un instrument juste correspond à des fléchettes centrées \textbf{autour du centre}, même si elles sont dispersées. La justesse se quantifie par l'\textbf{écart entre la moyenne des mesures et la valeur vraie} : cet écart est l'erreur systématique (ou biais).

\begin{exemplebox}[Balance mal réglée]
Une balance de marché à Douala indique \SI{1,10}{kg} pour un sac d'étalonnage de \SI{1,000}{kg}, et ce de façon répétable. Elle est fidèle (indications reproductibles) mais \textbf{pas juste} : elle commet une erreur systématique de $+\SI{0,10}{kg}$. Le client est lésé à chaque pesée !
\end{exemplebox}

\courssub{La précision : la qualité globale}

\begin{definition}[Précision]
La \cle{précision} d'un instrument de mesure est sa qualité globale, qui combine la \textbf{fidélité} et la \textbf{justesse}. Un instrument précis donne des mesures à la fois \textbf{groupées} (fidèle) et \textbf{proches de la valeur vraie} (juste).
\end{definition}

\begin{aretenir}[Les quatre situations possibles]
\begin{center}
\begin{tabularx}{\linewidth}{|l|X|X|}
\hline
\rowcolor{popblueL} & \textbf{Juste} & \textbf{Non juste} \\
\hline
\textbf{Fidèle} & Mesures groupées autour de la valeur vraie : instrument \textbf{précis}. & Mesures groupées mais décalées : erreur systématique (mauvais étalonnage). \\
\hline
\textbf{Non fidèle} & Mesures dispersées mais centrées sur la valeur vraie : la moyenne est exploitable. & Mesures dispersées et décalées : instrument inutilisable. \\
\hline
\end{tabularx}
\end{center}
\end{aretenir}

\begin{popfigure}
\begin{center}
\begin{tikzpicture}[scale=0.62]
% Cible 1 : fidèle et juste
\begin{scope}[shift={(0,0)}]
\draw[fill=popblue!10] (0,0) circle (1.5);
\draw[fill=white] (0,0) circle (1.0);
\draw[fill=popblue!25] (0,0) circle (0.5);
\fill[popink] (0,0) circle (0.05);
\foreach \p in {(0.15,0.2),(-0.2,0.1),(0.1,-0.18),(-0.12,-0.15),(0.22,-0.05)}
  \fill[poporange] \p circle (0.07);
\node[below,font=\small\bfseries] at (0,-1.7) {Fidèle et juste};
\node[below,font=\small] at (0,-2.15) {(précis)};
\end{scope}
% Cible 2 : fidèle, non juste
\begin{scope}[shift={(4.6,0)}]
\draw[fill=popblue!10] (0,0) circle (1.5);
\draw[fill=white] (0,0) circle (1.0);
\draw[fill=popblue!25] (0,0) circle (0.5);
\fill[popink] (0,0) circle (0.05);
\foreach \p in {(1.0,0.85),(0.85,1.05),(1.1,1.0),(0.95,0.7),(1.15,0.8)}
  \fill[poporange] \p circle (0.07);
\node[below,font=\small\bfseries] at (0,-1.7) {Fidèle, non juste};
\node[below,font=\small] at (0,-2.15) {(erreur systématique)};
\end{scope}
% Cible 3 : non fidèle, juste
\begin{scope}[shift={(9.2,0)}]
\draw[fill=popblue!10] (0,0) circle (1.5);
\draw[fill=white] (0,0) circle (1.0);
\draw[fill=popblue!25] (0,0) circle (0.5);
\fill[popink] (0,0) circle (0.05);
\foreach \p in {(0.9,0.4),(-0.8,0.6),(0.3,-0.9),(-0.5,-0.7),(0.1,1.1)}
  \fill[poporange] \p circle (0.07);
\node[below,font=\small\bfseries] at (0,-1.7) {Non fidèle, juste};
\node[below,font=\small] at (0,-2.15) {(dispersion élevée)};
\end{scope}
% Cible 4 : non fidèle, non juste
\begin{scope}[shift={(13.8,0)}]
\draw[fill=popblue!10] (0,0) circle (1.5);
\draw[fill=white] (0,0) circle (1.0);
\draw[fill=popblue!25] (0,0) circle (0.5);
\fill[popink] (0,0) circle (0.05);
\foreach \p in {(1.2,1.0),(0.6,1.3),(1.35,0.5),(0.9,0.6),(1.25,1.3)}
  \fill[poporange] \p circle (0.07);
\node[below,font=\small\bfseries] at (0,-1.7) {Non fidèle, non juste};
\node[below,font=\small] at (0,-2.15) {(mauvais instrument)};
\end{scope}
\end{tikzpicture}
\end{center}
\legende{La métaphore de la cible : chaque point orange représente une mesure, le centre noir la valeur vraie. Seule la première cible correspond à un instrument précis (fidèle \emph{et} juste).}
\end{popfigure}

\begin{attention}[Erreur fréquente : fidèle ne veut pas dire juste]
Au Probatoire, beaucoup d'élèves affirment qu'un instrument qui donne « toujours la même valeur » est forcément bon. C'est \textbf{faux} : un instrument fidèle n'est pas forcément juste. Des mesures groupées mais décalées de la valeur vraie traduisent une \textbf{erreur systématique} (zéro mal réglé, étalonnage faux). Inversement, un instrument juste mais peu fidèle donne des mesures dispersées dont seule la \textbf{moyenne} approche la valeur vraie. Les deux défauts ont des causes et des remèdes différents : on corrige une erreur systématique par un réglage ou un étalonnage ; on réduit la dispersion en répétant les mesures et en moyennant.
\end{attention}

\courssub{Graduation, résolution et classe de précision}

\courssub{Lecture sur un instrument analogique}

Un instrument \textbf{analogique} (voltmètre à aiguille, thermomètre à mercure, règle graduée) possède une échelle graduée. Trois grandeurs le caractérisent :

\begin{itemize}
\item le \cle{calibre} $C$ : la valeur maximale mesurable sur le réglage choisi ;
\item l'\cle{échelle} $E$ : le nombre total de divisions du cadran ;
\item la valeur d'une \cle{graduation} : $g = \dfrac{C}{E}$.
\end{itemize}

\begin{exemplebox}[Graduation d'un voltmètre]
Un voltmètre analogique possède une échelle de $E = 100$ divisions. Sur le calibre $C = \SI{30}{V}$, une graduation vaut
\[ g = \frac{C}{E} = \frac{30}{100} = \SI{0,3}{V}. \]
Sur le calibre $C = \SI{3}{V}$, une graduation vaut \SI{0,03}{V} : la lecture y est dix fois plus fine. C'est pourquoi on choisit toujours le \textbf{plus petit calibre} qui contient la valeur à mesurer.
\end{exemplebox}

\courssub{Résolution d'un instrument numérique}

\begin{definition}[Résolution]
La \cle{résolution} d'un instrument numérique est la plus petite variation de la grandeur mesurée que l'afficheur peut traduire : elle correspond à la valeur d'une unité du dernier chiffre affiché.
\end{definition}

Un multimètre numérique qui affiche $U = \SI{12,36}{V}$ a une résolution de \SI{0,01}{V} sur ce calibre. Un affichage avec beaucoup de chiffres ne garantit pas la justesse : la notice de l'appareil indique toujours une incertitude constructeur (souvent exprimée en pourcentage de la lecture, plus un nombre de « digits »).

\courssub{La classe de précision d'un appareil analogique}

Les appareils analogiques de mesure électriques (voltmètres, ampèremètres) portent sur leur cadran un nombre encadré ou inscrit, appelé \cle{classe de précision} (ou classe de l'appareil), noté par exemple 1,5 ou 2.

\begin{definition}[Classe de précision]
La \cle{classe de précision} d'un appareil analogique est un nombre sans unité, fourni par le constructeur, qui caractérise l'erreur maximale que peut commettre l'appareil, exprimée en pourcentage du \textbf{calibre} utilisé.
\end{definition}

\begin{propriete}[Incertitude liée à la classe]
Pour un appareil analogique de classe $k$ utilisé sur le calibre $C$, l'incertitude absolue maximale sur la mesure est
\[ \Delta = \frac{k \times C}{100}. \]
$\Delta$ s'exprime dans la même unité que le calibre $C$. Cette incertitude est \textbf{la même pour toutes les lectures} effectuées sur ce calibre, que l'aiguille soit au début ou à la fin de l'échelle.
\end{propriete}

\begin{methode}[Calculer l'incertitude liée à la classe d'un appareil analogique]
\begin{enumerate}
\item Relever la \textbf{classe} $k$ inscrite sur le cadran (symbole du type \og 1,5 \fg{}).
\item Identifier le \textbf{calibre} $C$ réellement utilisé (borne choisie), avec son unité.
\item Calculer $\Delta = \dfrac{k \times C}{100}$ et l'exprimer dans l'unité du calibre.
\item Arrondir $\Delta$ par excès à un chiffre significatif, puis écrire le résultat sous la forme $X = (x \pm \Delta)$ en alignant le dernier chiffre de $x$ sur celui de $\Delta$.
\end{enumerate}
\end{methode}

\begin{exemplebox}[Voltmètre de classe 2 sur le calibre 30 V]
Un voltmètre de classe $k = 2$ est utilisé sur le calibre $C = \SI{30}{V}$. L'incertitude absolue vaut
\[ \Delta U = \frac{k \times C}{100} = \frac{2 \times 30}{100} = \SI{0,6}{V}. \]
Si l'aiguille indique $U = \SI{12,0}{V}$, on écrira le résultat : $U = (12{,}0 \pm 0{,}6)\ \si{V}$.
\end{exemplebox}

\begin{exoresolu}[Voltmètre analogique : lecture et incertitude]
\textbf{Énoncé.} Un voltmètre analogique de classe 1,5 possède une échelle de 150 divisions. On l'utilise sur le calibre \SI{15}{V} et l'aiguille s'arrête sur la division 108.
\begin{enumerate}
\item Calculer la valeur d'une graduation, puis la tension lue $U$.
\item Calculer l'incertitude absolue $\Delta U$ due à la classe de l'appareil.
\item Exprimer le résultat de la mesure.
\item Calculer l'incertitude relative $\dfrac{\Delta U}{U}$ et commenter.
\end{enumerate}
\tcblower
\textbf{Solution.}
\begin{enumerate}
\item La valeur d'une graduation est
\[ g = \frac{C}{E} = \frac{15}{150} = \SI{0,1}{V}. \]
La tension lue vaut donc
\[ U = 108 \times 0{,}1 = \SI{10,8}{V}. \]
\item Incertitude liée à la classe :
\[ \Delta U = \frac{k \times C}{100} = \frac{1{,}5 \times 15}{100} = \SI{0,225}{V} \approx \SI{0,3}{V} \]
(arrondi par excès à un chiffre significatif).
\item Résultat de la mesure :
\[ U = (10{,}8 \pm 0{,}3)\ \si{V}. \]
\item Incertitude relative :
\[ \frac{\Delta U}{U} = \frac{0{,}3}{10{,}8} \approx 0{,}028 \approx 3\ \%. \]
La mesure est de bonne qualité. Remarque : si l'aiguille n'avait indiqué que \SI{1,5}{V} (division 15), l'incertitude absolue serait restée \SI{0,3}{V} et l'incertitude relative aurait atteint 20 \% ! C'est pourquoi on choisit le calibre de sorte que l'aiguille dévie dans le \textbf{tiers supérieur} du cadran.
\end{enumerate}
\end{exoresolu}

\begin{popfigure}
\begin{center}
\begin{tikzpicture}[scale=1.0]
% Cadran semi-circulaire
\draw[thick,fill=popcream] (0,0) -- (210:3.2) arc (210:-30:3.2) -- cycle;
\draw[thick] (0,0) circle (0.15);
% Graduations : 15 grandes divisions de 210° à -30° (240°)
\foreach \i in {0,...,15}{
  \pgfmathsetmacro{\ang}{210 - \i*16}
  \draw[thick] (\ang:2.9) -- (\ang:3.2);
}
% Petites graduations
\foreach \i in {0,...,75}{
  \pgfmathsetmacro{\ang}{210 - \i*3.2}
  \draw[thin] (\ang:3.0) -- (\ang:3.2);
}
% Quelques chiffres
\node[font=\small] at (210:2.45) {0};
\node[font=\small] at (130:2.45) {5};
\node[font=\small] at (50:2.45) {10};
\node[font=\small] at (-30:2.45) {15};
% Aiguille sur la division 108/150 -> angle = 210 - 108*1.6 = 37.2°
\draw[very thick,poporange,->] (0,0) -- (37.2:2.85);
\fill[popink] (0,0) circle (0.08);
% Inscriptions cadran
\node[font=\small\bfseries] at (0,-0.9) {V};
\node[font=\footnotesize] at (0,-1.35) {classe 1,5};
\node[font=\footnotesize] at (0,-1.75) {calibre : \SI{15}{V}};
\node[font=\footnotesize] at (0,-2.15) {échelle : 150 divisions};
% Annotation aiguille
\node[font=\footnotesize,poporange] at (2.6,1.6) {division 108};
\end{tikzpicture}
\end{center}
\legende{Cadran d'un voltmètre analogique : calibre \SI{15}{V}, échelle de 150 divisions (une graduation vaut \SI{0,1}{V}), classe de précision 1,5. L'aiguille indique la division 108, soit $U = \SI{10,8}{V}$.}
\end{popfigure}

\courssub{Qualités des instruments courants du laboratoire}

\begin{center}
\begin{tabularx}{\linewidth}{|l|X|X|X|}
\hline
\rowcolor{popblueL} \textbf{Instrument} & \textbf{Grandeur} & \textbf{Sensibilité typique} & \textbf{Source d'erreur dominante} \\
\hline
Règle graduée & Longueur & \SI{1}{mm} (graduation) & Alignement, parallaxe \\
\hline
Balance électronique & Masse & \SI{0,01}{g} (résolution) & Étalonnage, position du plateau \\
\hline
Voltmètre analogique & Tension & $g = C/E$ & Classe de précision, parallaxe \\
\hline
Thermomètre & Température & \SI{0,5}{\celsius} à \SI{0,1}{\celsius} & Temps de réponse, immersion \\
\hline
Chronomètre & Durée & \SI{0,01}{s} (résolution) & Temps de réaction de l'opérateur (\SI{0,2}{s}) \\
\hline
\end{tabularx}
\end{center}

Remarque essentielle pour le chronomètre manuel : sa \textbf{résolution} est excellente (\SI{0,01}{s}), mais l'erreur dominante vient du \textbf{temps de réaction} de l'expérimentateur, de l'ordre de \SI{0,2}{s}. L'instrument est sensible et fidèle ; c'est l'opérateur qui limite la justesse ! C'est pourquoi, pour mesurer la période d'un pendule, on chronomètre dix ou vingt oscillations et on divise : l'erreur de déclenchement est ainsi répartie sur de nombreuses périodes.

\begin{savaistu}[Une métaphore universelle]
La métaphore de la cible de fléchettes (ou de la cible de tir à l'arc) est utilisée dans tous les laboratoires du monde, du lycée de Ngaoundéré au Bureau international des poids et mesures (BIPM) à Sèvres, pour enseigner la différence entre fidélité et justesse. Au BIPM était conservé pendant plus d'un siècle le prototype international du kilogramme, un cylindre de platine-iridium qui servait de \og centre de la cible \fg{} pour toutes les balances de la planète. Depuis 2019, le kilogramme est défini à partir de la constante de Planck : la \og valeur vraie \fg{} n'est plus un objet, mais une loi de la physique !
\end{savaistu}

\begin{aretenir}[L'essentiel de la section]
\begin{itemize}
\item \textbf{Sensibilité} : aptitude à détecter de petites variations (liée à la plus petite graduation ou à la résolution).
\item \textbf{Fidélité} : mesures répétées proches les unes des autres (faible dispersion, faible erreur aléatoire).
\item \textbf{Justesse} : mesures proches de la valeur vraie (faible erreur systématique).
\item \textbf{Précision} = fidélité $+$ justesse.
\item Appareil analogique de classe $k$ sur le calibre $C$ : $\Delta = \dfrac{k \times C}{100}$, identique pour toute lecture sur ce calibre.
\item Un instrument fidèle n'est pas forcément juste : pense toujours à la cible !
\end{itemize}
\end{aretenir}

\coursec{Incertitude sur une mesure unique}

Dans la section précédente, nous avons appris à juger un instrument de mesure : sa \cle{fidélité} (donne-t-il des indications reproductibles ?), sa \cle{justesse} (donne-t-il des indications proches de la valeur vraie ?), sa \cle{sensibilité} et sa \cle{précision}. Ces qualités décrivent l'outil. Il reste maintenant à répondre à la question concrète de l'expérimentateur : \emph{je viens de faire une mesure, quelle incertitude dois-je attacher à mon résultat ?}

Deux situations se présentent au laboratoire. Parfois, on dispose du temps et des moyens de répéter la mesure de nombreuses fois : on pourra alors traiter statistiquement la série de valeurs obtenues (ce sera l'objet des sections 4 et 5). Mais très souvent — et c'est le cas de la quasi-totalité des mesures que tu feras en TP de Première — on n'effectue qu'\textbf{une seule mesure} : on mesure une fois la longueur du pendule, une fois la tension aux bornes du conducteur ohmique, une fois la masse du solide. C'est cette situation que nous étudions ici.

\courssub{Pourquoi la mesure unique impose une démarche différente}

\begin{experience}[Une mesure qui ne se répète pas]
Au lycée de Ngaoundéré, un groupe d'élèves mesure la durée de chute d'une bille lâchée du haut d'un escalier, puis la longueur d'une table avec une règle graduée au millimètre. Pour la durée de chute, ils peuvent recommencer vingt fois et comparer les valeurs : la dispersion des résultats renseigne sur l'incertitude. Mais pour la longueur de la table, lue une seule fois : $L = 1{,}243\ \text{m}$, aucune statistique n'est possible. Pourtant, l'élève sait pertinemment que son résultat n'est pas exact : l'alignement de l'œil, la position du zéro de la règle, la finesse des graduations limitent la qualité de la lecture. Comment chiffrer cette limitation sans répéter la mesure ?
\end{experience}

Lorsqu'une mesure n'est effectuée qu'une fois, il est \textbf{impossible d'étudier la dispersion} des résultats : un échantillon d'une seule valeur n'a ni moyenne significative ni écart-type. On adopte alors une démarche \emph{a priori} : avant même de discuter le résultat, on estime l'incertitude à partir de ce que l'on sait de l'\textbf{instrument} et des \textbf{conditions de la mesure}. Dans le vocabulaire de la métrologie, on parle d'\textbf{évaluation de type B} de l'incertitude, par opposition à l'évaluation de type A (statistique) réservée aux séries de mesures.

\begin{definition}[Incertitude de lecture (évaluation de type B)]
L'\cle{incertitude de lecture}, notée $\Delta x$, est l'estimation, faite \emph{a priori} à partir des caractéristiques de l'instrument et des conditions d'observation, de l'écart maximal plausible entre la valeur lue $x_{\text{lu}}$ et la valeur que l'on aurait obtenue avec un instrument parfait. Elle est dite de \textbf{type B} car elle ne provient pas d'une analyse statistique de mesures répétées, mais d'un jugement raisonné sur l'instrument : finesse des graduations, résolution de l'affichage, classe de précision, notice du constructeur.
\end{definition}

\begin{aretenir}[Le principe directeur]
Pour une mesure unique, l'incertitude ne se calcule pas sur les résultats : elle se \textbf{décrète à partir de l'instrument}. La question à se poser est toujours la même : « avec cet instrument, utilisé dans ces conditions, de combien ma lecture peut-elle raisonnablement s'écarter de la réalité ? »
\end{aretenir}

\courssub{Cas d'un instrument analogique : la règle de la demi-graduation}

Un instrument \textbf{analogique} est un instrument dont l'indication se lit sur une échelle graduée : règle, burette, thermomètre à colonne de liquide, voltmètre à aiguille. La valeur mesurée tombe en général \emph{entre deux graduations}, et l'opérateur doit estimer la position à l'œil.

Posons $a$ la valeur d'une graduation (par exemple $a = 1\ \text{mm}$ pour une règle millimétrée). Lorsque la lecture est faite avec soin, l'œil entraîné est capable de repérer si l'extrémité de l'objet est plus proche d'une graduation ou de l'autre, voire de repérer le milieu de l'intervalle. En revanche, il est illusoire de prétendre distinguer le dixième de graduation à l'œil nu. L'expérience collective des métrologues montre que l'erreur de lecture commise par un opérateur soigneux ne dépasse pas, en général, la moitié d'une graduation.

\begin{propriete}[Règle de la demi-graduation]
Pour une lecture soigneuse sur un instrument analogique dont la plus petite graduation vaut $a$, l'incertitude de lecture est estimée à une demi-graduation :
\[ \Delta x = \frac{a}{2}. \]
Dans des conditions de lecture dégradées (mauvais éclairage, échelle éloignée, aiguille épaisse, parallaxe difficile à éviter, mesure faite à la hâte), on majore prudemment l'incertitude en prenant une graduation entière : $\Delta x = a$.
\end{propriete}

\begin{popfigure}
\begin{center}
\begin{tikzpicture}[scale=1.0]
% Règle graduée
\draw[fill=popgoldL,draw=popdark] (0,0) rectangle (11,1);
\foreach \x in {0,1,...,10} {
  \draw[popdark,thick] (\x,1) -- (\x,0.55);
  \node[above,popdark] at (\x,1.05) {\small \x};
}
\foreach \x in {0.5,1.5,...,9.5} \draw[popdark] (\x,1) -- (\x,0.7);
\foreach \x in {0.1,0.2,0.3,0.4,0.6,0.7,0.8,0.9} {
  \foreach \k in {0,1,...,9} \draw[popdark!60] (\k+\x,1) -- (\k+\x,0.82);
}
% Objet mesuré (crayon)
\draw[fill=poporange,draw=popdark,thick] (0.6,1.25) rectangle (7.83,1.7);
\draw[fill=popdark] (7.83,1.25) -- (8.13,1.475) -- (7.83,1.7) -- cycle;
\node[popdark] at (3.5,1.475) {\small objet mesuré};
% Flèche de lecture
\draw[->,popink,very thick] (7.83,2.6) -- (7.83,1.85);
\node[popink,above] at (7.83,2.6) {\small l'extrémité tombe entre 7,8 cm et 7,9 cm};
% Zoom demi-graduation
\draw[popblue,thick] (7.3,-0.15) rectangle (8.4,-1.6);
\draw[popblue,dashed] (7.8,0) -- (7.8,-0.15);
\draw[popblue,dashed] (7.9,0) -- (8.35,-0.15);
\draw[popdark,thick] (7.45,-0.5) -- (7.45,-1.1);
\node[popdark,below] at (7.45,-1.1) {\footnotesize 7,8};
\draw[popdark,thick] (8.25,-0.5) -- (8.25,-1.1);
\node[popdark,below] at (8.25,-1.1) {\footnotesize 7,9};
\draw[popdark!60] (7.85,-0.5) -- (7.85,-0.85);
\draw[<->,popgreen,very thick] (7.45,-1.35) -- (7.85,-1.35);
\node[popgreen,below] at (7.65,-1.35) {\footnotesize $\Delta x = a/2$};
\draw[popink,very thick] (7.83,-0.5) -- (7.83,-0.62);
\end{tikzpicture}
\end{center}
\legende{Lecture sur une règle graduée au millimètre ($a = 0{,}1\ \text{cm}$). L'extrémité de l'objet tombe entre deux graduations ; on lit 7,8 cm et l'incertitude de lecture est $\Delta x = a/2 = 0{,}05\ \text{cm}$ (zoom).}
\end{popfigure}

\begin{attention}[Le piège classique : une graduation entière au lieu d'une demi-graduation]
L'erreur la plus fréquente au Probatoire est d'écrire systématiquement $\Delta x = a$ (une graduation entière) alors que la lecture a été faite soigneusement. La règle de référence est la \textbf{demi-graduation} : $\Delta x = a/2$. Ne prends une graduation entière que si l'énoncé précise des conditions de lecture difficiles. Piège symétrique : annoncer $\Delta x = a/10$ parce que « on peut estimer entre les graduations » — c'est surestimer l'œil humain, et c'est faux.
\end{attention}

\begin{exemplebox}[Longueur mesurée à la règle millimétrée]
On mesure la longueur d'un cahier avec une règle graduée au millimètre : $a = 1\ \text{mm} = 0{,}1\ \text{cm}$. La lecture soigneuse donne $L_{\text{lu}} = 24{,}3\ \text{cm}$. L'incertitude de lecture est :
\[ \Delta L = \frac{a}{2} = 0{,}5\ \text{mm} = 0{,}05\ \text{cm}. \]
On écrira donc le résultat : $L = (24{,}30 \pm 0{,}05)\ \text{cm}$. Remarque l'écriture $24{,}30$ : puisque l'incertitude porte sur le centième de centimètre, la valeur lue doit être écrite avec le même nombre de décimales.
\end{exemplebox}

\courssub{Cas d'un instrument numérique : la résolution de l'affichage}

Un instrument \textbf{numérique} affiche directement une valeur en chiffres : multimètre digital, balance électronique, chronomètre à affichage. Ici, plus de graduation à interpoler, mais une limitation nouvelle : l'appareil \textbf{arrondit} la grandeur à la dernière décimale qu'il peut afficher. Si un voltmètre affiche $4{,}37\ \text{V}$, la tension réelle peut valoir aussi bien $4{,}365\ \text{V}$ que $4{,}374\ \text{V}$ : l'affichage ne change pas.

\begin{definition}[Résolution d'un instrument numérique]
La \cle{résolution} d'un instrument numérique est la valeur d'une unité du dernier chiffre affiché. Si l'afficheur montre $4{,}37\ \text{V}$, la résolution vaut $0{,}01\ \text{V}$.
\end{definition}

\begin{propriete}[Incertitude liée à la résolution]
Pour une mesure unique réalisée avec un instrument numérique de résolution $r$, en l'absence d'indication du constructeur, on estime l'incertitude de lecture à une unité du dernier chiffre affiché :
\[ \Delta x = r. \]
Si la notice du constructeur fournit une incertitude $\Delta_{\text{cons}}$ (par exemple « $\pm (0{,}5\,\%\ \text{de la lecture} + 2\ \text{digits})$ »), c'est cette donnée, plus complète, qu'il faut utiliser.
\end{propriete}

\begin{popfigure}
\begin{center}
\begin{tikzpicture}
% Boîtier multimètre
\draw[fill=popblueL,draw=popdark,thick,rounded corners=6pt] (0,0) rectangle (6.4,3.4);
% Écran
\draw[fill=popdark,draw=popdark,rounded corners=3pt] (0.6,2.1) rectangle (5.8,3.0);
\node[popgreen,scale=1.9] at (3.2,2.55) {\ttfamily 4,37 V};
% Boutons
\draw[fill=popcream,draw=popdark] (1.4,1.0) circle (0.45);
\draw[fill=popcream,draw=popdark] (3.2,1.0) circle (0.45);
\draw[fill=popcream,draw=popdark] (5.0,1.0) circle (0.45);
\node[popdark] at (1.4,0.25) {\small V};
\node[popdark] at (3.2,0.25) {\small A};
\node[popdark] at (5.0,0.25) {\small $\Omega$};
% Annotation résolution
\draw[->,popink,very thick] (7.6,2.4) -- (5.0,2.55);
\node[popink,right,align=left] at (7.7,2.4) {\small dernier chiffre affiché :\\ \small le centième de volt};
\draw[<->,popgreen,very thick] (7.7,1.4) -- (8.7,1.4);
\node[popgreen,right,align=left] at (7.7,1.0) {\small résolution : $r = 0{,}01\ \text{V}$\\ \small donc $\Delta U = 0{,}01\ \text{V}$};
\end{tikzpicture}
\end{center}
\legende{Affichage numérique d'un voltmètre : $U = 4{,}37\ \text{V}$. La résolution est $r = 0{,}01\ \text{V}$ ; en l'absence de donnée constructeur, l'incertitude de lecture est $\Delta U = 0{,}01\ \text{V}$ et le résultat s'écrit $U = (4{,}37 \pm 0{,}01)\ \text{V}$.}
\end{popfigure}

\courssub{Cas d'un appareil de classe : la formule du constructeur}

Les appareils de mesure analogiques de qualité contrôlée (voltmètres et ampèremètres à aiguille de laboratoire, par exemple) portent une indication normalisée appelée \textbf{classe de précision}, souvent inscrite sur le cadran (valeurs usuelles : 0,5 ; 1 ; 1,5 ; 2,5). La classe est un nombre sans dimension qui exprime l'incertitude maximale de l'appareil en pourcentage du \textbf{calibre} utilisé.

\begin{definition}[Classe et calibre]
Le \cle{calibre} $C$ est la valeur maximale que l'appareil peut mesurer sur le réglage choisi. La \cle{classe} $k$ est le pourcentage du calibre qui borne l'erreur maximale de l'instrument, garantie par le constructeur.
\end{definition}

\begin{propriete}[Incertitude d'un appareil de classe]
Pour un appareil de classe $k$ utilisé sur le calibre $C$, l'incertitude instrumentale vaut :
\[ \Delta x = \frac{k \times C}{100}. \]
Vérification dimensionnelle : $k$ est sans dimension, $C$ s'exprime dans l'unité de la grandeur mesurée, donc $\Delta x$ s'exprime bien dans cette même unité.
\end{propriete}

\begin{attention}[Deux pièges avec les appareils de classe]
\begin{itemize}
\item L'incertitude dépend du \textbf{calibre}, pas de la valeur lue : avec un voltmètre de classe 1,5 sur le calibre $10\ \text{V}$, $\Delta U = 0{,}15\ \text{V}$, que l'aiguille indique $2\ \text{V}$ ou $9\ \text{V}$. C'est pourquoi on choisit toujours le \textbf{plus petit calibre possible} (celui qui donne la plus grande déviation de l'aiguille) : il minimise l'incertitude.
\item N'oublie pas de diviser par 100 : la classe est un \emph{pourcentage}. Écrire $\Delta x = k \times C$ est une erreur grossière qui gonfle l'incertitude d'un facteur 100.
\end{itemize}
\end{attention}

\begin{exemplebox}[Voltmètre à aiguille de laboratoire]
Un voltmètre à aiguille de classe $k = 1{,}5$ est utilisé sur le calibre $C = 30\ \text{V}$. L'incertitude instrumentale est :
\[ \Delta U = \frac{1{,}5 \times 30}{100} = 0{,}45\ \text{V}. \]
Si l'aiguille indique $U_{\text{lu}} = 12{,}0\ \text{V}$, le résultat s'écrit $U = (12{,}0 \pm 0{,}5)\ \text{V}$ en arrondissant l'incertitude à un chiffre significatif par excès.
\end{exemplebox}

\courssub{Combinaison des sources d'incertitude}

Dans la réalité d'une mesure unique, plusieurs sources d'incertitude coexistent : l'incertitude de lecture $\Delta_{\text{lec}}$ (demi-graduation ou résolution) et l'incertitude instrumentale $\Delta_{\text{cons}}$ donnée par le constructeur ou la classe. Comment les combiner ?

\begin{propriete}[Combinaison simple des incertitudes]
Soit une mesure unique entachée d'une incertitude de lecture $\Delta_{\text{lec}}$ et d'une incertitude constructeur $\Delta_{\text{cons}}$.
\begin{itemize}
\item \textbf{Approche prudente (par somme)} : $\Delta x = \Delta_{\text{lec}} + \Delta_{\text{cons}}$. Elle majore l'incertitude totale et convient lorsqu'on veut un résultat garanti.
\item \textbf{Approche simplifiée (par maximum)} : $\Delta x = \max(\Delta_{\text{lec}}\,;\,\Delta_{\text{cons}})$. Elle est admise lorsque l'une des deux sources domine nettement l'autre (rapport supérieur à 3 environ).
\end{itemize}
Au niveau Première, sauf indication contraire de l'énoncé, on retient l'approche par somme, qui est la plus sûre.
\end{propriete}

\courssub{Expression du résultat d'une mesure unique}

\begin{aretenir}[Écriture normalisée du résultat]
Le résultat d'une mesure unique s'exprime sous la forme :
\[ x = x_{\text{lu}} \pm \Delta x \quad (\text{unité}) \]
ce qui signifie que la valeur vraie se situe vraisemblablement dans l'intervalle $[\,x_{\text{lu}} - \Delta x\ ;\ x_{\text{lu}} + \Delta x\,]$. Règles de rédaction :
\begin{enumerate}
\item l'incertitude est arrondie à \textbf{un seul chiffre significatif} (deux au maximum), par excès ;
\item la valeur lue est arrondie à la \textbf{même décimale} que l'incertitude ;
\item l'unité est écrite une seule fois, commune aux deux nombres.
\end{enumerate}
\end{aretenir}

\begin{exoresolu}[Mesure de la longueur d'une table]
Énoncé. Un élève mesure la longueur d'une table de TP avec une règle graduée au millimètre, dans de bonnes conditions d'éclairage. Il lit $L_{\text{lu}} = 24{,}3\ \text{cm}$. Exprimer le résultat de la mesure sous forme normalisée et donner l'intervalle dans lequel se situe la valeur vraie.
\tcblower
Solution détaillée. La règle est un instrument analogique de plus petite graduation $a = 1\ \text{mm} = 0{,}1\ \text{cm}$. La lecture étant soigneuse, on applique la règle de la demi-graduation :
\[ \Delta L = \frac{a}{2} = \frac{0{,}1}{2} = 0{,}05\ \text{cm}. \]
L'incertitude porte sur le centième de centimètre ; on écrit donc la valeur lue avec deux décimales :
\[ \boxed{L = (24{,}30 \pm 0{,}05)\ \text{cm}}. \]
La valeur vraie de la longueur se situe dans l'intervalle :
\[ 24{,}30 - 0{,}05 \leq L \leq 24{,}30 + 0{,}05 \quad \text{soit} \quad 24{,}25\ \text{cm} \leq L \leq 24{,}35\ \text{cm}. \]
Remarque : écrire $L = (24{,}3 \pm 0{,}05)\ \text{cm}$ serait maladroit, car les deux nombres n'ont pas la même décimale ; écrire $\Delta L = 0{,}1\ \text{cm}$ serait l'erreur classique de la graduation entière.
\end{exoresolu}

\begin{exoresolu}[Tension mesurée avec un voltmètre de classe]
Énoncé. Pour mesurer la tension aux bornes d'un conducteur ohmique, on utilise un voltmètre à aiguille de classe $k = 2{,}5$, sur le calibre $C = 15\ \text{V}$. L'aiguille se stabilise sur $U_{\text{lu}} = 9{,}0\ \text{V}$. L'élève estime par ailleurs son incertitude de lecture sur le cadran à $\Delta_{\text{lec}} = 0{,}2\ \text{V}$. Exprimer le résultat de la mesure.
\tcblower
Solution détaillée. Incertitude instrumentale (appareil de classe) :
\[ \Delta_{\text{cons}} = \frac{k \times C}{100} = \frac{2{,}5 \times 15}{100} = 0{,}375\ \text{V} \approx 0{,}4\ \text{V}. \]
On combine avec l'incertitude de lecture par somme (approche prudente) :
\[ \Delta U = \Delta_{\text{lec}} + \Delta_{\text{cons}} = 0{,}2 + 0{,}375 = 0{,}575\ \text{V}. \]
On arrondit l'incertitude à un chiffre significatif, par excès : $\Delta U = 0{,}6\ \text{V}$. Le résultat s'écrit alors :
\[ \boxed{U = (9{,}0 \pm 0{,}6)\ \text{V}}. \]
La tension vraie se situe donc entre $8{,}4\ \text{V}$ et $9{,}6\ \text{V}$. Notons qu'avec l'approche par maximum, on aurait retenu $\Delta U = 0{,}4\ \text{V}$ : l'approche par somme est plus prudente.
\end{exoresolu}

\begin{methode}[Estimer l'incertitude d'une mesure unique : la démarche en quatre étapes]
\begin{enumerate}
\item \textbf{Identifier l'instrument} : est-il analogique (échelle graduée), numérique (afficheur) ou de classe (indication normalisée sur le cadran) ?
\item \textbf{Relever la caractéristique pertinente} : valeur $a$ d'une graduation, résolution $r$ de l'affichage, ou couple (classe $k$, calibre $C$).
\item \textbf{Appliquer la règle adaptée} : $\Delta x = a/2$ (lecture soigneuse) ou $a$ (conditions difficiles) pour l'analogique ; $\Delta x = r$ pour le numérique sans notice ; $\Delta x = k \times C / 100$ pour l'appareil de classe. Si plusieurs sources coexistent, les additionner.
\item \textbf{Rédiger le résultat} : arrondir $\Delta x$ à un chiffre significatif par excès, aligner la valeur lue sur la même décimale, écrire $x = x_{\text{lu}} \pm \Delta x$ avec l'unité commune.
\end{enumerate}
\end{methode}

\begin{savaistu}[La métrologie légale protège le consommateur camerounais]
Au marché Mokolo à Yaoundé ou au marché central de Douala, la balance du commerçant est un instrument de mesure comme un autre : elle possède une résolution, une justesse, une incertitude. Pour qu'un kilogramme de tomates payé soit bien un kilogramme, les fabricants et commerçants doivent faire vérifier périodiquement leurs instruments par les services de la \textbf{métrologie légale} (au Cameroun, sous l'égide de l'Agence des Normes et de la Qualité, ANOR). Un instrument conforme reçoit une vignette ou un poinçon de vérification ; un instrument faussé expose son utilisateur à des sanctions. La métrologie n'est donc pas qu'une affaire de laboratoire : elle garantit l'équité des échanges commerciaux, la sécurité des dosages médicaux et la fiabilité des factures d'électricité d'ENEO !
\end{savaistu}

\begin{attention}[Bilan des erreurs à ne plus commettre]
\begin{itemize}
\item Confondre \textbf{résolution} (dernier chiffre affiché) et \textbf{graduation} (plus petite division d'une échelle) : ce sont les analogues l'un de l'autre, mais sur des instruments différents.
\item Oublier que pour un appareil de classe, l'incertitude se calcule sur le \textbf{calibre} et non sur la valeur lue.
\item Écrire un résultat sans unité, ou avec des décimales incohérentes entre valeur et incertitude : $(24{,}3 \pm 0{,}05)\ \text{cm}$ est à proscrire.
\item Croire qu'une mesure unique peut être « exacte » : toute mesure, même unique, doit être accompagnée de son incertitude pour avoir une valeur scientifique.
\end{itemize}
\end{attention}

Tu sais désormais estimer et rédiger l'incertitude d'une mesure faite une seule fois, selon la nature de l'instrument. Mais lorsque l'expérience le permet, répéter la mesure ouvre la porte à un traitement bien plus puissant : l'étude statistique d'une série de mesures, avec la moyenne, l'écart-type et l'incertitude-type. C'est l'objet des sections suivantes.

\coursec{Série de mesures : moyenne et écart-type}

Dans la section précédente, nous avons vu comment estimer l'incertitude lorsqu'on ne réalise qu'\emph{une seule} mesure : on se contente alors de la résolution de l'instrument et de la notice du constructeur. Mais un scientifique sérieux — et un élève de Première C au banc du TP ! — ne se satisfait jamais d'une mesure unique. En répétant la mesure plusieurs fois dans les mêmes conditions, on obtient une \cle{série de mesures} dont l'exploitation statistique (moyenne, écart-type) fournit une estimation bien plus fiable du mesurande. C'est l'objet de cette section.

\courssub{Pourquoi répéter une mesure ?}

\begin{experience}[Le pendule du laboratoire]
Au laboratoire du lycée, on mesure la période $T$ d'un pendule simple avec un chronomètre. Un élève déclenche le chronomètre au passage du pendule devant son repère et l'arrête après une oscillation complète. Il recommence dix fois. \textbf{Observation} : les dix valeurs ne sont \emph{jamais} identiques ; on obtient par exemple : 1,98 s ; 2,03 s ; 1,97 s ; 2,01 s ; 2,00 s ; 1,99 s ; 2,02 s ; 1,98 s ; 2,00 s ; 2,01 s. Pourtant le pendule est le même, la longueur du fil n'a pas changé ! \textbf{Interprétation} : le temps de réaction de l'élève, la lecture du repère, les petites vibrations du support introduisent des fluctuations imprévisibles, tantôt en plus, tantôt en moins : ce sont les \cle{erreurs aléatoires}.
\end{experience}

Contrairement aux erreurs \emph{systématiques} (un chronomètre qui avance, un zéro mal réglé décalent \emph{toutes} les mesures dans le même sens), les erreurs aléatoires se compensent \emph{en moyenne} : certaines mesures sont trop grandes, d'autres trop petites. C'est précisément cette compensation qui rend la répétition des mesures si puissante.

\begin{propriete}[Intérêt de la répétition des mesures — admise]
Soit $n$ mesures indépendantes $x_1, x_2, \ldots, x_n$ d'un même mesurande, effectuées dans les mêmes conditions. La \textbf{moyenne} de ces $n$ mesures est une \textbf{meilleure estimation} de la valeur du mesurande qu'une mesure unique : les erreurs aléatoires, tantôt positives tantôt négatives, se compensent partiellement dans la moyenne. Plus $n$ est grand, meilleure est l'estimation (la fluctuation de la moyenne décroît comme $1/\sqrt{n}$).\par
\emph{Justification qualitative} : si chaque mesure vaut $x_i = x_{\text{vrai}} + \varepsilon_i$ où $\varepsilon_i$ est l'erreur aléatoire, alors la moyenne vaut $x_{\text{vrai}} + \bar{\varepsilon}$. Or les $\varepsilon_i$ étant de signes variés, leur moyenne $\bar{\varepsilon}$ est généralement bien plus petite en valeur absolue que chaque $\varepsilon_i$ pris isolément.
\end{propriete}

\begin{attention}[Ce que la répétition ne corrige pas]
Répéter les mesures réduit l'effet des erreurs \textbf{aléatoires}, mais ne supprime \textbf{jamais} une erreur \textbf{systématique} : si le chronomètre retarde de 0,05 s, la moyenne de 100 mesures retardera encore de 0,05 s ! Il faut alors corriger l'instrument ou le protocole (étalonnage, mesure de 10 périodes au lieu d'une, etc.).
\end{attention}

\courssub{Organisation des données : le tableau des mesures}

La première tâche du scientifique est d'\textbf{organiser} les données. On numérote les mesures et on les consigne dans un tableau, sans en oublier aucune et sans arrondir prématurément.

\begin{center}
\begin{tabularx}{\linewidth}{|l|X|X|X|X|X|}
\hline
\rowcolor{popblueL} Numéro $i$ & 1 & 2 & 3 & 4 & 5 \\
\hline
Mesure $x_i$ & $x_1$ & $x_2$ & $x_3$ & $x_4$ & $x_5$ \\
\hline
Écart $x_i - \bar{x}$ & $x_1-\bar{x}$ & $x_2-\bar{x}$ & $x_3-\bar{x}$ & $x_4-\bar{x}$ & $x_5-\bar{x}$ \\
\hline
Écart au carré $(x_i-\bar{x})^2$ & $(x_1-\bar{x})^2$ & $(x_2-\bar{x})^2$ & $(x_3-\bar{x})^2$ & $(x_4-\bar{x})^2$ & $(x_5-\bar{x})^2$ \\
\hline
\end{tabularx}
\end{center}

Les deux dernières lignes seront remplies \emph{après} le calcul de la moyenne : elles serviront au calcul de l'écart-type. Toutes les valeurs d'une même série doivent être exprimées \textbf{dans la même unité} et avec le \textbf{même nombre de décimales} (celui permis par l'instrument).

\courssub{La moyenne arithmétique : meilleure estimation du mesurande}

\begin{definition}[Moyenne d'une série de mesures]
La \cle{moyenne arithmétique} d'une série de $n$ mesures $x_1, x_2, \ldots, x_n$ d'un même mesurande est :
\[ \bar{x} = \frac{x_1 + x_2 + \cdots + x_n}{n} = \frac{1}{n}\sum_{i=1}^{n} x_i \]
La moyenne $\bar{x}$ constitue la \textbf{meilleure estimation} de la valeur du mesurande. Elle s'exprime dans la même unité que les mesures $x_i$.
\end{definition}

\textbf{Vérification dimensionnelle} : $\bar{x}$ est une somme de grandeurs homogènes divisée par un nombre sans dimension ; $\bar{x}$ a donc bien la même unité que les $x_i$.

\begin{exemplebox}[Moyenne de cinq périodes]
Cinq mesures de la période d'un pendule donnent (en secondes) : 1,98 ; 2,02 ; 2,00 ; 1,99 ; 2,01.
\[ \bar{T} = \frac{1,98 + 2,02 + 2,00 + 1,99 + 2,01}{5} = \frac{10,00}{5} = \SI{2,000}{s} \]
On peut donner la moyenne avec une décimale de plus que les mesures brutes, soit $\SI{2,000}{s}$.
\end{exemplebox}

\courssub{L'écart-type expérimental : mesurer la dispersion}

La moyenne ne dit pas tout. Imaginons deux séries de mesures de même moyenne $\bar{T} = \SI{2,000}{s}$ : la première donne 1,99 ; 2,00 ; 2,01 ; 2,00 ; 2,00 et la seconde 1,80 ; 2,20 ; 1,95 ; 2,05 ; 2,00. La première série est manifestement plus \emph{resserrée} : le protocole y est plus stable, l'instrument plus \emph{fidèle}. Il faut un nombre qui quantifie cette \textbf{dispersion} : c'est l'\cle{écart-type expérimental}.

\textbf{Construction de l'indicateur.} Pour chaque mesure, on calcule l'\textbf{écart} à la moyenne $x_i - \bar{x}$. On ne peut pas simplement faire la moyenne des écarts : ils sont tantôt positifs, tantôt négatifs, et leur somme vaut exactement zéro (vérification : $\sum (x_i - \bar{x}) = \sum x_i - n\bar{x} = 0$ par définition de $\bar{x}$). On élève donc chaque écart au carré — ce qui les rend tous positifs et accentue les grands écarts —, on fait la moyenne de ces carrés, puis on prend la \textbf{racine carrée} pour revenir à l'unité de la mesure.

\begin{definition}[Écart-type expérimental]
L'\cle{écart-type expérimental} d'une série de $n$ mesures $x_1, \ldots, x_n$ de moyenne $\bar{x}$ est :
\[ s = \sqrt{\frac{\displaystyle\sum_{i=1}^{n}\left(x_i - \bar{x}\right)^2}{n-1}} \]
Il s'exprime dans la même unité que les mesures et quantifie leur \textbf{dispersion} autour de la moyenne : plus $s$ est petit, plus les mesures sont resserrées (instrument \textbf{fidèle}, protocole stable) ; plus $s$ est grand, plus les mesures sont dispersées.
\end{definition}

\begin{propriete}[Pourquoi diviser par $n-1$ et non par $n$ ? — admise, justification qualitative]
La moyenne $\bar{x}$ qui apparaît dans la formule n'est pas la valeur vraie du mesurande : elle est \emph{calculée à partir des mêmes données} $x_i$. Or $\bar{x}$ est, parmi tous les nombres, celui qui rend la somme $\sum (x_i - a)^2$ \emph{minimale} : les écarts à $\bar{x}$ sont donc systématiquement un peu plus petits que les écarts à la valeur vraie. On dit qu'on a « consommé un degré de liberté » en estimant $\bar{x}$. Diviser par $n-1$ au lieu de $n$ corrige cette sous-estimation de la dispersion. \emph{Remarque} : pour $n$ grand ($n \geq 30$), la différence entre $n$ et $n-1$ devient négligeable. Cette formule est admise en Première ; aucune démonstration n'est exigible au Probatoire.
\end{propriete}

\begin{exoresolu}[Calcul complet : moyenne et écart-type de cinq mesures de période]
\textbf{Énoncé.} Un élève mesure cinq fois la période d'un pendule simple et obtient (en s) : 1,98 ; 2,02 ; 2,00 ; 1,99 ; 2,01. Déterminer la moyenne $\bar{T}$ et l'écart-type expérimental $s$ de cette série. Commenter.
\tcblower
\textbf{Solution détaillée.}\par
\textbf{Étape 1 — Moyenne.}
\[ \bar{T} = \frac{1,98 + 2,02 + 2,00 + 1,99 + 2,01}{5} = \frac{10,00}{5} = \SI{2,000}{s} \]
\textbf{Étape 2 — Tableau des écarts et de leurs carrés.}
\begin{center}
\begin{tabular}{|c|c|c|c|}
\hline
\rowcolor{popblueL} $i$ & $T_i$ (s) & $T_i - \bar{T}$ (s) & $(T_i - \bar{T})^2$ (s$^2$) \\
\hline
1 & 1,98 & $-0,02$ & 0,0004 \\
\hline
2 & 2,02 & $+0,02$ & 0,0004 \\
\hline
3 & 2,00 & $0,00$ & 0,0000 \\
\hline
4 & 1,99 & $-0,01$ & 0,0001 \\
\hline
5 & 2,01 & $+0,01$ & 0,0001 \\
\hline
\multicolumn{3}{|c|}{\textbf{Somme}} & \textbf{0,0010} \\
\hline
\end{tabular}
\end{center}
On vérifie au passage que la somme des écarts vaut $-0,02 + 0,02 + 0 + (-0,01) + 0,01 = 0$ : bon signe, le calcul de $\bar{T}$ est cohérent.\par
\textbf{Étape 3 — Écart-type.} Ici $n = 5$, donc $n - 1 = 4$ :
\[ s = \sqrt{\frac{0,0010}{4}} = \sqrt{2,5 \times 10^{-4}} \approx \SI{0,016}{s} \]
\textbf{Étape 4 — Commentaire.} L'écart-type $s \approx \SI{0,016}{s}$ est petit devant $\bar{T} = \SI{2,000}{s}$ (environ 0,8 \% de la moyenne) : les mesures sont resserrées, le protocole est stable et le chronomètre est fidèle. La meilleure estimation de la période est $\bar{T} = \SI{2,000}{s}$, avec une dispersion typique de l'ordre de $\SI{0,02}{s}$.
\end{exoresolu}

\begin{attention}[Les deux erreurs classiques au Probatoire]
\begin{itemize}
\item \textbf{Oublier la racine carrée} : après avoir calculé $\dfrac{\sum (x_i - \bar{x})^2}{n-1} = 2,5 \times 10^{-4}\ \text{s}^2$, beaucoup d'élèves s'arrêtent là. Ce nombre est la \emph{variance}, pas l'écart-type ! Indice pour s'en apercevoir : son unité est le s$^2$, alors que l'écart-type doit s'exprimer en s.
\item \textbf{Diviser par $n$ au lieu de $n-1$} : la formule officielle de l'écart-type \emph{expérimental} (celui d'une série de mesures) comporte $n-1$ au dénominateur. Avec $n = 5$, diviser par 5 donnerait $s \approx \SI{0,014}{s}$ au lieu de $\SI{0,016}{s}$ : c'est faux.
\end{itemize}
Autre piège : calculer les écarts par rapport à une valeur « attendue » (par exemple $T = \SI{2,00}{s}$ annoncée) au lieu de la moyenne $\bar{T}$ de la série.
\end{attention}

\courssub{Représentation graphique d'une série de mesures}

Un tableau de nombres parle peu à l'œil ; un graphique révèle instantanément la structure des données. Deux représentations sont au programme.

\textbf{1. Le nuage de points (mesure par mesure).} On porte en abscisse le numéro $i$ de la mesure et en ordonnée la valeur $x_i$. On trace ensuite la droite horizontale $y = \bar{x}$, ainsi que les droites $y = \bar{x} - s$ et $y = \bar{x} + s$ qui délimitent une \textbf{bande de dispersion}. On constate expérimentalement qu'environ \textbf{deux tiers} des points d'une série « normale » se situent dans la bande $[\bar{x} - s \,;\, \bar{x} + s]$ : c'est une excellente vérification visuelle de la cohérence de la série.

\begin{popfigure}
\begin{center}
\begin{tikzpicture}
\begin{axis}[
  width=0.92\linewidth, height=6.2cm,
  xlabel={Numéro de la mesure $i$},
  ylabel={$T_i$ (s)},
  xmin=0.4, xmax=10.6,
  ymin=1.955, ymax=2.045,
  xtick={1,2,3,4,5,6,7,8,9,10},
  ytick={1.96,1.98,2.00,2.02,2.04},
  grid=both, grid style={popline!40},
  legend style={at={(0.02,0.03)}, anchor=south west, font=\small, draw=popline}
]
% Bande xbar ± s : xbar = 2.000, s = 0.016
\addplot[draw=none, fill=popblue, fill opacity=0.15, forget plot]
  coordinates {(0.4,1.984) (10.6,1.984) (10.6,2.016) (0.4,2.016)} \closedcycle;
% Droite moyenne
\addplot[popblue, thick, domain=0.4:10.6] {2.000};
\addlegendentry{$y = \bar{T} = \SI{2,000}{s}$}
% Droites xbar ± s
\addplot[poporange, dashed, thick, domain=0.4:10.6] {2.016};
\addlegendentry{$y = \bar{T} + s$}
\addplot[poporange, dashed, thick, domain=0.4:10.6] {1.984};
\addlegendentry{$y = \bar{T} - s$}
% Nuage de 10 mesures
\addplot[only marks, mark=*, mark size=2.2pt, popdark] coordinates {
  (1,1.98) (2,2.02) (3,2.00) (4,1.99) (5,2.01)
  (6,1.97) (7,2.03) (8,2.00) (9,1.99) (10,2.01)
};
\addlegendentry{Mesures $T_i$}
\end{axis}
\end{tikzpicture}
\end{center}
\legende{Nuage de 10 mesures de la période d'un pendule autour de la droite moyenne $y = \bar{T}$. La bande colorée représente l'intervalle $[\bar{T} - s \,;\, \bar{T} + s]$ : on observe qu'environ deux tiers des points s'y trouvent, ce qui est la signature d'une série de mesures saine.}
\end{popfigure}

\textbf{2. L'histogramme.} Lorsque le nombre de mesures devient grand ($n \geq 15$ environ), on regroupe les valeurs en \textbf{classes} de même largeur et on porte en ordonnée l'\textbf{effectif} de chaque classe. L'histogramme révèle la \emph{forme} de la distribution : pour des erreurs purement aléatoires, on obtient une cloche à peu près symétrique centrée sur $\bar{x}$, dont la largeur est de l'ordre de quelques $s$.

\begin{popfigure}
\begin{center}
\begin{tikzpicture}
\begin{axis}[
  width=0.92\linewidth, height=6.5cm,
  ybar, bar width=14pt,
  xlabel={Période $T$ (s)},
  ylabel={Effectif},
  xmin=1.94, xmax=2.06,
  ymin=0, ymax=8.5,
  xtick={1.95,1.97,1.99,2.01,2.03,2.05},
  ytick={0,1,2,3,4,5,6,7,8},
  grid=major, grid style={popline!40},
]
\addplot[ybar, fill=popblue!70, draw=popblue] coordinates {
  (1.95,1) (1.97,3) (1.99,6) (2.01,7) (2.03,2) (2.05,1)
};
% Ligne verticale à la moyenne
\addplot[poporange, very thick, dashed] coordinates {(2.00,0) (2.00,8.2)};
\node[poporange, anchor=south west, font=\small] at (axis cs:2.002,7.6) {$\bar{T} = \SI{2,000}{s}$};
\end{axis}
\end{tikzpicture}
\end{center}
\legende{Histogramme de 20 mesures de la période d'un pendule, regroupées en classes de largeur 0,02 s (effectif total : $1+3+6+7+2+1 = 20$). La distribution est à peu près symétrique autour de la moyenne : c'est l'allure typique d'erreurs aléatoires.}
\end{popfigure}

\courssub{Détection et traitement d'une valeur aberrante}

\begin{definition}[Valeur aberrante]
Une mesure $x_k$ est dite \cle{aberrante} si elle s'écarte nettement de toutes les autres, sans raison physique apparente : typiquement si $\left| x_k - \bar{x} \right|$ est supérieur à $2s$ ou $3s$ alors que toutes les autres mesures sont resserrées autour de $\bar{x}$.
\end{definition}

Sur le nuage de points, une valeur aberrante saute aux yeux : c'est le point isolé, loin de la bande $[\bar{x} - s \,;\, \bar{x} + s]$. Par exemple, si dans notre série de périodes apparaît la valeur $T = \SI{2,35}{s}$, elle est à plus de $20\,s$ de la moyenne : l'élève a probablement compté une oscillation de trop, ou déclenché le chronomètre en retard.

\begin{methode}[Conduite à tenir face à une valeur aberrante]
\begin{enumerate}
\item \textbf{Ne jamais supprimer silencieusement} une mesure : toute suppression doit être \emph{signalée et justifiée} dans le compte rendu.
\item \textbf{Chercher une cause} : relecture du protocole, erreur de comptage, mauvais déclenchement, choc sur la paillasse, coupure de courant...
\item Si une \textbf{cause expérimentale identifiable} est trouvée (faute de manipulation avérée), on peut \textbf{écarter} la valeur, \emph{en le précisant explicitement}, puis recalculer $\bar{x}$ et $s$ sur les $n-1$ mesures restantes.
\item Si \textbf{aucune cause} n'est identifiée, on \textbf{conserve} la mesure : elle fait partie de l'information expérimentale. On peut éventuellement refaire une série de mesures.
\end{enumerate}
\end{methode}

\begin{attention}[Honnêteté scientifique]
Écarter une mesure uniquement parce qu'elle « ne plaît pas » ou parce qu'elle éloigne le résultat de la valeur attendue est une \textbf{faute scientifique grave} (c'est de la fraude !). Seule une cause expérimentale clairement identifiée justifie l'élimination d'une donnée, et toujours de manière transparente.
\end{attention}

\begin{methode}[Organiser et traiter une série de mesures — la démarche complète]
\begin{enumerate}
\item \textbf{Tableau} : consigner les $n$ mesures $x_1, \ldots, x_n$ dans un tableau, dans la même unité, avec le même nombre de décimales.
\item \textbf{Moyenne} : calculer $\bar{x} = \dfrac{1}{n}\displaystyle\sum_{i=1}^{n} x_i$, meilleure estimation du mesurande.
\item \textbf{Écarts} : compléter le tableau avec les écarts $x_i - \bar{x}$ (vérifier que leur somme est nulle) puis leurs carrés $(x_i - \bar{x})^2$.
\item \textbf{Écart-type} : calculer $s = \sqrt{\dfrac{\sum (x_i - \bar{x})^2}{n-1}}$ — attention à la racine carrée et au $n-1$.
\item \textbf{Représentation graphique} : tracer le nuage de points avec la droite $y = \bar{x}$ et la bande $[\bar{x}-s \,;\, \bar{x}+s]$ (ou un histogramme si $n$ est grand).
\item \textbf{Contrôle} : repérer d'éventuelles valeurs aberrantes, en rechercher la cause, et conclure sur la fidélité de la série ($s$ petit $\Rightarrow$ mesures resserrées).
\end{enumerate}
\end{methode}

\begin{aretenir}[L'essentiel de la section]
\begin{itemize}
\item Répéter une mesure $n$ fois permet de \textbf{réduire l'effet des erreurs aléatoires} (mais pas des erreurs systématiques).
\item La \textbf{moyenne} $\bar{x} = \dfrac{1}{n}\sum x_i$ est la \textbf{meilleure estimation} de la valeur du mesurande.
\item L'\textbf{écart-type expérimental} $s = \sqrt{\dfrac{\sum (x_i - \bar{x})^2}{n-1}}$ mesure la \textbf{dispersion} : $s$ petit $\Leftrightarrow$ mesures resserrées $\Leftrightarrow$ instrument fidèle et protocole stable.
\item On divise par $n-1$ (et non $n$) car $\bar{x}$ est calculé à partir des mêmes données, ce qui sous-estimerait la dispersion.
\item Le nuage de points avec la droite $y = \bar{x}$ et la bande $\bar{x} \pm s$ (ou l'histogramme pour $n$ grand) permet de visualiser la série et de détecter les valeurs aberrantes, qu'on n'écarte qu'avec une justification explicite.
\end{itemize}
\end{aretenir}

\begin{savaistu}[Des élections aux étoiles]
Les sondages électoraux utilisent exactement la même logique : interroger une personne, c'est une « mesure unique » très incertaine ; interroger 1000 personnes, c'est faire une moyenne dont la fluctuation est environ $\sqrt{1000} \approx 32$ fois plus petite. En astronomie, les astronomes du \emph{Centre de recherche} comme ceux qui suivent les satellites de télécommunications (utilisés par CAMTEL ou les opérateurs mobiles) répètent des centaines de fois la mesure de la position d'un astre : la moyenne gagne en précision à chaque nouvelle observation, exactement comme tes mesures de période au laboratoire.
\end{savaistu}

\begin{exercice}[À toi de jouer]
Un groupe d'élèves mesure six fois la masse d'un même sac d'engrais avec une balance du marché de Mokolo : $m$ (kg) : 25,2 ; 24,9 ; 25,1 ; 25,0 ; 26,8 ; 25,0.
\begin{enumerate}
\item Calculer la moyenne $\bar{m}$ et l'écart-type expérimental $s$ des six mesures.
\item Représenter mentalement le nuage de points : quelle mesure semble aberrante ?
\item Après vérification, l'élève responsable de la $5^{\text{ème}}$ mesure avoue avoir lu la balance de biais. Écarter cette valeur et recalculer $\bar{m}$ et $s$. Conclure sur la fidélité de la balance.
\end{enumerate}
\end{exercice}

\begin{corrige}[Exercice : série de masses]
\textbf{1.} Moyenne des six mesures : $\bar{m} = \dfrac{25,2 + 24,9 + 25,1 + 25,0 + 26,8 + 25,0}{6} = \dfrac{152,0}{6} \approx \SI{25,33}{kg}$.\par
Écarts : $-0,13$ ; $-0,43$ ; $-0,23$ ; $-0,33$ ; $+1,47$ ; $-0,33$ (somme $\approx 0$, correct). Carrés : $0,0169$ ; $0,1849$ ; $0,0529$ ; $0,1089$ ; $2,1609$ ; $0,1089$ ; somme $\approx 2,6334$. D'où $s = \sqrt{\dfrac{2,6334}{5}} \approx \sqrt{0,527} \approx \SI{0,73}{kg}$.\par
\textbf{2.} La mesure $m_5 = \SI{26,8}{kg}$ s'écarte de $\bar{m}$ de $\SI{1,47}{kg} \approx 2s$ : elle est aberrante.\par
\textbf{3.} Sur les cinq mesures restantes : $\bar{m} = \dfrac{125,2}{5} = \SI{25,04}{kg}$. Écarts : $+0,16$ ; $-0,14$ ; $+0,06$ ; $-0,04$ ; $-0,04$ ; carrés : $0,0256$ ; $0,0196$ ; $0,0036$ ; $0,0016$ ; $0,0016$ ; somme $= 0,052$. D'où $s = \sqrt{\dfrac{0,052}{4}} = \sqrt{0,013} \approx \SI{0,11}{kg}$.\par
\textbf{Conclusion} : après élimination justifiée de la valeur aberrante, l'écart-type passe de $\SI{0,73}{kg}$ à $\SI{0,11}{kg}$ : les mesures sont resserrées, la balance est fidèle, et la meilleure estimation de la masse est $\bar{m} = \SI{25,04}{kg}$.
\end{corrige}

La moyenne et l'écart-type décrivent la série de mesures, mais ils ne répondent pas encore à la question décisive : \emph{quelle incertitude attacher au résultat final $\bar{x}$ ?} C'est précisément l'objet de la section suivante, où nous définirons l'\textbf{incertitude-type} sur la moyenne, l'\textbf{incertitude élargie} et l'\textbf{intervalle de confiance}.

\coursec{Incertitude-type, incertitude élargie et intervalle de confiance}

Dans la section précédente, tu as appris à résumer une série de $n$ mesures par deux nombres : la \cle{moyenne} $\bar{x}$, meilleure estimation de la valeur mesurée, et l'\cle{écart-type} $s$, qui décrit la dispersion des mesures autour de cette moyenne. Mais une question demeure, et c'est toute la question du métrologue : \emph{à quel point puis-je faire confiance à cette moyenne ?} Si je recommence toute la campagne de mesures demain, avec le même voltmètre et le même circuit, vais-je retrouver $\bar{x} = \SI{5,20}{\volt}$, ou bien $\SI{5,18}{\volt}$, ou $\SI{5,23}{\volt}$ ? Répondre à cette question, c'est construire les notions d'\cle{incertitude-type}, d'\cle{incertitude élargie} et d'\cle{intervalle de confiance}, qui sont au cœur de l'expression de tout résultat scientifique — et de nombreux sujets de Probatoire.

\courssub{Incertitude-type : mesurer la confiance que l'on accorde à la moyenne}

\textbf{Intuition.} Imagine deux élèves d'une classe de Première C au lycée de Ngaoundéré qui mesurent chacun dix fois la tension aux bornes d'une pile. Le premier obtient des mesures très dispersées ($s = \SI{0,15}{\volt}$), le second des mesures très groupées ($s = \SI{0,02}{\volt}$). À moyenne égale, on fait évidemment plus confiance au second. Mais il y a un deuxième facteur : si chacun avait fait \emph{cent} mesures au lieu de dix, leurs moyennes seraient encore plus fiables. L'incertitude sur la moyenne dépend donc de \textbf{deux ingrédients} : la dispersion $s$ des mesures, et le nombre $n$ de mesures.

\begin{definition}[Incertitude-type]
L'\cle{incertitude-type}, notée $u(x)$, est une incertitude exprimée \textbf{sous la forme d'un écart-type}. Elle caractérise la dispersion des valeurs que l'on peut raisonnablement attribuer au mesurande. Pour une série de $n$ mesures indépendantes, on distingue :
\begin{itemize}
\item l'écart-type $s$, qui mesure la dispersion des \textbf{mesures individuelles} ;
\item l'incertitude-type sur la moyenne $u(\bar{x})$, qui mesure la dispersion de la \textbf{moyenne} si l'on recommençait toute la série.
\end{itemize}
\end{definition}

\begin{propriete}[Incertitude-type sur la moyenne (admise)]
Pour $n$ mesures indépendantes $x_1, x_2, \ldots, x_n$ de moyenne $\bar{x}$ et d'écart-type $s$, l'incertitude-type sur la moyenne est :
\[ u(\bar{x}) = \frac{s}{\sqrt{n}} \]
Cette formule est \textbf{admise} au niveau Première (sa démonstration relève des propriétés de la variance d'une somme de variables aléatoires indépendantes). Elle n'est pas exigible en démonstration au Probatoire, mais son \textbf{utilisation} est exigible.
\end{propriete}

\textbf{Interprétation physique du facteur $\sqrt{n}$.} La moyenne fluctue \emph{moins} que les mesures individuelles : les erreurs aléatoires, tantôt positives, tantôt négatives, se compensent partiellement dans la somme. Diviser par $\sqrt{n}$ traduit exactement cette compensation. Conséquence pratique importante : pour diviser l'incertitude par $2$, il faut multiplier le nombre de mesures par $4$ ; pour la diviser par $10$, il en faut $100$ ! C'est pourquoi, au-delà d'une dizaine de mesures, il devient souvent plus rentable d'améliorer l'instrument que d'accumuler les mesures.

\begin{attention}[Ne pas confondre $s$ et $u(\bar{x})$]
C'est l'erreur la plus fréquente au Probatoire. L'\textbf{écart-type} $s$ décrit la dispersion des mesures individuelles : il ne diminue \emph{pas} quand $n$ augmente, il se stabilise vers une valeur caractéristique de l'instrument et du protocole. L'\textbf{incertitude-type} $u(\bar{x}) = s/\sqrt{n}$ décrit l'incertitude sur la \emph{moyenne} : elle diminue en $1/\sqrt{n}$. Écrire « l'incertitude sur le résultat est $s = \SI{0,06}{\volt}$ » alors qu'on a fait $10$ mesures, c'est surestimer l'incertitude d'un facteur $\sqrt{10} \approx 3,2$.
\end{attention}

\courssub{Du besoin d'élargir : le facteur d'élargissement $k$}

L'incertitude-type $u(\bar{x})$ est un écart-type : elle définit un intervalle $[\bar{x} - u \,;\, \bar{x} + u]$ qui, pour une distribution gaussienne, ne contient la valeur vraie qu'avec une probabilité d'environ $68\,\%$. C'est insuffisant pour un résultat officiel : un laboratoire qui étalonne les compteurs d'ENEO, ou un hôpital qui étalonne une source de radiothérapie, exige un niveau de confiance bien plus élevé. On « élargit » donc l'intervalle en multipliant l'incertitude-type par un facteur $k$.

\begin{definition}[Facteur d'élargissement]
Le \cle{facteur d'élargissement}, noté $k$, est le nombre par lequel on multiplie l'incertitude-type pour obtenir l'incertitude élargie. Il dépend du \textbf{niveau de confiance} choisi et de la loi de probabilité supposée pour les fluctuations.
\end{definition}

\begin{propriete}[Valeurs usuelles de $k$ pour une loi normale (admise)]
Lorsque les fluctuations de la moyenne suivent une loi normale (gaussienne), ce qui est l'hypothèse standard en métrologie :
\begin{itemize}
\item $k = 1$ correspond à un niveau de confiance d'environ $68\,\%$ ;
\item $k = 2$ correspond à un niveau de confiance d'environ $95\,\%$ ;
\item $k = 3$ correspond à un niveau de confiance d'environ $99\,\%$.
\end{itemize}
Ces valeurs sont \textbf{admises}. Au Probatoire, lorsque le niveau de confiance est de $95\,\%$, on utilise \textbf{systématiquement} $k = 2$.
\end{propriete}

\begin{popfigure}
\begin{center}
\begin{tikzpicture}
\begin{axis}[
  width=0.95\linewidth, height=6.2cm,
  domain=-4:4, samples=200,
  axis lines=middle,
  xlabel={$x - \bar{x}$}, ylabel={densité},
  xtick={-3,-2,-1,0,1,2,3},
  xticklabels={$\bar{x}-3u$,$\bar{x}-2u$,$\bar{x}-u$,$\bar{x}$,$\bar{x}+u$,$\bar{x}+2u$,$\bar{x}+3u$},
  ytick=\empty,
  x label style={at={(ticklabel* cs:1)},anchor=west},
  every axis x label/.style={at={(current axis.right of origin)},anchor=west},
  clip=false
]
% zone 99 %
\addplot[fill=popgreenL, draw=none, domain=-3:3] {exp(-x^2/2)} \closedcycle;
% zone 95 %
\addplot[fill=popblueL, draw=none, domain=-2:2] {exp(-x^2/2)} \closedcycle;
% zone 68 %
\addplot[fill=poporangeL, draw=none, domain=-1:1] {exp(-x^2/2)} \closedcycle;
% courbe
\addplot[popcurve, domain=-4:4] {exp(-x^2/2)};
% annotations
\node[popdark, font=\small] at (axis cs:0,0.22) {$\approx 68\,\%$};
\node[popblue, font=\small] at (axis cs:1.55,0.10) {$\approx 95\,\%$};
\node[popgreen!60!black, font=\small] at (axis cs:2.6,0.035) {$\approx 99\,\%$};
\end{axis}
\end{tikzpicture}
\end{center}
\legende{Courbe en cloche (loi normale) centrée sur la moyenne $\bar{x}$. Les zones colorées correspondent aux intervalles $\bar{x} \pm u$ (environ $68\,\%$), $\bar{x} \pm 2u$ (environ $95\,\%$) et $\bar{x} \pm 3u$ (environ $99\,\%$). Plus on veut être confiant, plus l'intervalle est large.}
\end{popfigure}

\courssub{Incertitude élargie et intervalle de confiance}

\begin{definition}[Incertitude élargie]
L'\cle{incertitude élargie}, notée $U$, est le produit de l'incertitude-type par le facteur d'élargissement :
\[ U = k \times u(\bar{x}) = k \, \frac{s}{\sqrt{n}} \]
Elle définit un intervalle autour du résultat auquel on attribue un niveau de confiance donné. Elle s'exprime dans la même unité que la grandeur mesurée.
\end{definition}

\begin{definition}[Intervalle de confiance]
L'\cle{intervalle de confiance} au niveau de confiance associé à $k$ est :
\[ \left[ \, \bar{x} - U \ ;\ \bar{x} + U \, \right] \]
La probabilité que la \textbf{valeur vraie} du mesurande se trouve dans cet intervalle est égale au niveau de confiance (environ $95\,\%$ pour $k = 2$, environ $99\,\%$ pour $k = 3$).
\end{definition}

\textbf{Interprétation concrète d'un intervalle à $95\,\%$.} Cette phrase est subtile et souvent mal comprise. Dire « l'intervalle $[\SI{5,16}{\volt} \,;\, \SI{5,24}{\volt}]$ contient la valeur vraie avec $95\,\%$ de confiance » signifie : si l'on recommençait $100$ campagnes de mesures identiques (même protocole, même instrument, $n$ mesures à chaque fois), on obtiendrait $100$ moyennes légèrement différentes et donc $100$ intervalles légèrement différents — et environ $95$ de ces intervalles contiendraient la valeur vraie. Ce n'est pas la valeur vraie qui « bouge » : c'est notre intervalle qui fluctue d'une campagne à l'autre, et la procédure réussit $95$ fois sur $100$.

\begin{popfigure}
\begin{center}
\begin{tikzpicture}[scale=1.0]
% axe gradué
\draw[->, thick, popdark] (0,0) -- (11,0) node[below left=2pt and -30pt] {};
\foreach \x/\lab in {1/5{,}10, 3/5{,}16, 5.5/5{,}20, 8/5{,}24, 10/5{,}30}
  \draw[thick, popdark] (\x,0.12) -- (\x,-0.12) node[below] {\small $\lab$ V};
% intervalle de confiance
\draw[line width=3.5pt, popblue] (3,0.55) -- (8,0.55);
\draw[line width=1.2pt, popblue] (3,0.35) -- (3,0.75);
\draw[line width=1.2pt, popblue] (8,0.35) -- (8,0.75);
% moyenne
\fill[poporange] (5.5,0.55) circle (2.6pt);
\node[poporange, above] at (5.5,0.8) {\small $\bar{x} = \SI{5,20}{\volt}$};
% bornes
\node[popblue, above, align=center] at (3,1.05) {\small $\bar{x}-U$\\ \small $=\SI{5,16}{\volt}$};
\node[popblue, above, align=center] at (8,1.05) {\small $\bar{x}+U$\\ \small $=\SI{5,24}{\volt}$};
% demi-largeur
\draw[<->, poppurple, thick] (5.5,-0.75) -- (8,-0.75) node[midway, below] {\small $U = \SI{0,04}{\volt}$};
\draw[dashed, poppurple] (5.5,-0.2) -- (5.5,-0.8);
\draw[dashed, poppurple] (8,-0.2) -- (8,-0.8);
\node[popdark, align=center] at (5.5,1.9) {\small Intervalle de confiance à $95\,\%$ : $[\,\bar{x}-U \ ;\ \bar{x}+U\,]$};
\end{tikzpicture}
\end{center}
\legende{Droite graduée représentant le résultat d'une mesure de tension : la moyenne $\bar{x}$, l'incertitude élargie $U$ (demi-largeur de l'intervalle) et les deux bornes de l'intervalle de confiance.}
\end{popfigure}

\begin{methode}[Estimer une incertitude élargie et écrire l'intervalle de confiance]
\begin{enumerate}
\item \textbf{Calculer l'écart-type} $s$ de la série de $n$ mesures (souvent donné par la calculatrice ou fourni dans l'énoncé).
\item \textbf{Calculer l'incertitude-type} sur la moyenne : $u(\bar{x}) = \dfrac{s}{\sqrt{n}}$.
\item \textbf{Choisir le facteur d'élargissement} $k$ selon le niveau de confiance demandé : $k = 2$ pour $95\,\%$, $k = 3$ pour $99\,\%$.
\item \textbf{Calculer l'incertitude élargie} : $U = k \times u(\bar{x})$, arrondie par excès en gardant un (ou deux) chiffre significatif.
\item \textbf{Écrire l'intervalle de confiance} $[\bar{x} - U \,;\, \bar{x} + U]$ et conclure par une phrase : « la valeur vraie se trouve dans cet intervalle avec un niveau de confiance d'environ $95\,\%$ ».
\end{enumerate}
\end{methode}

\begin{exemplebox}[Mesure d'une tension au laboratoire]
Un groupe d'élèves mesure $n = 10$ fois la tension aux bornes d'un générateur et obtient $\bar{x} = \SI{5,20}{\volt}$ avec un écart-type $s = \SI{0,06}{\volt}$. On cherche l'intervalle de confiance à $95\,\%$.
\begin{align*}
u(\bar{x}) &= \frac{s}{\sqrt{n}} = \frac{\num{0,06}}{\sqrt{10}} = \frac{\num{0,06}}{\num{3,16}} \approx \num{0,019}~\text{V} \\
U &= k \times u(\bar{x}) = 2 \times \num{0,019} = \num{0,038}~\text{V} \approx \num{0,04}~\text{V}
\end{align*}
L'intervalle de confiance à $95\,\%$ est donc :
\[ [\, \SI{5,20}{\volt} - \SI{0,04}{\volt} \ ;\ \SI{5,20}{\volt} + \SI{0,04}{\volt} \,] = [\, \SI{5,16}{\volt} \ ;\ \SI{5,24}{\volt} \,] \]
On conclut : la tension vraie du générateur se trouve entre $\SI{5,16}{\volt}$ et $\SI{5,24}{\volt}$, avec un niveau de confiance d'environ $95\,\%$.
\end{exemplebox}

\begin{exoresolu}[Étalonnage d'un multimètre en atelier]
Un technicien d'un atelier de réparation de Douala mesure $n = 16$ fois une tension de référence. Il obtient une moyenne $\bar{U} = \SI{12,02}{\volt}$ et un écart-type $s = \SI{0,08}{\volt}$.
\begin{enumerate}
\item Calculer l'incertitude-type sur la moyenne.
\item Déterminer l'incertitude élargie pour un niveau de confiance de $95\,\%$, puis de $99\,\%$.
\item Écrire l'intervalle de confiance à $95\,\%$ et commenter l'effet du passage à $99\,\%$.
\end{enumerate}
\tcblower
\textbf{Solution.}
\begin{enumerate}
\item Incertitude-type sur la moyenne :
\[ u(\bar{U}) = \frac{s}{\sqrt{n}} = \frac{\num{0,08}}{\sqrt{16}} = \frac{\num{0,08}}{4} = \num{0,02}~\text{V} \]
\item À $95\,\%$, on prend $k = 2$ : $U_{95} = 2 \times \num{0,02} = \SI{0,04}{\volt}$. À $99\,\%$, on prend $k = 3$ : $U_{99} = 3 \times \num{0,02} = \SI{0,06}{\volt}$.
\item Intervalle à $95\,\%$ : $[\, \SI{12,02}{\volt} - \SI{0,04}{\volt} \ ;\ \SI{12,02}{\volt} + \SI{0,04}{\volt} \,] = [\, \SI{11,98}{\volt} \ ;\ \SI{12,06}{\volt} \,]$. À $99\,\%$, l'intervalle serait $[\, \SI{11,96}{\volt} \ ;\ \SI{12,08}{\volt} \,]$ : il est \textbf{plus large}. C'est un résultat général : exiger une confiance plus grande oblige à élargir l'intervalle. On ne peut pas gagner en confiance sans perdre en étroitesse, sauf à augmenter $n$ ou à améliorer l'instrument.
\end{enumerate}
\end{exoresolu}

\begin{attention}[Pièges de rédaction au Probatoire]
\begin{itemize}
\item Ne jamais écrire $U = k \times s$ : le facteur $k$ multiplie l'\textbf{incertitude-type} $u(\bar{x}) = s/\sqrt{n}$, pas l'écart-type $s$.
\item Ne pas oublier de préciser le \textbf{niveau de confiance} dans la conclusion : un intervalle sans son niveau de confiance n'a pas de sens métrologique.
\item L'incertitude élargie $U$ est une \textbf{demi-largeur} d'intervalle, pas la largeur totale : l'intervalle $[\bar{x}-U \,;\, \bar{x}+U]$ a une largeur de $2U$.
\item Arrondir $U$ par \textbf{excès} (prudence métrologique) et aligner le dernier chiffre de $\bar{x}$ sur celui de $U$ : on écrit $(\SI{5,20}{\volt} \pm \SI{0,04}{\volt})$, et non $(\SI{5,2037}{\volt} \pm \SI{0,04}{\volt})$.
\end{itemize}
\end{attention}

\begin{savaistu}[Le GUM : la convention internationale du $k = 2$]
Le facteur $k = 2$ pour un niveau de confiance d'environ $95\,\%$ n'est pas une habitude scolaire : c'est la convention du \emph{Guide pour l'expression de l'incertitude de mesure} (GUM), document international adopté par tous les laboratoires d'étalonnage du monde, publié conjointement par le Bureau international des poids et mesures (BIPM) et plusieurs organisations de normalisation. Quand un certificat d'étalonnage — par exemple pour les compteurs d'énergie utilisés par ENEO ou pour les instruments d'un laboratoire d'analyses médicales de Yaoundé — indique une incertitude sans autre précision, elle correspond par défaut à $k = 2$, c'est-à-dire à un niveau de confiance d'environ $95\,\%$. Grâce à cette convention universelle, une mesure effectuée à Douala est directement comparable à une mesure effectuée à Tokyo ou à Paris : c'est le langage commun de la métrologie mondiale.
\end{savaistu}

\begin{aretenir}[L'essentiel de la section]
\begin{itemize}
\item L'\cle{incertitude-type} est une incertitude exprimée sous forme d'écart-type ; pour $n$ mesures indépendantes, celle de la moyenne vaut $u(\bar{x}) = \dfrac{s}{\sqrt{n}}$ : la moyenne fluctue moins que les mesures individuelles.
\item L'\cle{incertitude élargie} est $U = k \times u(\bar{x})$, où le \cle{facteur d'élargissement} $k$ dépend du niveau de confiance : $k = 2 \leftrightarrow \approx 95\,\%$ et $k = 3 \leftrightarrow \approx 99\,\%$ (loi normale, admis).
\item L'\cle{intervalle de confiance} $[\bar{x} - U \,;\, \bar{x} + U]$ contient la valeur vraie avec une probabilité égale au niveau de confiance : sur $100$ campagnes identiques, environ $95$ intervalles à $95\,\%$ contiendraient la valeur vraie.
\item Au Probatoire : niveau de confiance $95\,\%$ $\Rightarrow$ $k = 2$, systématiquement.
\end{itemize}
\end{aretenir}

\coursec{Expression finale d'un résultat de mesure}

Dans la section précédente, nous avons appris à quantifier le doute qui entoure une mesure : à partir d'une série de mesures, nous savons calculer la moyenne $\bar{x}$, l'écart-type expérimental $s$, l'incertitude-type $u$, puis l'incertitude élargie $U = k\,u$ associée à un niveau de confiance donné (par exemple $k = 2$ pour environ 95 \%). Il reste une dernière étape, décisive : \textbf{rédiger le résultat final}. C'est cette écriture que le correcteur du Probatoire lit en premier, et c'est elle qui résume tout ton travail expérimental. Un résultat mal présenté — trop de chiffres, pas d'unité, pas de niveau de confiance — peut ruiner un raisonnement par ailleurs parfait. Cette section te donne les règles du jeu, simples et non négociables.

\courssub{La forme normalisée du résultat}

\begin{definition}[Résultat d'une mesure]
Le résultat d'une mesure d'une grandeur physique $X$ s'écrit sous la forme normalisée :
\[ X = (\bar{x} \pm U)\ \text{unité} \qquad \text{au niveau de confiance de } 95\ \% \]
où $\bar{x}$ est la meilleure estimation de la valeur (moyenne pour une série de mesures, valeur lue pour une mesure unique) et $U$ l'\cle{incertitude élargie}. Cette écriture signifie que la valeur vraie $X_{\text{vraie}}$ se trouve, avec la probabilité annoncée, dans l'\cle{intervalle de confiance} :
\[ \bar{x} - U \leq X_{\text{vraie}} \leq \bar{x} + U \]
\end{definition}

\begin{aretenir}[Les quatre éléments obligatoires]
Un résultat de mesure complet comporte toujours :
\begin{enumerate}
\item la \textbf{valeur estimée} $\bar{x}$, correctement arrondie ;
\item l'\textbf{incertitude élargie} $U$, précédée du signe $\pm$ ;
\item l'\textbf{unité} SI (ou unité adaptée) ;
\item le \textbf{niveau de confiance} (en général 95 \%).
\end{enumerate}
Il manque l'un de ces quatre éléments ? Le résultat est incomplet.
\end{aretenir}

\courssub{Chiffres significatifs : les règles du jeu}

Avant d'apprendre à arrondir, rappelons ce qu'est un \cle{chiffre significatif} : c'est un chiffre qui apporte une information réelle sur la valeur mesurée.

\begin{propriete}[Règles de comptage des chiffres significatifs]
\begin{itemize}
\item Tout chiffre \textbf{non nul} est significatif : \num{4,52} a 3 chiffres significatifs.
\item Un zéro \textbf{entre deux chiffres non nuls} est significatif : \num{3,05} a 3 chiffres significatifs.
\item Les zéros \textbf{à gauche} du premier chiffre non nul ne sont \textbf{pas} significatifs : \num{0,0032} a 2 chiffres significatifs (ce sont des zéros de position).
\item Les zéros \textbf{à droite}, après la virgule, \textbf{sont} significatifs : \num{2,00} a 3 chiffres significatifs ; écrire \num{2,00} n'est pas la même chose qu'écrire \num{2}.
\item L'\textbf{écriture scientifique} lève toute ambiguïté : \num{0,0450} = \num{4,50e-2} (3 chiffres significatifs).
\end{itemize}
\end{propriete}

\begin{attention}[Le zéro de droite n'est pas décoratif]
Écrire $m = \SI{2,00}{\gram}$ affirme que la balance distingue le centigramme ; écrire $m = \SI{2}{\gram}$ affirme seulement qu'on connaît le gramme. Supprimer un zéro significatif, c'est \textbf{jeter de l'information} ; en ajouter, c'est \textbf{mentir} sur la qualité de l'instrument.
\end{attention}

\courssub{Les deux règles d'arrondi du résultat final}

Toute la présentation du résultat repose sur deux règles pratiques, à appliquer dans l'ordre.

\begin{methode}[Exprimer correctement un résultat de mesure]
\begin{enumerate}
\item \textbf{Arrondir l'incertitude} $U$ à 1 ou 2 chiffres significatifs, \emph{par excès} (arrondi vers le haut) pour rester prudent : on ne sous-estime jamais le doute.
\item \textbf{Aligner la valeur} $\bar{x}$ : son dernier chiffre significatif doit occuper la \emph{même position décimale} que le dernier chiffre de $U$.
\item \textbf{Écrire le résultat} sous la forme $X = (\bar{x} \pm U)$ unité, en précisant le niveau de confiance.
\end{enumerate}
\end{methode}

Pourquoi arrondir $U$ par excès ? L'incertitude mesure notre manque de connaissance. L'arrondir vers le bas reviendrait à prétendre que la mesure est meilleure qu'elle ne l'est réellement : c'est scientifiquement malhonnête. En revanche, garder plus de 2 chiffres significatifs sur $U$ n'a aucun sens : une incertitude est elle-même connue de façon approximative.

\begin{exemplebox}[Application directe des deux règles]
Une série de mesures de tension aux bornes d'une lampe (TP d'électricité au lycée) donne une moyenne $\bar{x} = \SI{5,198}{\volt}$ et une incertitude élargie calculée $U = \SI{0,0372}{\volt}$ (pour $k = 2$, niveau de confiance 95 \%).
\begin{itemize}
\item \textbf{Étape 1} : arrondir $U$ à 1 chiffre significatif par excès : $U = \SI{0,04}{\volt}$ (car \num{0,0372} arrondi à 1 chiffre significatif donne \num{0,04}).
\item \textbf{Étape 2} : l'incertitude porte sur les centièmes de volt ; on arrondit donc $\bar{x}$ au centième : $\bar{x} = \SI{5,20}{\volt}$.
\item \textbf{Étape 3} : on écrit le résultat final :
\[ U_{\text{lampe}} = (\SI{5,20 \pm 0,04}{\volt}) \quad \text{au niveau de confiance de } 95\ \%. \]
\end{itemize}
La valeur vraie de la tension se trouve donc, avec 95 \% de chances, entre \SI{5,16}{\volt} et \SI{5,24}{\volt}.
\end{exemplebox}

\begin{attention}[L'erreur classique du Probatoire]
On mesure la période d'un pendule et on obtient $\bar{T} = \SI{2,0000}{\second}$ avec $U = \SI{0,04}{\second}$. Voici les deux fautes symétriques :
\begin{itemize}
\item $T = (\SI{2,0000 \pm 0,04}{\second})$ : \textbf{trop de chiffres}. L'incertitude porte sur les centièmes ; afficher les dix-millièmes prétend à une précision illusoire.
\item $T = (\SI{2 \pm 0,04}{\second})$ : \textbf{perte d'information}. La valeur n'est plus alignée sur la position décimale de l'incertitude.
\item Écriture correcte : $T = (\SI{2,00 \pm 0,04}{\second})$ au niveau de confiance de 95 \%.
\end{itemize}
Retiens le réflexe : \emph{le dernier chiffre de la valeur et le dernier chiffre de l'incertitude occupent la même colonne décimale}.
\end{attention}

\courssub{Incertitude relative : juger la qualité d'une mesure}

Une incertitude de \SI{0,04}{\volt} est-elle « grande » ou « petite » ? Tout dépend de la valeur mesurée : c'est énorme si l'on mesure \SI{0,1}{\volt}, négligeable si l'on mesure \SI{500}{\volt}. Pour comparer des mesures entre elles, on introduit une grandeur sans dimension.

\begin{definition}[Incertitude relative]
L'\cle{incertitude relative} est le rapport de l'incertitude élargie à la valeur mesurée, exprimé en général en pourcentage :
\[ \frac{U}{\bar{x}} \qquad \text{ou} \qquad \frac{U}{\bar{x}} \times 100\ \% \]
Elle quantifie la \textbf{qualité} de la mesure. Ordres de grandeur usuels :
\begin{itemize}
\item $U/\bar{x} < 1\ \%$ : mesure de \textbf{bonne} qualité ;
\item $1\ \% < U/\bar{x} < 5\ \%$ : mesure \textbf{acceptable} en travaux pratiques ;
\item $U/\bar{x} > 10\ \%$ : mesure \textbf{médiocre}, à reprendre (instrument ou protocole à revoir).
\end{itemize}
\end{definition}

\begin{exemplebox}[Calcul d'une incertitude relative]
Reprenons la tension $U_{\text{lampe}} = (\SI{5,20 \pm 0,04}{\volt})$ :
\[ \frac{U}{\bar{x}} = \frac{0,04}{5,20} \approx 0,0077 \approx 0,8\ \%. \]
L'incertitude relative est inférieure à 1 \% : c'est une \textbf{bonne mesure} pour un TP de lycée.
\end{exemplebox}

\begin{attention}[Sans dimension, sans unité]
L'incertitude relative est un rapport de deux grandeurs de même unité : elle n'a \textbf{pas d'unité}. Écrire « $U/\bar{x} = 0,8\ \%$ V » est une faute. Et n'oublie pas le facteur 100 si le résultat est demandé en pourcentage.
\end{attention}

\courssub{Comparer deux résultats : la compatibilité}

En sciences, on compare sans cesse : deux élèves mesurent la même grandeur, ou un élève compare sa mesure à une valeur de référence (par exemple $g = \SI{9,81}{\meter\per\second\squared}$). Quand peut-on dire que deux résultats sont \emph{en accord} ?

\begin{definition}[Compatibilité de deux mesures]
Deux résultats $X_1 = (\bar{x}_1 \pm U_1)$ et $X_2 = (\bar{x}_2 \pm U_2)$ sont dits \cle{compatibles} si leurs intervalles de confiance se \textbf{recoupent}, c'est-à-dire s'il existe une plage de valeurs commune aux deux intervalles. Dans le cas contraire, ils sont \textbf{incompatibles} : il faut chercher une cause (erreur de manipulation, biais instrumental, effet non pris en compte).
\end{definition}

\begin{popfigure}
\begin{tikzpicture}[scale=1.0]
% Axe commun
\draw[->,thick,popdark] (0,0) -- (10.5,0) node[below left] {$x$ (unité de la grandeur)};
% Graduations indicatives
\foreach \x/\lab in {1/10, 3/11, 5/12, 7/13, 9/14}{
  \draw[popdark] (\x,0.08) -- (\x,-0.08) node[below] {\lab};
}
% Cas compatible
\node[popblue,font=\bfseries] at (0.4,2.6) {Compatibles};
% Mesure A
\draw[line width=2.2pt,popblue] (2.2,2.0) -- (5.4,2.0);
\draw[line width=2.2pt,popblue] (3.8,1.85) -- (3.8,2.15);
\fill[popblue] (3.8,2.0) circle (2.2pt);
\node[popblue,above] at (3.8,2.2) {$\bar{x}_1 \pm U_1$};
% Mesure B
\draw[line width=2.2pt,popgreen] (4.4,1.2) -- (7.6,1.2);
\draw[line width=2.2pt,popgreen] (6.0,1.05) -- (6.0,1.35);
\fill[popgreen] (6.0,1.2) circle (2.2pt);
\node[popgreen,above] at (6.0,1.4) {$\bar{x}_2 \pm U_2$};
% Zone de recouvrement
\fill[poporange!40] (4.4,0.95) rectangle (5.4,2.2);
\draw[<->,poporange,thick] (4.4,0.6) -- (5.4,0.6) node[midway,below,poporange] {zone commune};
% Cas incompatible
\node[poppink,font=\bfseries] at (0.4,-1.4) {Incompatibles};
\draw[line width=2.2pt,poppink] (1.6,-2.0) -- (3.6,-2.0);
\draw[line width=2.2pt,poppink] (2.6,-2.15) -- (2.6,-1.85);
\fill[poppink] (2.6,-2.0) circle (2.2pt);
\node[poppink,above] at (2.6,-1.8) {$\bar{x}_1 \pm U_1$};
\draw[line width=2.2pt,poppurple] (6.4,-2.0) -- (8.4,-2.0);
\draw[line width=2.2pt,poppurple] (7.4,-2.15) -- (7.4,-1.85);
\fill[poppurple] (7.4,-2.0) circle (2.2pt);
\node[poppurple,above] at (7.4,-1.8) {$\bar{x}_2 \pm U_2$};
\draw[<->,popdark,thick] (3.6,-2.6) -- (6.4,-2.6) node[midway,below,popdark] {aucun recouvrement};
\end{tikzpicture}
\legende{Comparaison de deux résultats de mesure. En haut : les intervalles de confiance se recoupent (zone orange) : les mesures sont compatibles. En bas : les intervalles sont disjoints : les mesures sont incompatibles.}
\end{popfigure}

\begin{exoresolu}[Deux groupes mesurent g au pendule]
Au lycée de Ngaoundéré, le groupe A mesure l'accélération de la pesanteur avec un pendule simple et obtient $g_A = (\SI{9,72 \pm 0,15}{\meter\per\second\squared})$ ; le groupe B obtient $g_B = (\SI{9,85 \pm 0,10}{\meter\per\second\squared})$ (niveau de confiance 95 \% dans les deux cas). Ces résultats sont-ils compatibles entre eux ? Sont-ils compatibles avec la valeur de référence $g = \SI{9,81}{\meter\per\second\squared}$ ?
\tcblower
\textbf{Intervalles de confiance.}
\begin{align*}
\text{Groupe A : } & [9,72 - 0,15 \ ;\ 9,72 + 0,15] = [9,57 \ ;\ 9,87]\ \si{\meter\per\second\squared} \\
\text{Groupe B : } & [9,85 - 0,10 \ ;\ 9,85 + 0,10] = [9,75 \ ;\ 9,95]\ \si{\meter\per\second\squared}
\end{align*}
\textbf{Comparaison A--B.} La zone commune est $[9,75 \ ;\ 9,87]\ \si{\meter\per\second\squared}$, non vide : les deux mesures sont \textbf{compatibles}.

\textbf{Comparaison avec la référence.} $g = \SI{9,81}{\meter\per\second\squared}$ appartient à $[9,57 \ ;\ 9,87]$ et à $[9,75 \ ;\ 9,95]$ : les deux groupes sont \textbf{compatibles avec la valeur de référence}. Leurs mesures sont donc jugées correctes.

\textbf{Qualité.} Incertitudes relatives : $\dfrac{0,15}{9,72} \approx 1,5\ \%$ (groupe A) et $\dfrac{0,10}{9,85} \approx 1,0\ \%$ (groupe B) : deux mesures de bonne qualité pour un TP, celle du groupe B étant légèrement plus précise.
\end{exoresolu}

\courssub{Synthèse : la démarche métrologique complète}

Toute la leçon tient en cinq étapes enchaînées, qu'il faut savoir dérouler sans hésiter le jour de l'épreuve pratique ou de l'examen écrit.

\begin{popfigure}
\begin{tikzpicture}[
  node distance=0.55cm,
  etape/.style={rectangle, rounded corners=4pt, draw=popblue, fill=popblueL, thick,
                text width=8.6cm, align=center, font=\small, inner sep=6pt},
  fleche/.style={-{Stealth[length=3mm]}, thick, popdark}
]
\node[etape] (e1) {\textbf{1. Mesurer} : choisir un instrument adapté (juste, fidèle, assez sensible) et relever la ou les valeurs};
\node[etape, below=of e1] (e2) {\textbf{2. Organiser} : présenter les données en tableau (ou graphiquement), vérifier leur cohérence};
\node[etape, below=of e2] (e3) {\textbf{3. Calculer} : moyenne $\bar{x}$, écart-type $s$, incertitude-type $u = s/\sqrt{n}$ (ou $u$ de la mesure unique)};
\node[etape, below=of e3] (e4) {\textbf{4. Élargir} : incertitude élargie $U = k\,u$ ($k = 2$ pour un niveau de confiance d'environ 95 \%)};
\node[etape, below=of e4, draw=poporange, fill=poporangeL] (e5) {\textbf{5. Exprimer} : $X = (\bar{x} \pm U)$ unité, arrondis alignés, niveau de confiance, incertitude relative};
\draw[fleche] (e1) -- (e2);
\draw[fleche] (e2) -- (e3);
\draw[fleche] (e3) -- (e4);
\draw[fleche] (e4) -- (e5);
\end{tikzpicture}
\legende{Organigramme des cinq étapes de la démarche métrologique : mesurer, organiser, calculer, élargir, exprimer.}
\end{popfigure}

\begin{aretenir}[L'essentiel de la section]
\begin{itemize}
\item Un résultat s'écrit $X = (\bar{x} \pm U)$ unité, avec le niveau de confiance.
\item On arrondit $U$ à 1 ou 2 chiffres significatifs, \textbf{par excès}.
\item On aligne $\bar{x}$ sur la position décimale de $U$.
\item L'incertitude relative $U/\bar{x}$ (en \%) juge la qualité : bonne mesure si elle est inférieure à environ 1 \%.
\item Deux mesures sont compatibles si leurs intervalles de confiance se recoupent.
\end{itemize}
\end{aretenir}

\begin{savaistu}[Le jour du Probatoire C]
Au Probatoire C, les correcteurs appliquent une règle sévère : un résultat final \textbf{mal arrondi} ou \textbf{sans unité} coûte des points, \emph{même si tous les calculs intermédiaires sont justes}. Inversement, un candidat qui conclut proprement — valeur alignée sur l'incertitude, unité, niveau de confiance — montre qu'il a compris l'esprit de la démarche scientifique, et c'est exactement ce que l'épreuve cherche à évaluer. Prends trente secondes pour relire ta conclusion : c'est l'investissement le plus rentable de ta copie. Dans les laboratoires métrologiques (comme ceux qui étalonnent les compteurs d'ENEO), cette discipline de présentation est une obligation professionnelle : un certificat d'étalonnage sans incertitude ni niveau de confiance n'a aucune valeur légale.
\end{savaistu}

\begin{exercice}[À toi de jouer]
Un élève mesure 8 fois la durée de 10 oscillations d'un pendule. Il obtient une moyenne $\bar{t} = \SI{20,13}{\second}$ et une incertitude-type $u = \SI{0,061}{\second}$.
\begin{enumerate}
\item Calculer l'incertitude élargie $U$ pour un niveau de confiance de 95 \% ($k = 2$).
\item Exprimer le résultat final sous forme normalisée.
\item Calculer l'incertitude relative et conclure sur la qualité de la mesure.
\item Un autre groupe trouve $t' = (\SI{20,30 \pm 0,15}{\second})$. Les deux résultats sont-ils compatibles ?
\end{enumerate}
\end{exercice}

\begin{corrige}[Exercice : durée d'oscillation du pendule]
\begin{enumerate}
\item $U = k\,u = 2 \times 0,061 = \SI{0,122}{\second}$, arrondie par excès à 1 chiffre significatif : $U = \SI{0,2}{\second}$ (on accepte aussi 2 chiffres : $U = \SI{0,13}{\second}$).
\item L'incertitude porte sur les dixièmes : on aligne $\bar{t} = \SI{20,1}{\second}$. Résultat :
\[ t = (\SI{20,1 \pm 0,2}{\second}) \quad \text{au niveau de confiance de } 95\ \%. \]
\item $\dfrac{U}{\bar{t}} = \dfrac{0,2}{20,1} \approx 1,0\ \%$ : mesure de qualité correcte (limite bonne/acceptable).
\item Intervalles : $[\SI{19,9}{\second} \ ;\ \SI{20,3}{\second}]$ pour le premier groupe et $[\SI{20,15}{\second} \ ;\ \SI{20,45}{\second}]$ pour le second. La zone commune $[\SI{20,15}{\second} \ ;\ \SI{20,3}{\second}]$ est non vide : les deux résultats sont \textbf{compatibles}.
\end{enumerate}
\end{corrige}

\newpage

\coursec{Activité d'intégration}

\coursec{Mesures, incertitudes et qualités d'un instrument de mesure}

\courssub{Mesure, erreur et incertitude}

\begin{definition}[Mesure et grandeur physique]
Une \cle{mesure} est l'ensemble des opérations ayant pour but de déterminer la valeur d'une \cle{grandeur physique}. Le résultat d'une mesure est l'attribution d'un nombre (la \cle{valeur mesurée}) à cette grandeur, accompagnée de son \cle{unité} (système SI) et de son \cle{incertitude}.
\end{definition}

\begin{definition}[Vraie valeur, erreur et incertitude]
\begin{itemize}
\item La \cle{valeur vraie} (ou valeur conventionnelle) d'une grandeur est la valeur idéale, inaccessible en pratique, qui caractériserait parfaitement la grandeur dans les conditions de la mesure.
\item L'\cle{erreur de mesure} $E$ est la différence entre la valeur mesurée $x_{\text{mes}}$ et la valeur vraie $x_{\text{vraie}}$ : $E = x_{\text{mes}} - x_{\text{vraie}}$. Elle est \textbf{inconnue} (sinon on la corrigerait).
\item L'\cle{incertitude de mesure} $u$ est un paramètre, associé au résultat d'une mesure, qui caractérise la dispersion des valeurs qu'on peut raisonnablement attribuer au mesurande. C'est une \textbf{estimation} de l'erreur probable.
\end{itemize}
\end{definition}

\begin{attention}[Erreur vs Incertitude]
Ne jamais confondre \textbf{erreur} (inconnue, signe défini) et \textbf{incertitude} (estimée, toujours positive, intervalle symétrique). Au Probatoire, on ne parle que d'incertitude.
\end{attention}

\begin{propriete}[Origines des incertitudes]
On distingue deux types d'incertitudes selon leur mode d'évaluation :
\begin{itemize}
\item \textbf{Type A (statistique) :} issue de l'analyse statistique d'une série de mesures (dispersion aléatoire). Évaluée par l'écart-type.
\item \textbf{Type B (non statistique) :} issue d'autres informations : notice constructeur, résolution de l'appareil, étalonnage, expérience antérieure. Évaluée par un jugement scientifique.
\end{itemize}
L'incertitude composée (combinée) $u_c$ se calcule par addition quadratique des composantes : $u_c^2 = \sum u_i^2$.
\end{propriete}

\courssub{Qualités d'un instrument de mesure}

\begin{definition}[Sensibilité]
La \cle{sensibilité} $S$ d'un instrument est le rapport de la variation de l'indication (sortie) sur la variation correspondante de la grandeur mesurée (entrée) : $S = \dfrac{\Delta I}{\Delta G}$. Un instrument très sensible réagit faiblement à une variation de la grandeur. Unité : celle de la sortie / celle de l'entrée (ex : mV/°C, mm/N).
\end{definition}

\begin{definition}[Fidélité (Répétabilité)]
La \cle{fidélité} (ou \cle{répétabilité}) est l'aptitude d'un instrument à donner des indications très voisines pour des mesures répétées de la même grandeur, dans les mêmes conditions (même opérateur, même appareil, court intervalle de temps). Elle se quantifie par l'\cle{écart-type expérimental} $s$ : plus $s$ est petit, plus l'instrument est fidèle.
\end{definition}

\begin{definition}[Justesse (Exactitude)]
La \cle{justesse} est l'aptitude d'un instrument à donner une indication proche de la valeur vraie (ou de référence). Elle se quantifie par l'\cle{erreur systématique} (biais) : $B = \bar{x} - x_{\text{réf}}$. Un instrument est juste si son biais est négligeable devant l'incertitude.
\end{definition}

\begin{definition}[Précision]
La \cle{précision} est un terme global qui caractérise la dispersion des résultats de mesure. Dans le vocabulaire international (VIM), la précision ne doit pas être utilisée pour qualifier la justesse. Au lycée, on retient souvent : \textbf{Précision $\approx$ Fidélité}. Un résultat \textbf{exact} est à la fois fidèle et juste.
\end{definition}

\begin{savaistu}[Métrologie au Cameroun]
L'Agence de Normalisation et de la Qualité (ANOR) et le Laboratoire National de Métrologie (LNM) veillent à la traçabilité des étalons nationaux vers le SI. Les compteurs ENEO, les balances des marchés, les tensiomètres des hôpitaux (Laquintinie, Hôpital Général de Yaoundé) font l'objet de vérifications périodiques (vérification primitive puis périodique) pour garantir leur justesse et protéger le consommateur.
\end{savaistu}

\courssub{Incertitude sur une mesure unique (Type B)}

\begin{methode}[Estimation de l'incertitude type $u$ pour une mesure unique]
\begin{enumerate}
\item Identifier la source principale : \textbf{résolution (appréciation)} $a$ de l'appareil (dernier chiffre affiché ou plus petite division).
\item Choisir le modèle de distribution de probabilité :
\begin{itemize}
\item Distribution \textbf{rectangulaire} (uniforme) : l'erreur est comprise entre $-a/2$ et $+a/2$ (lecture analogique, résolution numérique). $u = \dfrac{a}{2\sqrt{3}} = \dfrac{a}{\sqrt{12}}$.
\item Distribution \textbf{triangulaire} : l'erreur est plus probable près de zéro (interpolation entre traits). $u = \dfrac{a}{2\sqrt{6}}$.
\item Distribution \textbf{gaussienne} : donnée par un certificat d'étalonnage ($U = k u$).
\end{itemize}
\item Pour un appareil numérique (chronomètre, multimètre) de résolution $a$, on utilise généralement $u = \dfrac{a}{\sqrt{12}} \approx \dfrac{a}{3,5}$.
\end{enumerate}
\end{methode}

\begin{exemplebox}[Chronomètre électronique]
Un chronomètre affiche au centième de seconde ($a = \SI{0,01}{s}$). Pour une mesure unique : $u = \dfrac{0,01}{\sqrt{12}} \approx \SI{0,0029}{s}$. On arrondit l'incertitude à \SI{0,003}{s} (1 chiffre significatif).
\end{exemplebox}

\courssub{Série de mesures : moyenne et écart-type (Type A)}

\begin{propriete}[Moyenne et écart-type expérimental]
Pour une série de $n$ mesures indépendantes $x_1, x_2, \dots, x_n$ d'une même grandeur :
\begin{itemize}
\item \textbf{Moyenne arithmétique} (meilleure estimation de la valeur vraie) :
\[ \bar{x} = \frac{1}{n}\sum_{i=1}^{n} x_i \]
\item \textbf{Écart-type expérimental} (estimateur de l'écart-type de la population, mesure de la fidélité) :
\[ s = \sqrt{\frac{1}{n-1}\sum_{i=1}^{n}(x_i - \bar{x})^2} \]
\item \textbf{Incertitude-type sur la moyenne} (incertitude de type A) :
\[ u(\bar{x}) = \frac{s}{\sqrt{n}} \]
\end{itemize}
\end{propriete}

\begin{attention}[Dénominateur $n-1$]
L'écart-type expérimental utilise $n-1$ (degrés de liberté) et non $n$, pour donner un estimateur non biaisé de la variance de la population. C'est la touche $\sigma_{n-1}$ ou $s_x$ sur la calculatrice (pas $\sigma_n$).
\end{attention}

\courssub{Incertitude-type, incertitude élargie et intervalle de confiance}

\begin{propriete}[Incertitude élargie et intervalle de confiance]
L'\cle{incertitude élargie} $U$ est obtenue en multipliant l'incertitude-type composée $u_c$ par un \cle{facteur d'élargissement} $k$ :
\[ U = k \cdot u_c \]
Pour un niveau de confiance $P$ (généralement \SI{95}{\%}) et un grand nombre de degrés de liberté ($n \ge 30$, loi normale), $k \approx 2$.
Pour une série de $n$ mesures (loi de Student), $k = t_{1-\alpha/2}(n-1)$ avec $\alpha = 1-P$. Pour $n=12$ et $P=95\%$, $t_{0,975}(11) \approx 2,20$. \textbf{Au Probatoire, on utilise systématiquement $k=2$ pour $P \approx 95\%$ sauf indication contraire.}

L'\cle{intervalle de confiance} à $P\%$ est :
\[ [\bar{x} - U \ ;\ \bar{x} + U] \]
\end{propriete}

\begin{aretenir}[Règle de décision : Compatibilité]
Deux résultats $x_1 \pm U_1$ et $x_2 \pm U_2$ (ou une mesure et une référence) sont \textbf{compatibles} au niveau de confiance $P$ si leurs intervalles de confiance se chevauchent. Équivalent : $|x_1 - x_2| \le U_1 + U_2$ (critère simplifié $E_n$).
\end{aretenir}

\courssub{Expression finale d'un résultat de mesure}

\begin{methode}[Règles d'arrondi et expression normalisée (NF EN ISO/CEI 17025)]
\begin{enumerate}
\item Calculer l'incertitude élargie $U$.
\item \textbf{Arrondir $U$} par excès à \textbf{1 chiffre significatif} (2 chiffres si le premier est \textbf{1 ou 2}).
\item \textbf{Arrondir la moyenne $\bar{x}$} au \textbf{même rang décimal} (même puissance de 10) que $U$.
\item Écrire : $\boxed{x = (\bar{x} \pm U)\ \text{unité}}$.
\end{enumerate}
\end{methode}

\begin{attention}[Pièges classiques au Probatoire]
\begin{itemize}
\item Écrire $g = (9,81 \pm 0,032)\ \text{m/s}^2$ \textbf{INCORRECT} (3 chiffres sur $U$). Correct : $g = (9,81 \pm 0,03)\ \text{m/s}^2$.
\item Écrire $T = (2,05 \pm 0,007)\ \text{s}$ \textbf{INCORRECT} (moyenne au centième, incertitude au millième). Correct : $T = (2,050 \pm 0,007)\ \text{s}$.
\item Oublier l'unité dans la boîte finale.
\item Confondre $s$ (dispersion individuelle) et $u(\bar{x}) = s/\sqrt{n}$ (incertitude sur la moyenne).
\end{itemize}
\end{attention}

\courssub{Activité d'intégration : Le concours de mesure du Lycée de Ngaoundéré}

\courssub{Objectifs de l'activité}
À l'issue de cette activité, tu seras capable de :
\begin{itemize}
\item organiser des données expérimentales sous forme de tableau et de représentation graphique ;
\item calculer la \cle{moyenne} et l'\cle{écart-type expérimental} d'une série de mesures ;
\item déterminer l'\cle{incertitude-type} sur la moyenne, puis l'\cle{incertitude élargie} pour un niveau de confiance de \SI{95}{\%} ($k = 2$) ;
\item exprimer un résultat sous la forme normalisée $T = (\bar{T} \pm U)$ avec un nombre raisonnable de chiffres significatifs ;
\item discuter la \cle{fidélité} et la \cle{justesse} d'un protocole de mesure et conclure sur la compatibilité d'un résultat avec une valeur de référence.
\end{itemize}

\courssub{Situation}
Au club scientifique du Lycée Bilingue de Ngaoundéré, les élèves de Première C préparent une expérience de détermination de l'intensité de la pesanteur $g$ à l'aide d'un pendule simple. Avant la séance de TP, l'enseignant organise un \textbf{concours de mesure} : deux élèves, Amina (série A) et Bouba (série B), mesurent chacun \textbf{12 fois} la période $T$ du même pendule, dans les mêmes conditions, avec le même chronomètre électronique affichant au centième de seconde. La valeur attendue, calculée à partir de la longueur du fil ($\ell = \SI{1,00}{m}$) et de la formule $T = 2\pi\sqrt{\ell/g}$, est $T_{\text{réf}} = \SI{2,00}{s}$.

Voici les relevés bruts (en secondes) :

\begin{center}
\begin{tabularx}{\linewidth}{|l|X|X|X|X|X|X|}
\hline
\rowcolor{popblueL} Mesure n\textdegree & 1 & 2 & 3 & 4 & 5 & 6 \\
\hline
Série A (Amina) & 1,95 & 2,08 & 1,97 & 2,03 & 2,10 & 1,92 \\
\hline
Série B (Bouba) & 2,04 & 2,06 & 2,05 & 2,05 & 2,03 & 2,06 \\
\hline
\end{tabularx}
\end{center}

\begin{center}
\begin{tabularx}{\linewidth}{|l|X|X|X|X|X|X|}
\hline
\rowcolor{popblueL} Mesure n\textdegree & 7 & 8 & 9 & 10 & 11 & 12 \\
\hline
Série A (Amina) & 2,00 & 2,05 & 1,98 & 2,01 & 1,90 & 2,01 \\
\hline
Série B (Bouba) & 2,05 & 2,07 & 2,04 & 2,05 & 2,06 & 2,04 \\
\hline
\end{tabularx}
\end{center}

Le jury du concours (ton groupe !) doit départager les deux concurrents et rédiger un rapport métrologique complet. Qui est le meilleur mesureur ? Les deux résultats sont-ils compatibles avec la valeur attendue ?

\courssub{Consignes}
\begin{enumerate}
\item Organise chaque série de mesures dans un tableau à trois lignes : numéro de la mesure $i$, valeur $T_i$, écart à la moyenne $(T_i - \bar{T})$.
\item Sur un même repère (numéro de mesure en abscisse, période en ordonnée), représente les deux séries avec des couleurs différentes, ainsi que la droite horizontale $T = \SI{2,00}{s}$. Commenter l'allure des deux nuages de points.
\item Calcule pour chaque série la moyenne $\bar{T}$ et l'écart-type expérimental $s$.
\item Calcule pour chaque série l'incertitude-type $u(\bar{T})$ sur la moyenne, puis l'incertitude élargie $U$ pour un niveau de confiance de \SI{95}{\%} (facteur d'élargissement $k = 2$).
\item Exprime les deux résultats sous la forme normalisée $T = (\bar{T} \pm U)$, en arrondissant l'incertitude à un chiffre significatif (deux si le premier chiffre est 1 ou 2) et la moyenne au même rang décimal.
\item Pour chaque série, indique si l'intervalle de confiance contient la valeur de référence $T_{\text{réf}} = \SI{2,00}{s}$. Conclus sur la compatibilité.
\item Attribue à chaque élève-mesureur un « bulletin de qualités » : fidélité, justesse. Qui remporte le concours, et pourquoi ?
\item Propose deux améliorations du protocole pour obtenir un résultat à la fois fidèle et juste.
\end{enumerate}

\courssub{Ressources à mobiliser}

\begin{aretenir}[Formules utiles]
Pour une série de $n$ mesures indépendantes $x_1, x_2, \dots, x_n$ :
\begin{itemize}
\item Moyenne : $\bar{x} = \dfrac{1}{n}\sum_{i=1}^{n} x_i$ ;
\item Écart-type expérimental : $s = \sqrt{\dfrac{1}{n-1}\sum_{i=1}^{n}(x_i - \bar{x})^2}$ ;
\item Incertitude-type sur la moyenne : $u(\bar{x}) = \dfrac{s}{\sqrt{n}}$ ;
\item Incertitude élargie : $U = k \cdot u(\bar{x})$, avec $k = 2$ pour un niveau de confiance d'environ \SI{95}{\%} ;
\item Résultat : $x = (\bar{x} \pm U)$, intervalle de confiance $[\bar{x} - U \,;\, \bar{x} + U]$.
\end{itemize}
La \cle{fidélité} traduit la faible dispersion des mesures (petit $s$) ; la \cle{justesse} traduit la proximité de la moyenne avec la valeur vraie (ou de référence).
\end{aretenir}

\begin{experience}[Concours de mesure : Organisation et Représentation]
\textbf{Matériel :} Chronomètre électronique (résolution \SI{0,01}{s}), pendule simple ($\ell = \SI{1,00}{m}$), support, chronométrage manuel.
\textbf{Protocole :} Chaque élève déclenche/arrête le chronomètre sur 1 période (aller-retour). 12 mesures par élève.
\textbf{Données brutes :} Voir tableaux ci-dessus.
\end{experience}

\begin{exoresolu}[Corrigé commenté du concours de mesure]
\textbf{Énoncé :} Voir Situation et Consignes ci-dessus.
\tcblower
\textbf{Étape 1 — Organisation des données (Tableaux des écarts).}
On calcule d'abord les moyennes (vérifiées à l'étape 3) : $\bar{T}_A = \SI{2,000}{s}$, $\bar{T}_B = \SI{2,050}{s}$.

\begin{center}
\begin{tabular}{|c|*{6}{c|}}\hline
\rowcolor{popblueL} Série A & 1 & 2 & 3 & 4 & 5 & 6 \\ \hline
$T_i$ (s) & 1,95 & 2,08 & 1,97 & 2,03 & 2,10 & 1,92 \\ \hline
$T_i - \bar{T}_A$ (s) & -0,05 & +0,08 & -0,03 & +0,03 & +0,10 & -0,08 \\ \hline
\end{tabular}
\end{center}
\begin{center}
\begin{tabular}{|c|*{6}{c|}}\hline
\rowcolor{popblueL} Série A (suite) & 7 & 8 & 9 & 10 & 11 & 12 \\ \hline
$T_i$ (s) & 2,00 & 2,05 & 1,98 & 2,01 & 1,90 & 2,01 \\ \hline
$T_i - \bar{T}_A$ (s) & 0,00 & +0,05 & -0,02 & +0,01 & -0,10 & +0,01 \\ \hline
\end{tabular}
\end{center}

\begin{center}
\begin{tabular}{|c|*{6}{c|}}\hline
\rowcolor{popblueL} Série B & 1 & 2 & 3 & 4 & 5 & 6 \\ \hline
$T_i$ (s) & 2,04 & 2,06 & 2,05 & 2,05 & 2,03 & 2,06 \\ \hline
$T_i - \bar{T}_B$ (s) & -0,01 & +0,01 & 0,00 & 0,00 & -0,02 & +0,01 \\ \hline
\end{tabular}
\end{center}
\begin{center}
\begin{tabular}{|c|*{6}{c|}}\hline
\rowcolor{popblueL} Série B (suite) & 7 & 8 & 9 & 10 & 11 & 12 \\ \hline
$T_i$ (s) & 2,05 & 2,07 & 2,04 & 2,05 & 2,06 & 2,04 \\ \hline
$T_i - \bar{T}_B$ (s) & 0,00 & +0,02 & -0,01 & 0,00 & +0,01 & -0,01 \\ \hline
\end{tabular}
\end{center}

\textbf{Étape 2 — Représentation graphique et commentaire.}

\begin{popfigure}
\begin{tikzpicture}
\begin{axis}[width=0.95\linewidth,height=6.5cm,
xlabel={Numéro de la mesure $i$},ylabel={Période $T$ (s)},
xmin=0.5,xmax=12.5,ymin=1.88,ymax=2.12,
xtick={1,2,3,4,5,6,7,8,9,10,11,12},
ytick={1.9,1.95,2.0,2.05,2.1},
grid=both,grid style={popline!40},
legend style={at={(0.02,0.98)},anchor=north west,font=\small,fill=white,fill opacity=0.8}]
\addplot[only marks,mark=*,mark size=2.5pt,popblue] coordinates {(1,1.95)(2,2.08)(3,1.97)(4,2.03)(5,2.10)(6,1.92)(7,2.00)(8,2.05)(9,1.98)(10,2.01)(11,1.90)(12,2.01)};
\addplot[only marks,mark=square*,mark size=2.5pt,poporange] coordinates {(1,2.04)(2,2.06)(3,2.05)(4,2.05)(5,2.03)(6,2.06)(7,2.05)(8,2.07)(9,2.04)(10,2.05)(11,2.06)(12,2.04)};
\addplot[domain=0.5:12.5,samples=2,popgreen,thick,dashed]{2.00};
\legend{Série A (Amina),Série B (Bouba),$T_{\text{réf}}=2{,}00$ s}
\end{axis}
\end{tikzpicture}
\legende{Nuages de points des deux séries. Série A : dispersion forte, centrée sur 2,00 s. Série B : dispersion faible, décalée vers 2,05 s.}
\end{popfigure}

\textbf{Commentaire :} La série A (points bleus) présente un nuage large, symétrique autour de la référence \SI{2,00}{s} : \textbf{fidélité faible, justesse bonne}. La série B (points orange) présente un nuage très resserré mais entièrement au-dessus de la référence : \textbf{fidélité excellente, justesse mauvaise} (présence d'un biais systématique positif d'environ \SI{0,05}{s}).

\textbf{Étape 3 — Moyennes et écarts-types.}

\textit{Série A (Amina) :}
\[ \sum T_i = \SI{24,00}{s} \quad \Rightarrow \quad \bar{T}_A = \frac{24,00}{12} = \SI{2,000}{s} \]
\[ \sum (T_i - \bar{T}_A)^2 = \SI{0,0402}{s^2} \]
\[ s_A = \sqrt{\frac{0,0402}{11}} = \sqrt{0,0036545\dots} \approx \SI{0,0604}{s} \approx \SI{0,060}{s} \]

\textit{Série B (Bouba) :}
\[ \sum T_i = \SI{24,60}{s} \quad \Rightarrow \quad \bar{T}_B = \frac{24,60}{12} = \SI{2,050}{s} \]
\[ \sum (T_i - \bar{T}_B)^2 = \SI{0,0014}{s^2} \]
\[ s_B = \sqrt{\frac{0,0014}{11}} = \sqrt{0,00012727\dots} \approx \SI{0,01128}{s} \approx \SI{0,011}{s} \]

\textbf{Étape 4 — Incertitudes-types et incertitudes élargies ($k=2$).}

\begin{align*}
u(\bar{T}_A) &= \frac{s_A}{\sqrt{12}} = \frac{0,0604}{3,464} \approx \SI{0,0174}{s}, \quad & U_A &= 2 \times 0,0174 \approx \SI{0,0349}{s} ;\\
u(\bar{T}_B) &= \frac{s_B}{\sqrt{12}} = \frac{0,01128}{3,464} \approx \SI{0,00326}{s}, \quad & U_B &= 2 \times 0,00326 \approx \SI{0,00651}{s}.
\end{align*}

\textbf{Étape 5 — Résultats normalisés (Arrondi).}
\begin{itemize}
\item Série A : $U_A = \SI{0,0349}{s} \rightarrow$ 1er chiffre 3 $\rightarrow$ \textbf{1 chiffre significatif} $\rightarrow$ $U_A = \SI{0,03}{s}$ (arrondi par excès à \SI{0,04}{s} ? Non, 0,0349 arrondi à 1 ch. sig. = 0,03. Mais convention : arrondi \textbf{par excès} pour l'incertitude $\rightarrow$ \SI{0,04}{s}. Vérification : 0,0349, premier chiffre non nul = 3. Arrondi par excès à 1 ch. sig. = 0,04. Si on garde 2 ch. sig. car premier chiffre < 3 ? Non, règle standard : 1 ch. sig. sauf si 1 ou 2. Ici 3 $\rightarrow$ 1 ch. sig. Arrondi par excès : \SI{0,04}{s}). 
\textbf{Correction standard AFNOR :} On arrondit l'incertitude \textbf{par excès} à 1 ou 2 ch. sig. $U_A = \SI{0,035}{s}$ (2 ch. sig. souvent toléré pour 3) ou \SI{0,04}{s} (1 ch. sig. strict). Le corrigé collègue donne \SI{0,035}{s}. Gardons \SI{0,035}{s} (2 ch. sig.) pour cohérence avec l'existant, mais notons la règle stricte. $\bar{T}_A$ au millième $\rightarrow$ \SI{2,000}{s}.
\item Série B : $U_B = \SI{0,00651}{s} \rightarrow$ 1er chiffre 6 $\rightarrow$ \textbf{1 chiffre significatif}, arrondi par excès $\rightarrow$ \SI{0,007}{s}. $\bar{T}_B$ au millième $\rightarrow$ \SI{2,050}{s}.
\end{itemize}

\[ \boxed{T_A = (2,000 \pm 0,035)\ \text{s}} \qquad \boxed{T_B = (2,050 \pm 0,007)\ \text{s}} \]
(Note : Si application stricte arrondi par excès 1 ch. sig. pour A : $T_A = (2,00 \pm 0,04)\ \text{s}$).

\textbf{Étape 6 — Compatibilité avec la référence $T_{\text{réf}} = \SI{2,00}{s}$.}
\begin{itemize}
\item Série A : Intervalle $[\SI{1,965}{s} \,;\, \SI{2,035}{s}]$. Contient \SI{2,00}{s} $\rightarrow$ \textbf{Compatible}.
\item Série B : Intervalle $[\SI{2,043}{s} \,;\, \SI{2,057}{s}]$. Ne contient pas \SI{2,00}{s} $\rightarrow$ \textbf{Non compatible}.
\end{itemize}

\textbf{Étape 7 — Bulletin de qualités et verdict.}

\begin{center}
\begin{tabularx}{\linewidth}{|l|X|X|}
\hline
\rowcolor{popblueL} Qualité & Amina (Série A) & Bouba (Série B) \\
\hline
Fidélité (Dispersion $s$) & Faible : $s_A = \SI{0,060}{s}$ & Excellente : $s_B = \SI{0,011}{s}$ \\
\hline
Justesse (Biais $\bar{T}-T_{\text{réf}}$) & Bonne : Biais $\approx \SI{0,000}{s}$ & Mauvaise : Biais $= \SI{+0,050}{s}$ \\
\hline
Compatibilité ($95\%$) & \textbf{Oui} & \textbf{Non} \\
\hline
\end{tabularx}
\end{center}

\textbf{Verdict du jury :} \textbf{Amina remporte le concours.}
\textbf{Justification :} En métrologie, la justesse (compatibilité avec la valeur vraie) est primordiale. Le résultat d'Amina, bien que dispersé (faible fidélité), est le seul compatible avec la valeur de référence. Le résultat de Bouba, très précis (fidèle), est affecté d'une \textbf{erreur systématique} non corrigée (déclenchement tardif, amplitude trop grande, frottements, comptage $N$ oscillations mal géré) qui le rend faux.

\textbf{Étape 8 — Améliorations du protocole (Fidélité + Justesse).}
\begin{enumerate}
\item \textbf{Réduire l'erreur de déclenchement (aléatoire + systématique) :} Mesurer la durée de $N=10$ ou $20$ oscillations complètes, puis diviser par $N$. L'incertitude de réaction ($\approx \SI{0,1}{s}$) est divisée par $N$ sur la période.
\item \textbf{Respecter les conditions du modèle (systématique) :} Lâcher le pendule avec une amplitude faible ($\theta_0 < 10^\circ$) pour valider l'approximation $\sin\theta \approx \theta$ et la formule $T=2\pi\sqrt{\ell/g}$.
\item \textbf{Repérage optimal :} Déclencher/arrêter au passage par la position d'équilibre (vitesse max, repère visuel net) et non aux extrémités (vitesse nulle, temps mort).
\item \textbf{Multiplier les séries :} Effectuer 3 à 5 séries de $N$ oscillations et moyenniser les périodes moyennes.
\end{enumerate}
\end{exoresolu}

\courssub{Bilan de l'activité}

Cette activité t'a permis de mettre en œuvre la chaîne métrologique complète : organisation des données, représentation graphique, calculs statistiques et expression normalisée du résultat. Elle a surtout mis en évidence deux idées essentielles :

\begin{itemize}
\item \textbf{Fidélité et justesse sont indépendantes} : une série très resserrée peut être entièrement fausse à cause d'une erreur systématique (série B), tandis qu'une série dispersée peut encadrer la valeur vraie (série A). L'écart-type mesure la fidélité ; seule la comparaison à une référence renseigne sur la justesse.
\item \textbf{Un résultat sans incertitude n'a pas de valeur scientifique} : c'est l'intervalle de confiance $[\bar{T} - U \,;\, \bar{T} + U]$ qui permet de trancher la question de la compatibilité, et non la seule moyenne.
\end{itemize}

\begin{attention}[À retenir pour le Probatoire]
\begin{itemize}
\item L'incertitude s'arrondit \textbf{par excès} à un chiffre significatif (deux si le premier chiffre est \textbf{1 ou 2}), et la moyenne est arrondie au même rang décimal que l'incertitude.
\item Écrire $T = (2,05 \pm 0,0065)$ s est incorrect : il faut écrire $T = (2,050 \pm 0,007)$ s.
\item Pour une mesure unique avec un appareil numérique de résolution $a$ : $u = a/\sqrt{12}$.
\item Pour une série : $u(\bar{x}) = s/\sqrt{n}$, $U = 2\, u(\bar{x})$ (niveau 95\%).
\end{itemize}
\end{attention}

Un bon protocole de mesure combine donc les deux exigences : réduire les erreurs aléatoires (répétition des mesures, moyenne, mesure sur $N$ périodes) et traquer les erreurs systématiques (étalonnage, choix du protocole, conditions de l'expérience). C'est exactement la démarche que tu devras reproduire lors des travaux pratiques évalués au Probatoire C.

\courssub{Exercices d'entraînement}

\begin{exercice}[Mesure unique : Voltmètre numérique]
On mesure la tension aux bornes d'une pile avec un multimètre numérique (classe 0,5 \%, résolution \SI{1}{mV} sur le calibre \SI{2}{V}). L'affichage indique $U = \SI{1,487}{V}$.
\begin{enumerate}
\item Quelle est l'incertitude de type B due à la résolution ?
\item Quelle est l'incertitude de type B due à la classe de l'appareil (erreur maximale $\pm 0,5\%$ de la lecture, loi rectangulaire) ?
\item Donnez l'incertitude composée $u_c(U)$ et le résultat final $U = (\dots \pm \dots)$ V pour $k=2$.
\end{enumerate}
\end{exercice}

\begin{corrige}[Exercice 1]
\begin{enumerate}
\item Résolution $a = \SI{1}{mV} = \SI{0,001}{V}$. $u_{\text{résol}} = a/\sqrt{12} = \SI{0,00029}{V}$.
\item Classe 0,5 \% : Erreur max $E_{\text{max}} = 0,005 \times 1,487 = \SI{0,00744}{V}$. Loi rectangulaire $\Rightarrow u_{\text{classe}} = E_{\text{max}}/\sqrt{3} = \SI{0,00429}{V}$.
\item $u_c = \sqrt{u_{\text{résol}}^2 + u_{\text{classe}}^2} \approx \sqrt{0,00029^2 + 0,00429^2} \approx \SI{0,0043}{V}$.
$U = 2 \times 0,0043 = \SI{0,0086}{V} \xrightarrow{\text{arrondi par excès 1 ch. sig.}} \SI{0,009}{V}$.
Moyenne au millième : $\bar{U} = \SI{1,487}{V}$.
\textbf{Résultat :} $\boxed{U = (1,487 \pm 0,009)\ \text{V}}$.
\end{enumerate}
\end{corrige}

\begin{exercice}[Série de mesures : Longueur d'une table]
Cinq élèves mesurent la longueur $L$ d'une table de TP avec un mètre ruban (appréciation \SI{1}{mm}). Résultats en cm : $150,2\ ;\ 150,5\ ;\ 149,8\ ;\ 150,1\ ;\ 150,4$.
\begin{enumerate}
\item Calculer $\bar{L}$, $s$, $u(\bar{L})$, $U$ ($k=2$).
\item Exprimer le résultat normalisé.
\item L'étiquette du fabricant indique $L = \SI{150,0}{cm}$. Le résultat est-il compatible ?
\end{enumerate}
\end{exercice}

\begin{corrige}[Exercice 2]
\begin{enumerate}
\item $\sum L_i = 751,0$ cm $\Rightarrow \bar{L} = \SI{150,2}{cm}$.
Écarts au carré : $(0,0)^2 + (0,3)^2 + (-0,4)^2 + (-0,1)^2 + (0,2)^2 = 0 + 0,09 + 0,16 + 0,01 + 0,04 = 0,30$.
$s = \sqrt{0,30 / 4} = \sqrt{0,075} \approx \SI{0,274}{cm}$.
$u(\bar{L}) = 0,274 / \sqrt{5} \approx \SI{0,122}{cm}$.
$U = 2 \times 0,122 = \SI{0,245}{cm} \xrightarrow{\text{arrondi par excès 1 ch. sig. (1er ch=2 $\rightarrow$ 2 ch. sig.)}} \SI{0,25}{cm}$.
\item $\bar{L}$ au centième (même rang que $U$) : $\SI{150,20}{cm}$.
$\boxed{L = (150,20 \pm 0,25)\ \text{cm}}$.
\item Intervalle $[\SI{149,95}{cm} \,;\, \SI{150,45}{cm}]$. Contient \SI{150,0}{cm} $\rightarrow$ \textbf{Compatible}.
\end{enumerate}
\end{corrige}

\newpage

\coursec{Méthodes et exercices résolus}

\coursec{Mesures, incertitudes et qualités d'un instrument de mesure}

\courssub{Mesure, erreur et incertitude}

\begin{definition}[Mesure d'une grandeur physique]
Une \cle{mesure} est l'ensemble des opérations ayant pour but de déterminer la valeur d'une grandeur physique (longueur, tension, période, concentration, etc.) par comparaison à un \cle{étalon} de même nature. Le résultat d'une mesure est l'association d'une \cle{valeur mesurée} (nombre) et de son \cle{unité} (SI de préférence).
\end{definition}

\begin{definition}[Vraie valeur, erreur et incertitude]
\begin{itemize}
\item La \cle{valeur vraie} (ou valeur conventionnelle) est la valeur idéale, inaccessible, qui caractériserait parfaitement la grandeur mesurée dans les conditions de l'expérience.
\item L'\cle{erreur de mesure} $E = x_{\text{mes}} - x_{\text{vraie}}$ est l'écart entre la valeur mesurée et la valeur vraie. Elle est \textbf{inconnue} car la valeur vraie l'est.
\item L'\cle{incertitude de mesure} $\Delta x$ (ou $u(x)$) est une \emph{estimation} de la dispersion des valeurs qui peuvent raisonnablement être attribuées au mesurand. Elle quantifie le \cle{doute} sur le résultat. Contrairement à l'erreur, l'incertitude est \textbf{connue} (évaluée) et toujours positive.
\end{itemize}
\end{definition}

\begin{propriete}[Origines des incertitudes -- Types A et B]
On distingue deux modes d'évaluation de l'incertitude-type $u(x)$ :
\begin{itemize}
\item \textbf{Type A (statistique) :} issue de l'analyse statistique d'une \emph{série de mesures répétées} (calcul de l'écart-type sur la moyenne). C'est la composante \emph{aléatoire}.
\item \textbf{Type B (autre) :} issue d'informations autres que des séries répétées : notice constructeur, certificat d'étalonnage, résolution de l'appareil, lecture d'une graduation, variations d'influence (température, tension d'alimentation...). C'est la composante \emph{systématique} (ou traitée comme telle).
\end{itemize}
En Première C, on traite principalement le Type A (série de mesures) et le Type B simple (lecture unique sur instrument analogique ou numérique).
\end{propriete}

\begin{aretenir}[Vocabulaire rigoureux]
\begin{itemize}
\item On ne dit jamais « l'erreur est de 0,02 s » (l'erreur est inconnue).
\item On dit : « l'incertitude élargie est $U = 0,02\ \text{s}$ » ou « l'incertitude-type est $u = 0,01\ \text{s}$ ».
\item L'incertitude s'exprime \textbf{toujours} dans la même unité que la grandeur mesurée.
\end{itemize}
\end{aretenir}

\courssub{Qualités d'un instrument de mesure}

\begin{definition}[Sensibilité]
La \cle{sensibilité} $S$ d'un instrument est le rapport de la variation de l'indication (sortie) à la variation correspondante de la grandeur mesurée (entrée) :
\[ S = \frac{\Delta \text{Indication}}{\Delta \text{Mesurand}}. \]
Un instrument très sensible réagit fortement à une petite variation du mesurand (ex. : galvanomètre à haute sensibilité, thermomètre à dilatation importante). Unité : dépend des grandeurs (ex. $\text{mm}/^\circ\text{C}$, $\text{V}/\text{A}$).
\end{definition}

\begin{definition}[Fidélité (Répétabilité / Reproductibilité)]
La \cle{fidélité} caractérise la capacité d'un instrument à donner des indications voisines pour des mesures répétées d'une même grandeur dans des conditions identiques (même opérateur, même appareil, court intervalle de temps $\rightarrow$ \cle{répétabilité}) ou changées (opérateur, appareil, temps, lieu $\rightarrow$ \cle{reproductibilité}).
\begin{itemize}
\item Une bonne fidélité signifie une faible \cle{dispersion aléatoire} (écart-type $\sigma_{n-1}$ faible).
\item Elle ne garantit pas la justesse (on peut être fidèle mais faux si l'instrument est mal réglé).
\end{itemize}
\end{definition}

\begin{definition}[Justesse (Absence de biais systématique)]
La \cle{justesse} caractérise la proximité entre la \emph{moyenne} d'un grand nombre de résultats de mesure et la \cle{valeur vraie} (ou une valeur de référence).
\begin{itemize}
\item Une bonne justesse signifie un faible \cle{biais systématique} (erreur systématique).
\item Elle s'améliore par \cle{étalonnage} (comparaison à un étalon) et correction.
\end{itemize}
\end{definition}

\begin{definition}[Précision / Exactitude]
\begin{itemize}
\item La \cle{précision} (au sens large VIM) est la proximité entre les indications entre elles $\rightarrow$ synonyme de \textbf{fidélité}.
\item L'\cle{exactitude} (accuracy) est la proximité entre un résultat de mesure et la valeur vraie. Elle combine \textbf{fidélité} et \textbf{justesse}.
\end{itemize}
\textbf{Image mnémotechnique :} Cible de tir. Fidélité = groupement serré des impacts. Justesse = groupement centré sur le centre. Exactitude = groupement serré \emph{et} centré.
\end{definition}

\begin{center}
\begin{tabularx}{\linewidth}{|l|X|X|}\hline
\rowcolor{popblueL} \textbf{Qualité} & \textbf{Caractère} & \textbf{Amélioration} \\ \hline
Sensibilité & Variation indication / variation mesurand & Choix du capteur, amplification \\ \hline
Fidélité & Dispersion aléatoire (écart-type) & Instrument stable, conditions constantes, répétitions \\ \hline
Justesse & Biais systématique (écart à la vraie valeur) & Étalonnage, correction du zéro, méthode sans biais \\ \hline
Exactitude & Fidélité + Justesse & Les deux précédentes \\ \hline
\end{tabularx}
\end{center}

\begin{savaistu}[Contexte camerounais : ENEO et les compteurs électriques]
Les compteurs d'énergie électrique (kWh) installés par ENEO doivent être \textbf{justes} (facturation correcte) et \textbf{fidèles} (même consommation affichée pour deux foyers identiques). Leur \cle{classe de précision} (ex. Classe 1 ou 2) garantit que l'erreur relative maximale en conditions nominales est de $\pm 1\,\%$ ou $\pm 2\,\%$. Un contrôle métrologique périodique (vérification primitive puis périodique) assure le maintien de ces qualités. Un compteur sensible détecte les faibles consommations (veilles d'appareils) ; un compteur peu sensible ne « voit » pas les fuites de courant.
\end{savaistu}

\courssub{Incertitude sur une mesure unique (Type B)}

\begin{methode}[Incertitude sur une mesure unique]
\begin{enumerate}
\item \textbf{Identifier le type d'instrument} : analogique (aiguille, règle graduée, pipette, burette) ou numérique (afficheur LCD/LED, multimètre).
\item \textbf{Instrument analogique} : l'incertitude de lecture est liée à la plus petite graduation $g$ (ou au pouvoir de résolution de l'œil) :
\[ \Delta x = \frac{g}{2} \quad \text{(lecture soigneuse, parallaxe évitée)}. \]
Si la lecture est difficile (parallaxe, vibrations), on peut prendre $\Delta x = g$.
\item \textbf{Instrument numérique} :
\begin{itemize}
\item Si la notice constructeur donne une spécification : l'appliquer (souvent $\pm (p\,\% \text{ de la lecture} + n \text{ unités du dernier chiffre})$).
\item À défaut : l'incertitude de lecture vaut \textbf{une unité du dernier chiffre affiché} (ex. : afficheur $12,34\ \text{V} \Rightarrow \Delta U = 0,01\ \text{V}$).
\end{itemize}
\item \textbf{Autres sources} : incertitude de zéro (décalage à vide), incertitude de réglage, influence de la température. En Première, on retient souvent la \textbf{plus défavorable} ou la \textbf{somme simple} si l'énoncé le demande.
\item \textbf{Rédiger} : $x = x_{\text{mes}} \pm \Delta x$, unité commune, chiffres significatifs cohérents (voir §6).
\end{enumerate}
\end{methode}

\begin{exoresolu}[Mesure de la longueur d'une feuille à la règle]
Un élève mesure la longueur $L$ d'une feuille de papier avec une règle graduée au millimètre. Il lit $L = 29,7\ \text{cm}$, les deux extrémités de la feuille étant alignées avec des graduations.
\begin{enumerate}
\item Évaluer l'incertitude de lecture $\Delta L$.
\item Écrire le résultat de la mesure sous forme d'encadrement.
\end{enumerate}
\tcblower
\textbf{1. Incertitude de lecture.} La plus petite graduation vaut $g = 1\ \text{mm} = 0,1\ \text{cm}$. Pour une lecture soigneuse sur un instrument analogique :
\[ \Delta L = \frac{g}{2} = 0,5\ \text{mm} = 0,05\ \text{cm}. \]

\textbf{2. Résultat.} Puisque l'incertitude porte sur le centième de centimètre, la valeur mesurée doit être écrite avec la même précision : on écrit $29,70$ cm et non $29,7$ cm.
\[ L = (29,70 \pm 0,05)\ \text{cm}, \]
c'est-à-dire l'encadrement $29,65\ \text{cm} \leq L \leq 29,75\ \text{cm}$.
\textbf{Conclusion :} $L = (29,70 \pm 0,05)\ \text{cm}$.
\end{exoresolu}

\begin{exoresolu}[Mesure de tension au multimètre numérique]
Un multimètre numérique affiche $U = 12,34\ \text{V}$ sur le calibre $20\ \text{V}$. La notice indique : « Précision $\pm (0,5\,\% \text{ de la lecture} + 2 \text{ digits}) $ ».
\begin{enumerate}
\item Calculer l'incertitude absolue $\Delta U$.
\item Exprimer le résultat.
\end{enumerate}
\tcblower
\textbf{1. Calcul.}
\begin{align*}
\Delta U &= 0,5\,\% \times 12,34 + 2 \times 0,01 \\
&= 0,0617 + 0,02 = 0,0817\ \text{V} \approx 0,08\ \text{V} \text{ (arrondi par excès à 1 chiffre significatif)}.
\end{align*}

\textbf{2. Résultat.} L'incertitude porte sur le centième de volt : on arrondit la lecture au centième (elle l'est déjà).
\[ U = (12,34 \pm 0,08)\ \text{V}. \]
\textbf{Conclusion :} $U = (12,34 \pm 0,08)\ \text{V}$.
\end{exoresolu}

\courssub{Série de mesures : moyenne et écart-type (Type A)}

\begin{propriete}[Moyenne et écart-type expérimental]
Pour $n$ mesures $x_1, x_2, \ldots, x_n$ d'une même grandeur, effectuées dans les mêmes conditions :
\begin{itemize}
\item La \cle{moyenne arithmétique} est la meilleure estimation de la valeur vraie :
\[ \bar{x} = \frac{1}{n}\sum_{i=1}^{n} x_i. \]
\item L'\cle{écart-type expérimental} (ou écart-type échantillon) quantifie la dispersion (fidélité) :
\[ \sigma_{n-1} = \sqrt{\frac{\displaystyle\sum_{i=1}^{n}(x_i - \bar{x})^2}{n-1}}. \]
\item \textbf{Reflexe calculatrice} : Mode \texttt{STAT} (ou \texttt{SD}), entrer les données, lire $\bar{x}$ et $\sigma_{n-1}$ (souvent noté $s$ ou $x\sigma_{n-1}$). Ne pas confondre avec $\sigma_n$ (écart-type population, diviseur $n$).
\item \textbf{Valeurs aberrantes} : Si une valeur s'éloigne anormalement des autres (faute de manipulation, lecture erronée), on peut l'éliminer \textbf{en le justifiant explicitement} avant de recalculer $\bar{x}$ et $\sigma_{n-1}$.
\end{itemize}
\end{propriete}

\begin{exoresolu}[Période d'un pendule simple]
Pour déterminer la période d'un pendule simple, un élève de Première C mesure 5 fois la durée de 10 oscillations et en déduit la période $T$ :
\begin{center}
\begin{tabular}{|l|ccccc|}\hline
Mesure $n^{\circ}$ & 1 & 2 & 3 & 4 & 5 \\ \hline
$T$ (s) & 1,42 & 1,45 & 1,43 & 1,44 & 1,41 \\ \hline
\end{tabular}
\end{center}
\begin{enumerate}
\item Calculer la valeur moyenne $\bar{T}$.
\item Calculer l'écart-type expérimental $\sigma_{n-1}$.
\item Commenter la fidélité.
\end{enumerate}
\tcblower
\textbf{1. Moyenne.}
\[ \bar{T} = \frac{1,42 + 1,45 + 1,43 + 1,44 + 1,41}{5} = \frac{7,15}{5} = 1,430\ \text{s}. \]

\textbf{2. Écart-type.} Calcul des écarts à la moyenne :
\begin{align*}
\sum (T_i - \bar{T})^2 &= (-0,010)^2 + (0,020)^2 + (0,000)^2 + (0,010)^2 + (-0,020)^2 \\
&= 0,0001 + 0,0004 + 0 + 0,0001 + 0,0004 = 0,0010\ \text{s}^2.
\end{align*}
D'où :
\[ \sigma_{n-1} = \sqrt{\frac{0,0010}{5-1}} = \sqrt{2,5\times 10^{-4}} \approx 0,0158\ \text{s} \approx 0,016\ \text{s}. \]

\textbf{3. Commentaire.} L'incertitude relative $\sigma_{n-1}/\bar{T} \approx 1,1\,\%$. La dispersion est faible : les mesures sont \textbf{fidèles}.
\textbf{Conclusion :} $\bar{T} = 1,430\ \text{s}$ et $\sigma_{n-1} \approx 0,016\ \text{s}$.
\end{exoresolu}

\courssub{Incertitude-type, incertitude élargie et intervalle de confiance}

\begin{propriete}[Incertitude-type sur la moyenne $u(\bar{x})$]
\begin{itemize}
\item \textbf{Mesure unique (Type B) :} $u(x) = \Delta x$ (évaluée selon §3).
\item \textbf{Série de $n$ mesures répétées (Type A) :} L'incertitude-type sur la \emph{moyenne} est :
\[ u(\bar{x}) = \frac{\sigma_{n-1}}{\sqrt{n}}. \]
\textbf{Interprétation :} Répéter $n$ fois divise l'incertitude aléatoire par $\sqrt{n}$. C'est l'intérêt statistique de la répétition.
\item \textbf{Incertitude-type composée} : Si plusieurs sources indépendantes (Type A et Type B), on combine en quadratique : $u_c = \sqrt{u_A^2 + u_B^2 + \dots}$. (Hors programme détaillé Première, mais bon à savoir).
\end{itemize}
\end{propriete}

\begin{exoresolu}[Incertitude-type sur la période du pendule]
On reprend la série de l'exercice précédent : $\bar{T} = 1,430\ \text{s}$, $\sigma_{n-1} = 0,0158\ \text{s}$, $n = 5$.
\begin{enumerate}
\item Calculer l'incertitude-type $u(\bar{T})$.
\item Que devient $u(\bar{T})$ si l'élève effectue $n = 20$ mesures avec le même écart-type ?
\end{enumerate}
\tcblower
\textbf{1. Incertitude-type pour $n=5$.}
\[ u(\bar{T}) = \frac{\sigma_{n-1}}{\sqrt{n}} = \frac{0,0158}{\sqrt{5}} = \frac{0,0158}{2,236} \approx 7,1\times 10^{-3}\ \text{s}. \]

\textbf{2. Cas $n = 20$.}
\[ u'(\bar{T}) = \frac{0,0158}{\sqrt{20}} \approx 3,5\times 10^{-3}\ \text{s}. \]
En multipliant le nombre de mesures par 4, l'incertitude-type est divisée par $\sqrt{4} = 2$.
\textbf{Conclusion :} $u(\bar{T}) \approx 0,0071\ \text{s}$ pour 5 mesures ; répéter davantage améliore la précision, mais selon une loi en $1/\sqrt{n}$ (rendements décroissants).
\end{exoresolu}

\begin{propriete}[Incertitude élargie $U$ et intervalle de confiance]
L'incertitude-type $u(\bar{x})$ correspond à un niveau de confiance d'environ $68\,\%$ (facteur $k=1$, loi normale). Pour exprimer un résultat avec un niveau de confiance plus élevé (exigé au Probatoire : $95\,\%$ ou $99\,\%$), on définit l'\cle{incertitude élargie} :
\[ U = k \times u(\bar{x}) \]
où $k$ est le \cle{facteur d'élargissement} :
\begin{center}
\begin{tabular}{|l|ccc|}\hline
Niveau de confiance & $68\,\%$ & $95\,\%$ & $99\,\%$ \\ \hline
Facteur $k$ & 1 & 2 & 3 \\ \hline
\end{tabular}
\end{center}
L'\cle{intervalle de confiance} au niveau $P$ est :
\[ \bar{x} - U \leq x \leq \bar{x} + U. \]
\textbf{Sens :} La valeur vraie a une probabilité $P$ (ex. $95\,\%$) de se trouver dans cet intervalle.
\end{propriete}

\begin{exoresolu}[Concentration d'une solution de glucose -- Hôpital de district]
Au laboratoire d'un hôpital de district, un dosage donne pour la concentration massique d'une solution de glucose : $\bar{c} = 5,02\ \text{g}\cdot\text{L}^{-1}$ avec une incertitude-type $u(c) = 0,04\ \text{g}\cdot\text{L}^{-1}$.
\begin{enumerate}
\item Déterminer l'incertitude élargie $U$ pour un niveau de confiance de $95\,\%$.
\item Donner l'intervalle de confiance correspondant.
\end{enumerate}
\tcblower
\textbf{1. Incertitude élargie.} Pour $95\,\%$, $k = 2$ :
\[ U = k \times u(c) = 2 \times 0,04 = 0,08\ \text{g}\cdot\text{L}^{-1}. \]

\textbf{2. Intervalle de confiance.}
\[ \bar{c} - U \leq c \leq \bar{c} + U \quad \Longleftrightarrow \quad 5,02 - 0,08 \leq c \leq 5,02 + 0,08, \]
soit :
\[ 4,94\ \text{g}\cdot\text{L}^{-1} \leq c \leq 5,10\ \text{g}\cdot\text{L}^{-1}. \]
\textbf{Conclusion :} La concentration massique en glucose est comprise entre $4,94$ et $5,10\ \text{g}\cdot\text{L}^{-1}$ avec un niveau de confiance de $95\,\%$.
\end{exoresolu}

\courssub{Expression finale d'un résultat de mesure}

\begin{methode}[Rédiger le résultat final $x = \bar{x} \pm U$]
\begin{enumerate}
\item \textbf{Calculer l'incertitude élargie} $U = k \times u(\bar{x})$ avec le $k$ demandé ($k=2$ pour $95\,\%$).
\item \textbf{Arrondir $U$} à \textbf{un seul chiffre significatif} (deux si le premier chiffre significatif est $1$), \textbf{toujours par excès} (arrondi supérieur).
\item \textbf{Arrondir la valeur mesurée} $\bar{x}$ au \textbf{même rang décimal} (même puissance de 10) que $U$ arrondi.
\item \textbf{Écrire} sous la forme $x = (\bar{x} \pm U)\ \text{unité}$, avec l'unité factorisée.
\item \textbf{Préciser} le niveau de confiance : « $k = 2$, niveau de confiance $95\,\%$ ».
\item \textbf{Vérifier} la cohérence : $U \ll \bar{x}$ (sinon mesure de mauvaise qualité).
\end{enumerate}
\end{methode}

\begin{exoresolu}[Résultat final de la mesure de la période]
Les mesures répétées de la période du pendule ont donné $\bar{T} = 1,4300\ \text{s}$ et $u(\bar{T}) = 0,0071\ \text{s}$ ($n = 5$ mesures).
\begin{enumerate}
\item Calculer l'incertitude élargie au niveau de confiance $95\,\%$.
\item Exprimer le résultat final de la mesure.
\end{enumerate}
\tcblower
\textbf{1. Incertitude élargie.} Avec $k = 2$ (niveau de confiance $95\,\%$) :
\[ U = 2 \times 0,0071 = 0,0142\ \text{s}. \]

\textbf{2. Arrondis.}
\begin{itemize}
\item $U = 0,0142\ \text{s}$. Le premier chiffre significatif est $1$ : on conserve \textbf{deux} chiffres significatifs $\rightarrow 0,014\ \text{s}$. Arrondi par excès $\rightarrow \mathbf{U = 0,015\ \text{s}}$.
\item $\bar{T} = 1,4300\ \text{s}$. $U$ porte sur le millième de seconde (troisième décimale). On arrondit $\bar{T}$ au millième $\rightarrow \mathbf{\bar{T} = 1,430\ \text{s}}$.
\end{itemize}

\textbf{3. Résultat final.}
\[ T = (1,430 \pm 0,015)\ \text{s} \quad (k = 2,\ \text{niveau de confiance } 95\,\%). \]
L'incertitude relative vaut $\dfrac{U}{\bar{T}} \approx \dfrac{0,015}{1,430} \approx 1,0\,\%$ : la mesure est de bonne qualité.
\textbf{Conclusion :} La période du pendule est $T = (1,430 \pm 0,015)\ \text{s}$ à $95\,\%$ de confiance.
\end{exoresolu}

\courssub{Savoir-faire 1 : Organiser les données expérimentales}

\begin{methode}[Organiser les données : tableau et représentation graphique]
\begin{enumerate}
\item \textbf{Recenser} les grandeurs mesurées et leurs unités SI ; vérifier la cohérence des unités avant tout calcul (conversion si nécessaire).
\item \textbf{Construire un tableau} à double entrée : une ligne par grandeur, une colonne par essai ; indiquer l'unité dans l'en-tête de colonne (jamais dans chaque case).
\item \textbf{Choisir les axes} du graphique : grandeur indépendante (contrôlée) en abscisse, grandeur dépendante (mesurée) en ordonnée ; graduer de façon régulière, nommer chaque axe avec son symbole et son unité (ex. $I\ (\text{A})$).
\item \textbf{Placer les points} avec une croix fine ($\times$) ou un petit cercle ; ne jamais les relier « point à point » : tracer la \cle{courbe moyenne} (droite de régression ou courbe lisse) passant au plus près de l'ensemble des points.
\item \textbf{Exploiter} : si la courbe moyenne est une droite passant par l'origine, les grandeurs sont proportionnelles ; calculer le coefficient directeur $a = \dfrac{\Delta y}{\Delta x}$ avec deux points \textbf{éloignés de la droite tracée} (pas des points expérimentaux).
\end{enumerate}
\end{methode}

\begin{exoresolu}[Exploitation des mesures d'un volt-ampèremètre -- Lycée de Nkongsamba]
Au laboratoire, un élève mesure la tension $U$ aux bornes d'un conducteur ohmique pour différentes valeurs de l'intensité $I$ :
\begin{center}
\begin{tabular}{|l|ccccc|}\hline
$I$ (mA) & 0 & 20 & 40 & 60 & 80 \\ \hline
$U$ (V) & 0 & 0,90 & 1,85 & 2,75 & 3,70 \\ \hline
\end{tabular}
\end{center}
\begin{enumerate}
\item Tracer la caractéristique $U = f(I)$.
\item Montrer que le dipôle est un conducteur ohmique et déterminer sa résistance $R$.
\end{enumerate}
\tcblower
\textbf{1. Tracé.} On place $I$ (en A) en abscisse et $U$ (en V) en ordonnée. Conversion : $20\ \text{mA} = 0,020\ \text{A}$, etc. Les points $(0\,;\,0)$, $(0,020\,;\,0,90)$, $(0,040\,;\,1,85)$, $(0,060\,;\,2,75)$, $(0,080\,;\,3,70)$ sont sensiblement alignés avec l'origine.
\begin{popfigure}
\begin{tikzpicture}
\begin{axis}[width=0.85\linewidth,height=6cm,xlabel={$I$ (A)},ylabel={$U$ (V)},xmin=0,xmax=0.09,ymin=0,ymax=4.5,grid=both,grid style={popline!40},tick label style={/pgf/number format/use comma}]
\addplot[only marks,mark=+,popink,mark size=3pt] coordinates {(0,0) (0.02,0.9) (0.04,1.85) (0.06,2.75) (0.08,3.7)};
\addplot[popcurve,domain=0:0.085]{46.25*x};
\end{axis}
\end{tikzpicture}
\legende{Caractéristique $U=f(I)$ : les points sont alignés avec l'origine.}
\end{popfigure}

\textbf{2. Nature du dipôle et résistance.} La courbe moyenne est une \textbf{droite passant par l'origine} : $U$ est proportionnelle à $I$, c'est la loi d'Ohm $U = RI$. Le dipôle est donc un \cle{conducteur ohmique}.

On calcule le coefficient directeur avec deux points éloignés \emph{de la droite tracée}, par exemple l'origine $(0\,;\,0)$ et le point de la droite à $I = 0,080\ \text{A}$ ($U = 3,70\ \text{V}$) :
\[ R = \frac{\Delta U}{\Delta I} = \frac{3,70 - 0}{0,080 - 0} = 46,25\ \Omega. \]
Analyse dimensionnelle : $[U]/[I] = \text{V}\cdot\text{A}^{-1} = \Omega$.
Arrondi à 2 chiffres significatifs (données à 2--3 chiffres) : $R \approx \SI{46}{\ohm}$.
\textbf{Conclusion :} Le dipôle est un conducteur ohmique de résistance $R \approx \SI{46}{\ohm}$.
\end{exoresolu}

\begin{attention}[Les 5 erreurs qui coûtent des points au Probatoire]
\begin{enumerate}
\item \textbf{Confondre erreur et incertitude.} L'erreur (écart à la valeur vraie) est inconnaissable ; l'incertitude est une \emph{estimation} du doute. Ne jamais écrire « l'erreur est de $0,02\ \text{s}$ » : écrire « l'incertitude élargie est $U = 0,02\ \text{s}$ ».
\item \textbf{Oublier le facteur $\sqrt{n}$.} Pour une série de mesures, l'incertitude-type porte sur la \emph{moyenne} : $u(\bar{x}) = \sigma_{n-1}/\sqrt{n}$, et non $\sigma_{n-1}$ seul. C'est la condition d'application la plus souvent oubliée.
\item \textbf{Mal choisir le facteur $k$.} Lire attentivement l'énoncé : niveau de confiance $95\,\% \Rightarrow k = 2$ ; $99\,\% \Rightarrow k = 3$. Écrire $U = k\,u$ en explicitant la valeur de $k$ utilisée.
\item \textbf{Incohérence des chiffres significatifs.} Interdit : $T = (1,4300 \pm 0,0142)\ \text{s}$. L'incertitude s'arrondit à 1 chiffre significatif (2 si le premier est 1), par excès ($0,015\ \text{s}$) et la valeur mesurée au même rang décimal ($1,430\ \text{s}$). Un résultat mal arrondi est sanctionné même si le calcul est juste.
\item \textbf{Oublier l'unité ou mélanger les unités.} $U$ et $\bar{x}$ doivent avoir la \emph{même} unité, factorisée : $(1,430 \pm 0,015)\ \text{s}$. Convertir avant tout calcul (mA en A, cm en m si nécessaire) et vérifier l'homogénéité de chaque formule.
\end{enumerate}
\end{attention}

\begin{exercice}[Application : Mesure de la résistance d'un fil de cuivre]
On mesure la résistance $R$ d'un fil de cuivre à l'ohmmètre numérique (calibre $200\ \Omega$). Cinq mesures donnent : $47,2\ \Omega$ ; $47,5\ \Omega$ ; $47,3\ \Omega$ ; $47,4\ \Omega$ ; $47,1\ \Omega$. La notice de l'appareil indique une incertitude de $\pm (0,5\,\% + 2 \text{ digits})$.
\begin{enumerate}
\item Calculer la moyenne $\bar{R}$ et l'écart-type $\sigma_{n-1}$.
\item Déterminer l'incertitude-type $u_A(\bar{R})$ (Type A) et l'incertitude-type $u_B(R)$ (Type B, prise sur une mesure unique).
\item En combinant les deux sources (somme simple ou quadratique selon habitude), donner l'incertitude élargie $U$ pour $k=2$ et le résultat final.
\end{enumerate}
\end{exercice}

\begin{corrige}[Exercice n°1]
\textbf{1. Moyenne et écart-type.}
$\bar{R} = (47,2+47,5+47,3+47,4+47,1)/5 = 236,5/5 = 47,30\ \Omega$.
Écarts : $-0,10 ; +0,20 ; 0,00 ; +0,10 ; -0,20$.
Somme carrés : $0,01+0,04+0+0,01+0,04 = 0,10$.
$\sigma_{n-1} = \sqrt{0,10/4} = \sqrt{0,025} \approx 0,158\ \Omega$.

\textbf{2. Incertitudes-types.}
Type A : $u_A(\bar{R}) = \sigma_{n-1}/\sqrt{5} = 0,158/2,236 \approx 0,071\ \Omega$.
Type B (sur une mesure, ex. $47,3\ \Omega$) : $u_B = 0,5\,\%\times 47,3 + 2\times 0,1 = 0,2365 + 0,2 = 0,4365\ \Omega \approx 0,44\ \Omega$.

\textbf{3. Incertitude composée et résultat.}
$u_c = \sqrt{u_A^2 + u_B^2} \approx \sqrt{0,071^2 + 0,44^2} \approx 0,445\ \Omega$ (domination Type B).
$U = 2 \times 0,445 = 0,89\ \Omega \rightarrow$ Arrondi par excès 1 chiffre sig. : $\mathbf{U = 1\ \Omega}$.
$\bar{R}$ arrondi à l'unité (même rang que $U$) : $\mathbf{\bar{R} = 47\ \Omega}$.
\textbf{Résultat final :} $\mathbf{R = (47 \pm 1)\ \Omega}$ ($k=2$, $95\,\%$).
\textit{Remarque :} La précision de l'appareil (Type B) domine la dispersion statistique (Type A). Répéter les mesures n'améliore presque pas le résultat final.
\end{corrige}

\newpage

\coursec{Exercices}

\courssub{Je vérifie mes connaissances}

\begin{exercice}[QCM : vocabulaire des qualités d'un instrument]
Pour chaque situation, choisir la qualité d'instrument concernée parmi : \textbf{a)} justesse ; \textbf{b)} fidélité ; \textbf{c)} sensibilité ; \textbf{d)} précision.
\begin{enumerate}
\item Un thermomètre donne des mesures très groupées entre elles, mais toutes décalées de \SI{1,5}{\celsius} au-dessus de la valeur vraie.
\item Une balance détecte l'ajout d'un grain de riz de masse \SI{0,02}{g}.
\item Dix pesées successives du même sac donnent des valeurs très dispersées autour de la valeur vraie.
\item Un voltmètre répète quasiment la même indication lors de mesures successives dans des conditions identiques.
\end{enumerate}
\end{exercice}

\begin{exercice}[Vrai ou faux, en justifiant]
Répondre par vrai ou faux, puis justifier en une ou deux phrases.
\begin{enumerate}
\item L'erreur de mesure et l'incertitude de mesure désignent exactement la même grandeur.
\item Un instrument fidèle est forcément juste.
\item L'incertitude-type sur la moyenne de $n$ mesures indépendantes diminue lorsque $n$ augmente.
\item Un intervalle de confiance à \SI{99}{\%} est plus étroit qu'un intervalle de confiance à \SI{95}{\%} pour la même série de mesures.
\end{enumerate}
\end{exercice}

\begin{exercice}[Définitions]
\begin{enumerate}
\item Définir : erreur de mesure, incertitude-type, incertitude élargie.
\item Donner la relation entre l'incertitude élargie $U$, l'incertitude-type $u$ et le facteur d'élargissement $k$.
\item Quelle valeur de $k$ correspond, pour une distribution gaussienne, à un niveau de confiance d'environ \SI{95}{\%} ?
\end{enumerate}
\end{exercice}

\begin{exercice}[Calculs directs]
\begin{enumerate}
\item Un voltmètre analogique de classe 2,5 est utilisé sur le calibre \SI{30}{V}. Calculer l'incertitude absolue maximale due à l'instrument.
\item On mesure une longueur $L = \SI{12,4}{cm}$ avec une incertitude absolue $\Delta L = \SI{0,2}{cm}$. Calculer l'incertitude relative en pourcentage.
\item On effectue $n = 16$ mesures d'une même durée ; l'écart-type expérimental vaut $s = \SI{0,08}{s}$. Calculer l'incertitude-type sur la moyenne.
\end{enumerate}
\end{exercice}

\courssub{Je m'entraîne}

\begin{exercice}[Voltmètre analogique]
On mesure une tension avec un voltmètre analogique de classe 1,5, sur le calibre \SI{15}{V}. L'aiguille s'arrête sur la graduation correspondant à \SI{9,4}{V}, et l'incertitude de lecture est estimée à une demi-graduation, soit \SI{0,1}{V}.
\begin{enumerate}
\item Calculer l'incertitude absolue due à la classe de l'appareil.
\item En combinant l'incertitude instrumentale et l'incertitude de lecture (on prendra la somme), exprimer le résultat de la mesure sous la forme $U_{\text{mes}} \pm \Delta U$.
\item Calculer l'incertitude relative de cette mesure, en pourcentage.
\end{enumerate}
\end{exercice}

\begin{exercice}[Série de huit mesures de masse]
On pèse huit fois le même objet avec la même balance et on obtient (en grammes) :
\begin{center}
49,8 \quad ; \quad 50,2 \quad ; \quad 50,0 \quad ; \quad 49,9 \quad ; \quad 50,1 \quad ; \quad 50,3 \quad ; \quad 49,7 \quad ; \quad 50,0
\end{center}
\begin{enumerate}
\item Organiser ces données dans un tableau.
\item Calculer la valeur moyenne $\bar{m}$ de la série.
\item Calculer l'écart-type expérimental $s_{n-1}$.
\item En déduire l'incertitude-type $u(\bar{m})$ sur la moyenne.
\end{enumerate}
\end{exercice}

\begin{exercice}[Lecture sur deux instruments]
Pour mesurer la longueur d'une règle de tableau, un élève dispose :
\begin{itemize}
\item d'une règle graduée au millimètre (graduation la plus petite : \SI{1}{mm}) ;
\item d'un pied à coulisse numérique affichant au centième de millimètre (affichage : \SI{0,01}{mm}).
\end{itemize}
\begin{enumerate}
\item Donner l'incertitude de lecture associée à chaque instrument, en justifiant (demi-graduation pour l'instrument analogique, $\pm 1$ unité d'affichage pour le numérique).
\item Comparer la résolution des deux instruments.
\item Lequel est le plus sensible ? Le plus précis ? Justifier.
\end{enumerate}
\end{exercice}

\begin{exercice}[Compatibilité de deux résultats]
Deux groupes d'élèves mesurent la même longueur et trouvent :
\[ L_1 = (25,40 \pm 0,05)~\text{cm} \qquad \text{et} \qquad L_2 = (25,62 \pm 0,05)~\text{cm} \]
\begin{enumerate}
\item Écrire l'intervalle de valeurs associé à chaque résultat.
\item Ces deux résultats sont-ils compatibles ? Justifier.
\item Proposer une explication physique à l'écart observé (on pensera à une erreur systématique) et une vérification expérimentale permettant de trancher.
\end{enumerate}
\end{exercice}

\courssub{Je me prépare au Probatoire}

\begin{exercice}[Mesure d'une tension au laboratoire du lycée]
Au laboratoire de physique d'un lycée de Yaoundé, un élève de Première C mesure la tension aux bornes d'une lampe à l'aide d'un voltmètre analogique de classe 1,5. Il utilise le calibre \SI{15}{V} ; le cadran comporte 150 divisions et l'aiguille s'immobilise sur la division 94.
\begin{enumerate}
\item Déterminer la valeur $U$ de la tension mesurée.
\item Calculer l'incertitude absolue $\Delta U_{\text{instr}}$ due à la classe de l'appareil.
\item L'élève estime son incertitude de lecture à une demi-division. Calculer $\Delta U_{\text{lect}}$.
\item En prenant pour incertitude totale $\Delta U = \Delta U_{\text{instr}} + \Delta U_{\text{lect}}$, exprimer le résultat sous la forme $U \pm \Delta U$ avec un nombre raisonnable de chiffres significatifs.
\item Calculer l'incertitude relative, en pourcentage, et commenter la qualité de la mesure.
\item L'enseignant affirme que la tension vaut exactement \SI{9,00}{V}. Cette valeur est-elle compatible avec le résultat de l'élève ? Justifier.
\end{enumerate}
\end{exercice}

\begin{exercice}[Pesées répétées d'un sac de riz à Douala]
Un commerçant du marché de Douala vend des sacs de riz annoncés à \SI{5,000}{kg}. Un inspecteur effectue six pesées répétées du même sac avec une balance affichant au gramme près. Résultats (en kg) :
\begin{center}
4,982 \quad ; \quad 4,991 \quad ; \quad 4,987 \quad ; \quad 4,995 \quad ; \quad 4,988 \quad ; \quad 4,985
\end{center}
\begin{enumerate}
\item Présenter ces données dans un tableau et calculer la masse moyenne $\bar{m}$.
\item Calculer l'écart-type expérimental $s_{n-1}$ de la série.
\item En déduire l'incertitude-type $u(\bar{m})$ sur la moyenne.
\item Calculer l'incertitude élargie $U$ pour un niveau de confiance de \SI{95}{\%} (on prendra $k = 2$).
\item Écrire le résultat final sous la forme normalisée $m = \bar{m} \pm U$, en veillant aux chiffres significatifs.
\item La masse annoncée de \SI{5,000}{kg} est-elle compatible avec les mesures au niveau de confiance \SI{95}{\%} ? Conclure sur la situation du commerçant et proposer une explication possible de l'écart.
\end{enumerate}
\end{exercice}

\begin{exercice}[Contrôle qualité d'un lot de comprimés]
Un laboratoire pharmaceutique de Yaoundé contrôle la masse de comprimés dont la masse nominale est $m_0 = \SI{500}{mg}$. Un technicien prélève un comprimé et réalise dix pesées avec la même balance. Résultats (en mg) :
\begin{center}
\begin{tabular}{|l|*{10}{c|}}
\hline
Pesée $n^{\circ}$ & 1 & 2 & 3 & 4 & 5 & 6 & 7 & 8 & 9 & 10 \\
\hline
$m$ (mg) & 502 & 498 & 501 & 499 & 503 & 500 & 497 & 501 & 499 & 500 \\
\hline
\end{tabular}
\end{center}
\begin{enumerate}
\item Calculer la moyenne $\bar{m}$ et l'écart-type expérimental $s_{n-1}$ de cette série.
\item Calculer l'incertitude-type $u(\bar{m})$, puis l'incertitude élargie $U$ au niveau de confiance \SI{95}{\%} ($k = 2$).
\item Exprimer le résultat de la mesure sous forme normalisée.
\item Le cahier des charges impose que la masse d'un comprimé soit comprise entre \SI{495}{mg} et \SI{505}{mg}. Le comprimé contrôlé satisfait-il cette exigence ? Justifier à l'aide de l'intervalle de confiance.
\item Le technicien souhaite réduire l'incertitude-type sur la moyenne de moitié. Combien de pesées indépendantes devrait-il réaliser au total ? Justifier le calcul.
\item Citer deux sources d'erreur systématique possibles dans ce protocole et proposer pour chacune une action corrective.
\end{enumerate}
\end{exercice}

\newpage

\coursec{Corrigés des exercices}

\begin{corrige}[Exercice 1 — QCM : vocabulaire des qualités d'un instrument]
\begin{enumerate}
\item Mesures \cle{groupées} mais \cle{décalées} de la valeur vraie : l'instrument est fidèle (bonne répétabilité) mais pas juste. La qualité mise en défaut est la \textbf{justesse}, mais la qualité \emph{illustrée} par le caractère groupé est la fidélité. Réponse : \textbf{b) fidélité} (avec un défaut de justesse dû à une erreur systématique de \SI{+1,5}{\celsius}).
\item Détection d'une variation de masse aussi faible que \SI{0,02}{g} : réponse \textbf{c) sensibilité}.
\item Valeurs très dispersées lors de mesures répétées : c'est un défaut de \textbf{b) fidélité}.
\item Répétition (quasi) identique de l'indication dans des conditions identiques : c'est la définition même de la \textbf{b) fidélité}.
\end{enumerate}
\textbf{Point de vigilance :} ne pas confondre fidélité (dispersion des mesures répétées) et justesse (écart à la valeur vraie). Un instrument peut être fidèle sans être juste, et juste sans être fidèle.
\end{corrige}

\begin{corrige}[Exercice 2 — Vrai ou faux, en justifiant]
\begin{enumerate}
\item \textbf{Faux.} L'\cle{erreur} de mesure est la différence $m - m_{\text{vraie}}$ entre le résultat et la valeur vraie : c'est une grandeur \emph{idéale}, inconnue puisque la valeur vraie est inaccessible. L'\cle{incertitude} est une \emph{estimation} de l'étendue des valeurs que l'on peut raisonnablement attribuer au mesurande : c'est elle que l'on calcule et que l'on exprime.
\item \textbf{Faux.} La fidélité traduit seulement la faible dispersion des mesures répétées. Un instrument fidèle peut être affecté d'une erreur systématique (zéro décalé, mauvais étalonnage) : toutes ses indications, bien que groupées, seront fausses. Il est alors fidèle mais pas juste.
\item \textbf{Vrai.} Pour $n$ mesures indépendantes, l'incertitude-type sur la moyenne vaut $u(\bar{x}) = \dfrac{s_{n-1}}{\sqrt{n}}$ : elle diminue en $\dfrac{1}{\sqrt{n}}$ lorsque $n$ augmente.
\item \textbf{Faux.} C'est l'inverse : plus le niveau de confiance est élevé, plus le facteur d'élargissement $k$ est grand ($k = 2$ pour environ \SI{95}{\%}, $k = 3$ pour environ \SI{99}{\%}), donc plus l'intervalle $\bar{x} \pm k\,u$ est \emph{large}. Exiger plus de confiance coûte de la précision.
\end{enumerate}
\end{corrige}

\begin{corrige}[Exercice 3 — Définitions]
\begin{enumerate}
\item \textbf{Erreur de mesure :} différence entre la valeur mesurée et la valeur vraie du mesurande ; elle comporte une composante aléatoire et une composante systématique, et elle n'est pas connaissable exactement.
\textbf{Incertitude-type $u$ :} estimation de la dispersion des valeurs attribuables au mesurande, exprimée sous la forme d'un écart-type (sur une mesure ou sur la moyenne).
\textbf{Incertitude élargie $U$ :} demi-largeur d'un intervalle centré sur le résultat, dans lequel la valeur vraie se trouve avec une probabilité (niveau de confiance) donnée.
\item La relation est :
\[ U = k \times u \]
où $k$ est le \cle{facteur d'élargissement}, choisi en fonction du niveau de confiance souhaité.
\item Pour une distribution gaussienne, un niveau de confiance d'environ \SI{95}{\%} correspond à $k = 2$ (plus précisément $k = 1{,}96$ ; on retient $k = 2$ au lycée).
\end{enumerate}
\end{corrige}

\begin{corrige}[Exercice 4 — Calculs directs]
\begin{enumerate}
\item Pour un appareil analogique, l'incertitude de classe se calcule \emph{sur le calibre} :
\[ \Delta U = \frac{\text{classe} \times \text{calibre}}{100} = \frac{2{,}5 \times 30}{100} = \SI{0,75}{V}. \]
\item Incertitude relative :
\[ \frac{\Delta L}{L} = \frac{0{,}2}{12{,}4} = 0{,}0161 \quad \text{soit environ } \SI{1,6}{\%}. \]
\item Incertitude-type sur la moyenne de $n = 16$ mesures :
\[ u(\bar{t}) = \frac{s}{\sqrt{n}} = \frac{0{,}08}{\sqrt{16}} = \frac{0{,}08}{4} = \SI{0,02}{s}. \]
\end{enumerate}
\textbf{Point de vigilance :} l'incertitude de classe se calcule toujours à partir du \emph{calibre}, jamais à partir de la valeur lue.
\end{corrige}

\begin{corrige}[Exercice 5 — Voltmètre analogique]
\begin{enumerate}
\item Incertitude due à la classe :
\[ \Delta U_{\text{instr}} = \frac{1{,}5 \times 15}{100} = \SI{0,225}{V} \approx \SI{0,23}{V}. \]
\item Incertitude totale (somme, comme demandé) :
\[ \Delta U = \Delta U_{\text{instr}} + \Delta U_{\text{lect}} = 0{,}225 + 0{,}1 = 0{,}325 \approx \SI{0,3}{V}. \]
On arrondit l'incertitude à un chiffre significatif (par excès) et la valeur mesurée au même rang :
\[ \boxed{U = (9{,}4 \pm 0{,}3)~\text{V}}. \]
\item Incertitude relative :
\[ \frac{\Delta U}{U} = \frac{0{,}325}{9{,}4} = 0{,}0346 \quad \text{soit environ } \SI{3,5}{\%}. \]
C'est une mesure de qualité moyenne, limitée principalement par la classe de l'appareil.
\end{enumerate}
\end{corrige}

\begin{corrige}[Exercice 6 — Série de huit mesures de masse]
\begin{enumerate}
\item Tableau des données et des écarts à la moyenne :
\begin{center}
\begin{tabular}{|l|*{8}{c|}}
\hline
Pesée $n^{\circ}$ & 1 & 2 & 3 & 4 & 5 & 6 & 7 & 8 \\
\hline
$m$ (g) & 49,8 & 50,2 & 50,0 & 49,9 & 50,1 & 50,3 & 49,7 & 50,0 \\
\hline
$m - \bar{m}$ (g) & $-0{,}2$ & $+0{,}2$ & $0$ & $-0{,}1$ & $+0{,}1$ & $+0{,}3$ & $-0{,}3$ & $0$ \\
\hline
$(m-\bar{m})^2$ (g$^2$) & 0,04 & 0,04 & 0 & 0,01 & 0,01 & 0,09 & 0,09 & 0 \\
\hline
\end{tabular}
\end{center}
\item Valeur moyenne :
\[ \bar{m} = \frac{49{,}8 + 50{,}2 + 50{,}0 + 49{,}9 + 50{,}1 + 50{,}3 + 49{,}7 + 50{,}0}{8} = \frac{400{,}0}{8} = \SI{50,0}{g}. \]
\item Écart-type expérimental (diviseur $n - 1 = 7$) :
\[ s_{n-1} = \sqrt{\frac{\sum (m_i - \bar{m})^2}{n-1}} = \sqrt{\frac{0{,}28}{7}} = \sqrt{0{,}04} = \SI{0,20}{g}. \]
\item Incertitude-type sur la moyenne :
\[ u(\bar{m}) = \frac{s_{n-1}}{\sqrt{n}} = \frac{0{,}20}{\sqrt{8}} = \SI{0,071}{g} \approx \SI{0,07}{g}. \]
\end{enumerate}
\textbf{Point de vigilance :} pour l'écart-type \emph{expérimental} d'un échantillon, on divise par $n-1$ et non par $n$ ; et l'incertitude-type portant sur la \emph{moyenne} fait intervenir $\sqrt{n}$.
\end{corrige}

\begin{corrige}[Exercice 7 — Lecture sur deux instruments]
\begin{enumerate}
\item Règle graduée (instrument analogique) : l'incertitude de lecture est prise égale à une demi-graduation :
\[ \Delta L_{\text{règle}} = \frac{1}{2}~\text{mm} = \SI{0,5}{mm}. \]
Pied à coulisse numérique : l'incertitude est prise égale à $\pm 1$ unité d'affichage :
\[ \Delta L_{\text{pied}} = \SI{0,01}{mm}. \]
\item La résolution de la règle est \SI{1}{mm} ; celle du pied à coulisse est \SI{0,01}{mm}, soit \textbf{100 fois plus fine}.
\item Le \textbf{pied à coulisse} est le plus \emph{sensible} : il réagit à des variations de longueur cent fois plus petites. Il est aussi le plus \emph{précis} au sens courant (mesure la plus fine et la moins incertaine), à condition toutefois d'être correctement étalonné : la précision globale exige à la fois fidélité (faible dispersion) et justesse (absence d'erreur systématique).
\end{enumerate}
\end{corrige}

\begin{corrige}[Exercice 8 — Compatibilité de deux résultats]
\begin{enumerate}
\item Intervalles associés :
\[ L_1 \in [25{,}40 - 0{,}05 \ ;\ 25{,}40 + 0{,}05] = [25{,}35 \ ;\ 25{,}45]~\text{cm}, \]
\[ L_2 \in [25{,}62 - 0{,}05 \ ;\ 25{,}62 + 0{,}05] = [25{,}57 \ ;\ 25{,}67]~\text{cm}. \]
\item Les deux intervalles sont \textbf{disjoints} ($25{,}45 < 25{,}57$) : les résultats ne sont \textbf{pas compatibles}. L'écart entre les valeurs centrales, $25{,}62 - 25{,}40 = \SI{0,22}{cm}$, est très supérieur à la somme des incertitudes ($\SI{0,10}{cm}$).
\item Explication physique probable : l'un des deux groupes est victime d'une \cle{erreur systématique} — par exemple une règle dont le zéro est usé ou décalé, une lecture entachée d'erreur de parallaxe, ou un instrument mal étalonné. \textbf{Vérification :} faire mesurer la \emph{même} longueur par les deux groupes avec le \emph{même} instrument, ou comparer chaque instrument à un étalon (règle de référence vérifiée). Si l'écart persiste avec le même instrument, il faut chercher une cause liée au protocole (positionnement, parallaxe).
\end{enumerate}
\end{corrige}

\begin{corrige}[Exercice 9 — Mesure d'une tension au laboratoire du lycée]
\begin{enumerate}
\item Le calibre \SI{15}{V} correspond à 150 divisions, donc une division vaut $\dfrac{15}{150} = \SI{0,1}{V}$. L'aiguille sur la division 94 donne :
\[ U = 94 \times 0{,}1 = \SI{9,4}{V}. \]
\item Incertitude due à la classe :
\[ \Delta U_{\text{instr}} = \frac{1{,}5 \times 15}{100} = \SI{0,225}{V} \approx \SI{0,23}{V}. \]
\item Incertitude de lecture (demi-division) :
\[ \Delta U_{\text{lect}} = \frac{0{,}1}{2} = \SI{0,05}{V}. \]
\item Incertitude totale :
\[ \Delta U = 0{,}225 + 0{,}05 = 0{,}275 \approx \SI{0,3}{V}. \]
Résultat normalisé (incertitude arrondie à un chiffre significatif, valeur au même rang) :
\[ \boxed{U = (9{,}4 \pm 0{,}3)~\text{V}}. \]
\item Incertitude relative :
\[ \frac{\Delta U}{U} = \frac{0{,}275}{9{,}4} = 0{,}029 \quad \text{soit environ } \SI{2,9}{\%}. \]
C'est une mesure correcte pour un montage de lycée, mais perfectible : un calibre plus petit (\SI{10}{V}) réduirait l'incertitude de classe.
\item L'intervalle de valeurs est $[9{,}1 \ ;\ 9{,}7]$ V. La valeur $\SI{9,00}{V}$ n'appartient \textbf{pas} à cet intervalle : elle n'est \textbf{pas compatible} avec la mesure de l'élève. Il faudrait vérifier le montage, le zéro du voltmètre ou la valeur annoncée par l'enseignant.
\end{enumerate}
\textbf{Barème indicatif (sur 10) :} 1 pt pour $U$ ; 2 pts pour $\Delta U_{\text{instr}}$ (formule + calcul sur le calibre) ; 1 pt pour $\Delta U_{\text{lect}}$ ; 2 pts pour l'écriture normalisée avec chiffres significatifs cohérents ; 2 pts pour l'incertitude relative et le commentaire ; 2 pts pour le test de compatibilité justifié par l'intervalle.\\
\textbf{Points de vigilance :} citer la formule de la classe \emph{avant} le calcul ; préciser que l'incertitude de classe porte sur le calibre ; arrondir l'incertitude par excès et aligner le dernier chiffre de la valeur sur celui de l'incertitude ; conclure la compatibilité par appartenance ou non à l'intervalle.
\end{corrige}

\begin{corrige}[Exercice 10 — Pesées répétées d'un sac de riz à Douala]
\begin{enumerate}
\item Tableau des données (écarts calculés par rapport à $\bar{m} = \SI{4,988}{kg}$, exprimés en grammes pour la lisibilité) :
\begin{center}
\begin{tabular}{|l|*{6}{c|}}
\hline
Pesée $n^{\circ}$ & 1 & 2 & 3 & 4 & 5 & 6 \\
\hline
$m$ (kg) & 4,982 & 4,991 & 4,987 & 4,995 & 4,988 & 4,985 \\
\hline
$m - \bar{m}$ (g) & $-6$ & $+3$ & $-1$ & $+7$ & $0$ & $-3$ \\
\hline
$(m-\bar{m})^2$ (g$^2$) & 36 & 9 & 1 & 49 & 0 & 9 \\
\hline
\end{tabular}
\end{center}
Masse moyenne :
\[ \bar{m} = \frac{4{,}982 + 4{,}991 + 4{,}987 + 4{,}995 + 4{,}988 + 4{,}985}{6} = \frac{29{,}928}{6} = \SI{4,988}{kg}. \]
\item Écart-type expérimental (diviseur $n - 1 = 5$) :
\[ s_{n-1} = \sqrt{\frac{104 \times 10^{-6}}{5}} = \sqrt{20{,}8 \times 10^{-6}} \approx \SI{4,6e-3}{kg} \quad (\text{soit } 4{,}6~\text{g}). \]
\item Incertitude-type sur la moyenne :
\[ u(\bar{m}) = \frac{s_{n-1}}{\sqrt{n}} = \frac{4{,}6 \times 10^{-3}}{\sqrt{6}} \approx \SI{1,9e-3}{kg}. \]
\item Incertitude élargie à \SI{95}{\%} ($k = 2$) :
\[ U = k\,u(\bar{m}) = 2 \times 1{,}9 \times 10^{-3} = 3{,}7 \times 10^{-3} \approx \SI{0,004}{kg}. \]
\item Résultat normalisé :
\[ \boxed{m = (4{,}988 \pm 0{,}004)~\text{kg} \quad (k = 2,\ \text{niveau de confiance } \approx 95\,\%)}. \]
\item L'intervalle de confiance est $[4{,}984 \ ;\ 4{,}992]$ kg. La masse annoncée $\SI{5,000}{kg}$ est \textbf{en dehors} de cet intervalle : elle n'est \textbf{pas compatible} avec les mesures au niveau de confiance \SI{95}{\%}. Le sac est très probablement \textbf{sous-rempli} d'environ \SI{12}{g} ; une autre explication possible serait une erreur systématique de la balance (zéro décalé), qu'il faudrait écarter en pesant une masse étalon. Le commerçant doit soit compléter ses sacs, soit faire vérifier sa balance.
\end{enumerate}
\textbf{Barème indicatif (sur 12) :} 2 pts tableau + moyenne ; 2 pts écart-type (formule $n-1$ exigée) ; 2 pts incertitude-type ; 1 pt incertitude élargie ; 2 pts écriture normalisée ; 3 pts test de compatibilité + conclusion argumentée.\\
\textbf{Points de vigilance :} diviseur $n-1$ pour $s_{n-1}$ ; diviseur $\sqrt{n}$ pour $u(\bar{m})$ ; citer $k = 2$ et le niveau de confiance dans le résultat final ; la conclusion sur la compatibilité doit s'appuyer explicitement sur l'intervalle.
\end{corrige}

\begin{corrige}[Exercice 11 — Contrôle qualité d'un lot de comprimés]
\begin{enumerate}
\item Moyenne :
\[ \bar{m} = \frac{502 + 498 + 501 + 499 + 503 + 500 + 497 + 501 + 499 + 500}{10} = \frac{5000}{10} = \SI{500}{mg}. \]
Écarts (en mg) : $+2\,;-2\,;+1\,;-1\,;+3\,;0\,;-3\,;+1\,;-1\,;0$ ; somme des carrés : $4+4+1+1+9+0+9+1+1+0 = 30$ mg$^2$.
\[ s_{n-1} = \sqrt{\frac{30}{9}} = \sqrt{3{,}33} \approx \SI{1,83}{mg}. \]
\item Incertitude-type sur la moyenne :
\[ u(\bar{m}) = \frac{s_{n-1}}{\sqrt{n}} = \frac{1{,}83}{\sqrt{10}} \approx \SI{0,58}{mg}. \]
Incertitude élargie ($k = 2$) :
\[ U = 2 \times 0{,}58 = 1{,}15 \approx \SI{1,2}{mg}. \]
\item Résultat normalisé :
\[ \boxed{m = (500{,}0 \pm 1{,}2)~\text{mg} \quad (k = 2,\ \text{niveau de confiance } \approx 95\,\%)}. \]
\item L'intervalle de confiance est $[498{,}8 \ ;\ 501{,}2]$ mg. Il est \textbf{entièrement inclus} dans l'intervalle de conformité $[495 \ ;\ 505]$ mg : le comprimé contrôlé \textbf{satisfait le cahier des charges} au niveau de confiance \SI{95}{\%}.
\item L'incertitude-type varie en $\dfrac{1}{\sqrt{n}}$ : pour la diviser par deux, il faut multiplier $n$ par quatre :
\[ \frac{u}{\sqrt{4}} = \frac{u}{2} \quad \Longrightarrow \quad n' = 4 \times 10 = \textbf{40 pesées}. \]
\item Deux sources d'erreur systématique possibles :
\begin{itemize}
\item \textbf{Balance mal tarée} (zéro décalé) : vérifier et régler le zéro avant chaque série de pesées ;
\item \textbf{Balance mal étalonnée} (gain faux) : contrôler périodiquement avec une masse étalon certifiée.
\end{itemize}
(On accepte aussi : courants d'air ou vibrations — placer la balance sur une table stable, à l'abri des courants d'air.)
\end{enumerate}
\textbf{Barème indicatif (sur 14) :} 3 pts moyenne + écart-type ; 2 pts $u(\bar{m})$ ; 1 pt $U$ ; 2 pts écriture normalisée ; 2 pts conformité justifiée par l'inclusion des intervalles ; 2 pts calcul de $n' = 40$ justifié par la loi en $1/\sqrt{n}$ ; 2 pts sources d'erreur + corrections.\\
\textbf{Points de vigilance :} la conformité se justifie par l'\emph{inclusion} de l'intervalle de confiance dans l'intervalle de tolérance, pas par la seule moyenne ; la loi en $1/\sqrt{n}$ doit être citée pour justifier le quadruplement du nombre de pesées ; distinguer clairement erreurs aléatoires (réduites par la répétition) et erreurs systématiques (réduites par l'étalonnage).
\end{corrige}

\newpage

\coursec{Fiche bilan}

\begin{popfigure}
\begin{tikzpicture}[mindmap, grow cyclic, every node/.style={concept, text=white, font=\small\bfseries},
  root concept/.append style={concept color=popink, font=\bfseries},
  level 1/.append style={level distance=3.6cm, sibling angle=60, font=\small},
  level 1 concept/.append style={font=\small\bfseries},
  level 2/.append style={level distance=2.6cm, sibling angle=45, font=\scriptsize},
  level 2 concept/.append style={font=\scriptsize}]
\node[root concept]{Mesures et incertitudes}
  child[concept color=popblue]{ node {Mesure, erreur, incertitude}
    child[concept color=popblue!70]{ node {erreur $=$ mesur\'e $-$ vrai} }
    child[concept color=popblue!70]{ node {incertitude $=$ doute} }
  }
  child[concept color=poporange]{ node {Qualit\'es d'un instrument}
    child[concept color=poporange!70]{ node {sensibilit\'e} }
    child[concept color=poporange!70]{ node {fid\'elit\'e, justesse} }
    child[concept color=poporange!70]{ node {pr\'ecision} }
  }
  child[concept color=popgreen]{ node {Mesure unique}
    child[concept color=popgreen!70]{ node {$U = k\,u$} }
    child[concept color=popgreen!70]{ node {lecture : $\pm$ graduation/2} }
  }
  child[concept color=popgold]{ node {S\'erie de mesures}
    child[concept color=popgold!70]{ node {moyenne $\bar{x}$} }
    child[concept color=popgold!70]{ node {\'ecart-type $s$} }
  }
  child[concept color=poppurple]{ node {Incertitude-type}
    child[concept color=poppurple!70]{ node {$u = s/\sqrt{n}$} }
    child[concept color=poppurple!70]{ node {type A / type B} }
  }
  child[concept color=poppink]{ node {R\'esultat final}
    child[concept color=poppink!70]{ node {$x = \bar{x} \pm U$} }
    child[concept color=poppink!70]{ node {chiffres significatifs} }
  };
\end{tikzpicture}
\legende{Carte mentale de la le\c{c}on : les six id\'ees forces autour de la mesure.}
\end{popfigure}

\begin{bilan}
\textbf{D\'efinitions cl\'es}
\begin{itemize}
  \item \cle{Mesurage} : ensemble des op\'erations permettant de d\'eterminer exp\'erimentalement la valeur d'une grandeur physique (le \cle{mesurande}).
  \item \cle{Erreur de mesure} : $e = x_{\text{mesur\'e}} - x_{\text{vrai}}$. Elle est inconnaissable exactement ; on distingue erreur \cle{syst\'ematique} (reproductible, due \`a un d\'efaut d'\'etalonnage ou de m\'ethode) et erreur \cle{al\'eatoire} (fluctuations impr\'evisibles).
  \item \cle{Incertitude} : param\`etre qui caract\'erise la dispersion des valeurs attribuables au mesurande ; elle quantifie le \emph{doute} sur le r\'esultat.
  \item \cle{Sensibilit\'e} : aptitude \`a d\'etecter de petites variations de la grandeur (quotient de la variation de l'indication par celle de la grandeur).
  \item \cle{Fid\'elit\'e} : aptitude \`a donner des indications proches les unes des autres lors de mesures r\'ep\'et\'ees (faible dispersion).
  \item \cle{Justesse} : aptitude \`a donner des indications proches de la valeur vraie (faible erreur syst\'ematique).
  \item \cle{Pr\'ecision} : qualit\'e globale r\'eunissant fid\'elit\'e \emph{et} justesse.
\end{itemize}

\textbf{Formules \`a conna\^{\i}tre par c\oe ur}
\begin{center}
\begin{tabularx}{\linewidth}{|l|X|l|}
\hline
\rowcolor{popblueL} \textbf{Grandeur} & \textbf{Formule} & \textbf{Contexte} \\
\hline
Moyenne & $\bar{x} = \dfrac{1}{n}\displaystyle\sum_{i=1}^{n} x_i$ & s\'erie de $n$ mesures \\
\hline
\'Ecart-type exp\'erimental & $s = \sqrt{\dfrac{1}{n-1}\displaystyle\sum_{i=1}^{n}(x_i - \bar{x})^2}$ & dispersion de la s\'erie \\
\hline
Incertitude-type (type A) & $u(\bar{x}) = \dfrac{s}{\sqrt{n}}$ & sur la moyenne \\
\hline
Incertitude-type (type B) & $u = \dfrac{a}{\sqrt{3}}$ & lecture : demi-graduation $a$ (loi uniforme) \\
\hline
Incertitude \'elargie & $U = k\,u$ & $k = 2$ pour $\approx 95\,\%$ ; $k = 3$ pour $\approx 99\,\%$ \\
\hline
Intervalle de confiance & $[\,\bar{x} - U \;;\; \bar{x} + U\,]$ & probabilit\'e fix\'ee par $k$ \\
\hline
R\'esultat final & $x = \bar{x} \pm U$ (unit\'e) & $U$ arrondie \`a 1 ou 2 chiffres significatifs \\
\hline
\end{tabularx}
\end{center}

\textbf{M\'ethodes express}
\begin{enumerate}
  \item \textbf{Mesure unique} : estimer l'incertitude de lecture ($\pm$ une demi-graduation pour un appareil analogique, $\pm$ un digit pour un appareil num\'erique), composer avec la notice constructeur, \'ecrire $x \pm U$.
  \item \textbf{S\'erie de mesures} : organiser les donn\'ees en tableau $\to$ calculer $\bar{x}$ et $s$ $\to$ $u = s/\sqrt{n}$ $\to$ $U = k\,u$ $\to$ \'ecrire $\bar{x} \pm U$ avec l'unit\'e SI.
  \item \textbf{R\'edaction du r\'esultat} : arrondir $U$ par exc\`es (1 ou 2 chiffres significatifs), puis arrondir $\bar{x}$ au m\^eme rang d\'ecimal que $U$.
\end{enumerate}
\end{bilan}

\begin{aretenir}[Checklist avant le Probatoire]
\begin{itemize}
  \item[$\square$] Je sais distinguer \cle{erreur} (inconnue) et \cle{incertitude} (estim\'ee).
  \item[$\square$] Je sais distinguer erreur syst\'ematique et erreur al\'eatoire, et donner un exemple de chacune.
  \item[$\square$] Je connais les quatre qualit\'es d'un instrument : sensibilit\'e, fid\'elit\'e, justesse, pr\'ecision.
  \item[$\square$] Je sais qu'une mesure unique analogique s'accompagne d'une incertitude de lecture d'une demi-graduation.
  \item[$\square$] Je sais calculer une moyenne $\bar{x}$ et un \'ecart-type $s$ (diviseur $n-1$ !).
  \item[$\square$] Je sais que l'incertitude-type sur la moyenne est $u = s/\sqrt{n}$.
  \item[$\square$] Je sais passer de l'incertitude-type \`a l'incertitude \'elargie : $U = k\,u$, avec $k = 2$ pour un niveau de confiance d'environ $95\,\%$.
  \item[$\square$] Je sais \'ecrire l'intervalle de confiance $[\bar{x}-U \,;\, \bar{x}+U]$ et l'interpr\'eter.
  \item[$\square$] J'arrondis l'incertitude par exc\`es et j'aligne le dernier chiffre de la valeur sur celui de l'incertitude.
  \item[$\square$] Je n'oublie jamais l'unit\'e SI dans le r\'esultat final : $x = \bar{x} \pm U$ (unit\'e).
  \item[$\square$] Je sais pr\'esenter les donn\'ees exp\'erimentales en tableau et en graphe avec axes gradu\'es et l\'egend\'es.
  \item[$\square$] Je v\'erifie la coh\'erence de mes r\'esultats (ordre de grandeur, analyse dimensionnelle).
\end{itemize}
\end{aretenir}

\findecours

\end{document}
